%! TeX root = ./main.tex
% vim: set spelllang=en,fr:
\selectlanguage{french}
\chapter{Synesth: réconciliation segmentale avec hôtes non-échantillonnés}\label{chap:synesth}

L’une des limites de la méthode proposée au chapitre précédent est qu’elle ignore les interactions possibles entre les unités étudiées et des hôtes qui ne feraient pas partie de l’échantillon choisi, par exemple s’ils appartiennent à une lignée qui s’est éteinte depuis lors.
Dans le contexte de la réconciliation avec transferts horizontaux, de telles interactions sont possibles: une unité pourrait être transférée d’un hôte \(A\), ancêtre d’un hôte échantillonné, vers un hôte \(B\) inconnu, puis à nouveau transféré vers un descendant \(C \tdesc A\).
Sous le modèle précédent, ce scénario serait interprété comme une série d’autant de spéciations et de pertes que d’hôtes ancestraux qui séparent \(A\) de \(C\), ce qui peut avoir tendance à écarter artificiellement ce type de possibilité.
Cette problématique a été étudiée par plusieurs auteurs dans le contexte de la réconciliation bilatérale, par exemple récemment par \textcite{Weiner2021ImprovedDuplicationTransfer}.
Nous verrons dans ce chapitre comment étendre le modèle pour prendre en compte cette possibilité.

Nous montrons également comment intégrer un nouveau type d’événement de \emph{scission} dans le modèle, suivant lequel une synténie se partitionne en deux synténies disjointes.
Cet événement additionnel, suggéré par Russell~\textsc{Schwartz} lors de discussions à la conférence RECOMB~2022, permet de mieux prendre en compte la réalité évolutive des synténies.

Finalement, un autre enjeu important est le besoin d’estimer des coûts pour les différents types d’événements du modèle.
En pratique, la valeur exacte des coûts a une influence importante sur la teneur des scénarios inférés, bien que les procédures d’estimation de ces coûts pourraient avoir une certaine incertitude.
Il devient alors pertinent de qualifier précisément l’impact des variations des coûts sur les résultats.
Une possibilité coûteuse et imprécise consiste à échantillonner l’espace des coûts et à exécuter l’algorithme pour chacun des choix, une méthode qui s’apparente au \emph{grid search} en apprentissage automatique.
Une approche préférable consiste à caractériser l’espace des solutions Pareto-optimales, c’est-à-dire des solutions dont le nombre d’événements de chaque type fait qu’il existe un certain ensemble de coûts qui les rendent meilleures que les autres.
En nous appuyant sur les travaux de \textcite{LibeskindHadas2014ParetoOptimalPhylogenetic}, nous montrerons qu’il est possible de brosser efficacement le \emph{paysage} des scénarios de réconciliation segmentale pour une certaine instance, indépendamment de tout choix spécifique de coûts.
Pour ce faire, nous procéderons à une étude combinatoire de l’espace des réconciliations segmentales sous ce modèle étendu.

% TODO: Migrate notations to notations.tex

\newcommand{\ancest}{\leq}
\newcommand{\separ}{\mathrel{\|}}
% \newcommand{\rootOf}{\mathrm{r}}
% \newcommand{\parentOf}{\mathrm{p}}
% \newcommand{\treeNodes}{\mathrm{V}}
% \newcommand{\treeLeaves}{\mathrm{L}}
\newcommand{\treeInner}{\mathrm{I}}
% \newcommand{\treeEdges}{\mathrm{E}}
% \newcommand{\treeEdge}[2]{(#1,#2)}
\newcommand{\treePath}[3][]{\mathrm{P}_{#1}(#2, #3)}
% \newcommand{\treeDist}[3][]{\mathrm{D}_{#1}(#2, #3)}
% \newcommand{\LCA}{\mathrm{lca}}
\newcommand{\children}{\mathrm{ch}}
% \newcommand{\leftChild}{_\mathrm{\ell}}
% \newcommand{\rightChild}{_\mathrm{r}}
\newcommand{\Child}{_\mathrm{c}}
% \newcommand{\geneFamily}{\Gamma}
% \newcommand{\geneFamilySet}{\mathcal{F}}
\newcommand{\augS}{\speciesTree^*}
% \newcommand{\F}{{\mathcal F}}
% \newcommand{\taxa}{{\Sigma}}
% \newcommand{\speciesOf}{s}
% \newcommand{\syntenyFamily}{\mathcal{X}}
% \newcommand{\syntenyTree}{\mathcal{T}}
% \newcommand{\syntenyOf}{x}
\newcommand{\restr}[2]{#1|_{#2}}
\newcommand{\allEvents}{\mathcal{E}}
% \newcommand{\eventSpe}{\mathrm{Spe}}
% \newcommand{\eventDup}{\mathrm{Dup}}
\newcommand{\eventCut}{\mathrm{Cut}}
\newcommand{\eventTrDup}{\mathrm{TrDup}}
\newcommand{\eventTrCut}{\mathrm{TrCut}}
\newcommand{\eventTrAll}{\mathrm{Tr}}
% \newcommand{\eventLoss}{\mathrm{Loss}}
% \newcommand{\eventGain}{\mathrm{Gain}}
\newcommand{\Hist}{\mathcal{H}}
% \newcommand{\speciesTree}{S}
% \newcommand{\eventOf}{e}
\newcommand{\eventCost}{\delta}
\newcommand{\histCost}{\mathrm{c}}
\newcommand{\targetChild}{_\mathrm{t}}
\newcommand{\conservChild}{_\mathrm{k}}
\newcommand{\solSpace}{\mathbb{H}}
\newcommand{\minSolSpace}{\solSpace^{\mathrm{min}}}
\newcommand{\setNat}{\mathbb{N}}
% \newcommand{\setReal}{\mathbb{R}}
\newcommand{\eventVec}{\mathrm{ev}}
\newcommand{\paretoMin}[1]{\min{}_{#1}\,}
\newcommand{\minVecSet}{\eventVec^{\mathrm{min}}}
\newcommand{\makeNode}[1]{[#1]}
\newcommand{\optPath}{\mathrm{path}}
\newcommand{\histExt}{\preceq}
\newcommand{\gainAt}{\mathrm{gain}}
\newcommand{\syntenyMinAt}{\syntenyOf^{\mathrm{min}}}
\newcommand{\minCost}{c^{\mathrm{min}}}
\newcommand{\visLeaves}{\treeLeaves_V}

%%%%%%%%%%%%%%%%%%%%%%%

\selectlanguage{english}
\articlehat{
    \textfrench{La suite de ce chapitre est une reproduction partielle de l’article intitulé \enquote{Synesth: Comprehensive syntenic reconciliation with unsampled lineages}.
    Cet article a été écrit entre octobre~2022 et décembre~2023 et a été publié en avril~2024 dans le journal Algorithms chez MDPI.}
}
\title{Synesth: Comprehensive syntenic\\reconciliation with unsampled lineages}
\author{Mattéo Delabre \and Nadia El-Mabrouk}
\institute{Département d'informatique et de recherche opérationnelle, Université de Montréal}
\makearticlehead

\begin{abstract}
    We present Synesth, the most comprehensive and flexible tool for tree reconciliation allowing for segmental events --- involving multiple genes --- including duplications, transfers, fissions, and transient events going through unsampled species.
    We examine the combinatorial properties of this extended tree reconciliation model and study various associated optimization problems related to various granularities of the solution space.
    First, the infinite set of histories is reduced to a finite but exponential set of Pareto-optimal event histories (in terms of counts of each event type), then to a polynomial set of Pareto-optimal event vectors, eventually ending with minimum event cost histories given an event cost function. An inductive algebraic characterization of the output using a different algebra for each space granularity leads to efficient dynamic programming algorithms, ultimately ending with an \(O(mn)\) time-complexity algorithm for computing a minimum event-cost history (\(m\) and \(n\): number of nodes of the two trees).
    This time complexity matches that of the fastest algorithms for classical (i.e. without syntenies and associated events) reconciliation with transfers.
    We show how Synesth can be applied for inferring Pareto-optimal evolutionary scenarios of CRISPR-Cas systems in a set of bacterial genomes.
    An implementation of the algorithm is available at: \url{https://github.com/UdeM-LBIT/superrec2/tree/algo2024}.

    \keywords{Tree reconciliation \and Algebraic dynamic programming \and Multi-objective optimization \and Horizontal gene transfer \and Synteny \and Fission \and Unsampled lineages.}
\end{abstract}

\section{Introduction}

A \emph{gene/species trees reconciliation} is an embedding of a gene tree into a species tree explicating the topological difference between the two trees through a sequence of events shaping the gene family inside the species tree.
Gene/species tree reconciliation has been widely studied since the 80s~\parencite{Goodman1979FittingGeneLineage}, first focusing on duplication~(D) and loss~(L) events and then extending to horizontal gene transfers~(T)~\parencite{Bansal2012EfficientAlgorithmsReconciliation,Donati2015EfficientTreeReconciliation,Tofigh2011SimultaneousIdentificationDuplications} and other events such as hybridization or incomplete lineage sorting (review in~\textcite{ElMabrouk2019GeneFamilyEvolution}). One of the major drawbacks of classical reconciliation is that gene families are considered separately from one another, which is not appropriate for genes organized in \emph{syntenies}, i.e. colocalized genes likely evolving together through segmental events.

Although some work has been done to infer the evolution of adjacencies~\parencite{Duchemin2017DecostarReconstructingAncestral}, group individual events into segmental ones~\parencite{Duchemin2017PhylogenyDependenciesDependencies}, or minimize \enquote{duplication episodes}~\parencite{Dondi2019ReconcilingMultipleGenes,Paszek2018EfficientAlgorithmsGenomic}, none of these methods are intended to explicitly look for evolutionary scenarios minimizing segmental events. We presented the first algorithm generalizing reconciliation to synteny trees (i.e. with leaves representing syntenies rather than single genes) and segmental events for the DL distance in~\textcite{Delabre2020EvolutionSegmentalDuplications} and extended it to horizontal transfers with \emph{SuperDTL} in~\textcite{Anselmetti2022ReconciliationSegmentalDuplication}.

Here, we present Synesth (for \emph{SYNteny Evolution in SegmenTal Histories}), an extended reconciliation model accounting for fissions whereby part of a synteny is detached to another locus or species, in addition to losses, gains, duplications and transfers.
Moreover, as the choice of the species included in a study, or equivalently the absence of species which are not chosen, unsampled or extinct, has a serious impact on the output of a reconciliation algorithm~\parencite{Jacox2016EcceteraComprehensiveGene,Szoelloesi2013LateralGeneTransfer,Weiner2021ImprovedDuplicationTransfer}, we account for transfers going through unsampled species. 

Among the infinite space of possible histories explicating the evolution of a synteny tree inside a species tree both given as input, we are interested in selecting the most likely ones.
A possible way is by assigning costs (or probabilities) to event types and selecting histories minimizing (resp. maximizing) the overall cost (resp. probability).
The resulting histories can vary significantly depending on the costs, which are usually challenging to determine and strongly depend on the taxa under study~\parencite{LibeskindHadas2014ParetoOptimalPhylogenetic,David2011RapidEvolutionaryInnovation}.
Instead, our goal here is to provide a landscape of possible histories independently of costs, before eventually choosing a solution for a particular choice of costs.

We subdivide the space of histories by formally defining nested classes.
First, the space is reduced to a finite but exponential set of Pareto-optimal histories (in terms of counts of each event type), then to a polynomial set of Pareto-optimal vectors, eventually ending with minimum-cost solutions.
We develop efficient and exact algorithms using an algebraic dynamic programming approach with a different algebra for each class.
The inductive characterization of Pareto-optimal histories given in \cref{sec:synesth:pareto-histories} leads to a polynomial-time dynamic programming algorithm to output all Pareto-optimal vectors (\cref{sec:synesth:pareto-vectors}) and ultimately to an \(\bigo(mn)\) time-complexity algorithm (\(m\) and \(n\): number of nodes of the two trees) to output a minimum event cost history (\Cref{sec:synesth:min-cost-histories}), matching time complexity of the fastest algorithms for classical DTL reconciliation~\parencite{Weiner2021ImprovedDuplicationTransfer}.

In \cref{sec:synesth:results}, we apply Synesth to CRISPR-associated (Cas) gene syntenies.
Taking advantage of previous work~\parencite{LibeskindHadas2022TreeReconciliationMethods,LibeskindHadas2014ParetoOptimalPhylogenetic}, we output a comprehensive representation of the solution landscape comprised of \enquote{cost regions}, such that any two event cost assignments within the same region lead to the same histories.

We first introduce the required notations in \cref{sec:synesth:defs}, the evolutionary model in \cref{sec:synesth:def-model} and the solution space subdivision and considered optimization problems in \cref{sec:synesth:opt-probs}.

\section{Preliminary notation}\label{sec:synesth:defs}

% TODO: Merge with introductory notation chapter

All trees are considered rooted.
Given a tree \(T\), we denote by \(\rootOf(T)\) its root, by \(\treeNodes(T)\) its node set, by \(\treeLeaves(T) \subseteq \treeNodes(T)\) its leaf set, and we let \(\treeInner(T) = \treeNodes(T) - \treeLeaves(T)\) be the set of its internal nodes.
A node \(v'\) is an \emph{ancestor} of \(v\)  if \(v'\) is on the path from \(\rootOf(T)\) to \(v\); the \emph{parent} \(\parentOf(v)\) of \(v\), of which \(v\) is a \emph{child}, directly precedes \(v\) on this path.
Conversely, \(v\) is a \emph{descendant} of \(v'\).
This ancestor--descendant relation is denoted \(\ancest\) and forms a partial order on nodes, in which the root is minimal and the leaves are maximal.
Any pair of nodes \(v\) and \(v'\) not ordered by this relation are said to be \emph{separated}, which we denote \(v \separ v'\).
Notice that a node \(v\) is both an ancestor and a descendant of itself; whenever this case needs to be excluded, we will talk about strict ancestors and strict descendants.

We denote by \(\treeEdges(T)\) the set of edges of \(T\), each represented by a pair of nodes \(\treeEdge{\parentOf(v)}{v}\).
For any two nodes \(v\) and \(v'\) of \(T\), there exists a unique path from \(v\) to \(v'\) that we denote \(\treePath[T]{v}{v'} \subseteq \treeEdges(T)\).
The distance between \(v\) and \(v'\) is defined as \(\treeDist[T]{v}{v'} = |\treePath[T]{v}{v'}|\).
Given a node \(v\) of \(T\), \(\subtree{T}{v}\) is the \emph{subtree} of \(T\) rooted at \(v\) (i.e. containing only the descendants of \(v\)).
The \emph{lowest common ancestor}~(LCA) of a subset \(V\) of nodes, denoted \(\LCA_T(V)\), is the ancestor of all nodes in \(V\) which is the most distant from the root.

A tree~\(T'\) is said to be an \emph{extension} of a tree~\(T\) if~\(T'\) can be obtained from~\(T\) by a sequence of operations among:
(1) \emph{subdividing} an edge \((u, w)\) by adding a new node \(v\) and replacing \((u, w)\) by two edges \((u, v)\) and \((v, w)\);
(2) \emph{grafting} a new node \(v\) below an existing node \(u\) by adding the edge \((u, v)\);
(3) \emph{rerooting} the tree to a new node \(u\) by adding the edge \((u, \rootOf(T))\).

The set of children of any node \(v\) is denoted by \(\children(v)\).
If \(|\children(v)| = 1\), then \(v\) is said to be \emph{unary} and we denote its only child by \(v\Child\).
If \(|\children(v)| = 2\), then it is said to be \emph{binary} and, unless specified otherwise, we denote its children by \(v\leftChild\) and \(v\rightChild\) in no particular order.
A \emph{binary tree} is a tree where all internal (non-leaf) nodes are binary.
If all internal nodes are unary or binary, then the tree is \emph{partially binary}.

If \(\geneFamilySet\) is a set of gene families, then a \emph{synteny} on \(\geneFamilySet\) is a subset~\(X \subseteq \geneFamilySet\) representing a group of genes assumed to have jointly evolved.
Notice that we ignore the relative order of genes in the genomic region, as well as the physical distance between genes and regions.
The genes of a synteny are considered to all belong to different gene families (i.e. repeated gene families inside a synteny are forbidden),
therefore a gene is simply identified by the family \(\Gamma \in \F\) it belongs to.

A \emph{species tree} on a set \(\taxa\) of species is a tree \(S\) with a bijection \(\speciesOf : \treeLeaves(S) \rightarrow \Sigma\).
For a set of syntenies \(\syntenyFamily\), a \emph{synteny tree} \(\tuple{T, \syntenyOf, \speciesOf}\) is a tuple where \(T\) is a tree, \(\syntenyOf : \treeLeaves(T) \rightarrow \syntenyFamily\) and \(\speciesOf : \treeLeaves(T) \rightarrow \Sigma\), are two functions, the second indicating the species to which each synteny belongs.

Finally, the restriction of a function \(f\) to a subset \(A\) of its domain is denoted \(\restr{f}{A}\).

\section{Evolutionary histories for syntenies}\label{sec:synesth:def-model}

We model the evolution of syntenies through the following segmental events:
codivergence with the host species~(\enquote{\(\eventSpe\)}),
duplication of a synteny segment~(\enquote{\(\eventDup\)}),
fission of a synteny~(\enquote{\(\eventCut\)}),
transfer of a duplicated or cut segment (resp. \enquote{\(\eventTrDup\)} or \enquote{\(\eventTrCut\)}),
and gain or loss of a segment~(\enquote{\(\eventGain\)}, \enquote{\(\eventLoss\)}).
Losses can be partial in the sense that only a subset of genes in a segment are lost.
Evolutionary histories are sequences of such events, as formally defined below (see example in \cref{fig:synesth:example}).

\begin{definition}[Evolutionary history for syntenies]\label{def:synesth:hist}
    A history \(\Hist\) on a binary species tree \(\speciesTree\) is a tuple \(\tuple{H, \eventOf, \syntenyOf, \speciesOf}\), where \(H\) is a partially binary tree.
    Each node \(v \in \treeNodes(H)\) is labeled with a species \(\speciesOf(v) \in \treeNodes(\speciesTree)\) and a synteny (i.e. a set of gene families) \(\syntenyOf(v)\).
    Each internal node is additionally labeled with an event \(\eventOf(v) \in \allEvents = \{\eventSpe,\allowbreak \eventDup,\allowbreak \eventCut,\allowbreak \eventTrDup,\allowbreak \eventTrCut,\allowbreak \eventGain,\allowbreak \eventLoss\}\) acting on \(\syntenyOf(v)\) and \(\speciesOf(v)\).
    These labelings satisfy the following conditions:
    \begin{enumerate}
        \item\label{def:synesth:hist-spe}
            If \(\eventOf(v) = \eventSpe\),
            with \(\children(v) = \{v', v''\}\) and \(\sigma = \speciesOf(v)\),
            then \(\syntenyOf(v) = \syntenyOf(v') = \syntenyOf(v'')\),
            \(\speciesOf(v') = \sigma\leftChild\),
            and \(\speciesOf(v'') = \sigma\rightChild\).

    	\item If \(\eventOf(v) \in \{\eventDup, \eventCut, \eventTrDup, \eventTrCut\}\),
            with \(\children(v) = \{v\targetChild, v\conservChild\}\):

    	\begin{enumerate}
            \item if \(\eventOf(v) \in \{\eventDup, \eventTrDup\}\),
                then \(\syntenyOf(v\targetChild) \subseteq \syntenyOf(v) = \syntenyOf(v\conservChild)\);

            \item if \(\eventOf(v) \in \{\eventCut, \eventTrCut\}\),
                then \(\syntenyOf(v\targetChild) \cup \syntenyOf(v\conservChild) = \syntenyOf(v)\)
                and \(\syntenyOf(v\targetChild) \cap \syntenyOf(v\conservChild) = \emptyset\)
                (but \(\syntenyOf(v\targetChild) = \emptyset\) or \(\syntenyOf(v\conservChild) = \emptyset\) is allowed);

            \item if \(\eventOf(v) \in \{\eventDup, \eventCut\}\),
                then \(\speciesOf(v) = \speciesOf(v\targetChild) = \speciesOf(v\conservChild)\);

            \item if \(\eventOf(v) \in \{\eventTrDup, \eventTrCut\}\),
                then \(\speciesOf(v) \separ \speciesOf(v\targetChild)\),
                and \(\speciesOf(v) = \speciesOf(v\conservChild)\).
        \end{enumerate}

        \item If \(\eventOf(v) \in \{\eventGain, \eventLoss\}\)
            with \(\children(v) = \{v\Child\}\),
            then \(\speciesOf(v\Child) = \speciesOf(v)\) and:

        \begin{enumerate}
            \item if \(\eventOf(v) = \eventGain\),
                then \(\syntenyOf(v\Child) \supsetneq \syntenyOf(v)\).
                
            \item if \(\eventOf(v) = \eventLoss\),
                then \(\syntenyOf(v\Child) \subsetneq \syntenyOf(v)\)
                (a loss is \emph{full} if \(\syntenyOf(v\Child) = \emptyset\),
                and \emph{partial} otherwise);
        \end{enumerate}

        \item\label{def:synesth:hist-gain-once}
            For each gene family \(\geneFamily\), exactly one \(\eventGain\) event in \(H\) involves \(\geneFamily\).

        \item The only nodes \(v\) of the history such that \(\syntenyOf(v) = \emptyset\) are its root and its leaves.
    \end{enumerate}
    
    Finally, we denote by \(\eventTrAll\) any event in \(\{\eventTrDup, \eventTrCut\}\).
\end{definition}

Condition (4) excludes \emph{convergent gains}, i.e., histories where the same gene family is gained separately in different parts of the history.
Moreover, as in~\textcite{Weiner2021ImprovedDuplicationTransfer}, we will allow for transfers from unsampled or extinct species by augmenting the species tree. For example, in \cref{fig:synesth:example}, the synteny \(\{1,2\}\) in species \(A\) is the result of a transfer from an unsampled species.
In the following, unless specified otherwise, all histories are on \(\augS\), as formally defined below.

\begin{definition}[Augmented species tree]\label{def:augmented-species}
    A tree~\(\speciesTree\) can be augmented into~\(\augS\) by adding \emph{unsampled leaves} as follows:
    (1) subdivide each edge \((v, v')\) of \(\treeEdges(\speciesTree)\) into two edges \((v, z)\) and \((z, v')\) through a new node \(z\),
    (2) connect each \(z\) to a new unsampled leaf, and
    (3) create a new root \(\rootOf(\augS)\) whose two children are \(\rootOf(\speciesTree)\) and a new unsampled leaf.
    Edges leading to unsampled leaves are called \emph{unsampled edges}.
\end{definition}

Finally, notice that \cref{def:synesth:hist} does not require the species involved in a transfer to be contemporary, nor does it forbid biologically infeasible cyclic histories, such as one resulting from a transfer from a species \(A\) to a species \(B\), and then back from species \(B\) to an ancestor of species \(A\)
(for a more precise definition of acyclicity, see~\textcite{Tofigh2011SimultaneousIdentificationDuplications}).

\begin{figure}[htb]
    \centering
    \sidecols{.4\textwidth}{.5\textwidth}{
    \centering
    \begin{forest}
        forked edges,
        for tree={
            grow'=0,
            s sep'+=.625cm,
            l=.8em,
            l sep=0pt,
        },
        before typesetting nodes={
            delay={
                where content={}{coordinate}{},
            },
            where n children=0{
                tier=extant,
                label/.process={Ow{content}{right:#1}},
                content=,
            }{},
        },
        [
            [\subnode{synteny-1}{1, 2}]
            [
                [
                    [
                        [\subnode{synteny-2}{2, 3}]
                        [\subnode{synteny-3}{2, 3}]
                    ]
                    [
                        [\subnode{synteny-4}{1, 3}]
                        [
                            [\subnode{synteny-5}{1, 2}]
                            [\subnode{synteny-6}{3}]
                        ]
                    ]
                ]
                [\subnode{synteny-7}{1, 2}]
            ]
        ]
    \end{forest}\hspace{2em}%
    \begin{forest}
        forked edges,
        for tree={
            grow=west,
            s sep'+=.525cm,
            l=.8em,
            l sep=0pt,
        },
        extant leaf/.style={
            tier=extant,
        },
        unsampled leaf/.style={
            tier=unsampled,
            edge=densely dotted,
        },
        before typesetting nodes={
            delay={
                where content={}{coordinate}{},
            },
            where n children=0{
                label/.process={Ow{content}{left:\textbf{#1}}},
                content=,
            }{},
        },
        [
            [, unsampled leaf]
            [
                [
                    [
                        [
                            [\subnode{species-A}{A}, extant leaf]
                            [, unsampled leaf]
                        ]
                        [
                            [\subnode{species-B}{B}, extant leaf]
                            [, unsampled leaf]
                        ]
                    ]
                    [, unsampled leaf]
                ]
                [
                    [\subnode{species-C}{C}, tier=extant]
                    [, unsampled leaf]
                ]
            ]
        ]
    \end{forest}
    \begin{tikzpicture}[
        overlay,
        remember picture,
        link/.style={
            black!30, thick, dashed,
            out=180,
            in=0,
            out distance=.35cm,
            in distance=.35cm,
        }
    ]
        \path[link] (species-A) edge (synteny-1.east);
        \path[link] (species-A) edge (synteny-2.east);
        \path[link] (species-B) edge (synteny-3.east);
        \path[link] (species-A) edge (synteny-4.east);
        \path[link] (species-B) edge (synteny-5.east);
        \path[link] (species-C) edge (synteny-6.east);
        \path[link] (species-C) edge (synteny-7.east);
    \end{tikzpicture}

    \(\syntenyTree\)
    \hspace{6em}
    \(\augS\)

    (1)
}{
    \centering
    {\small\begin{tikzpicture}[y={(.65pt,0)},x={(0,-.85pt)}]
% hosts
\path[host background] (0.0,75.175) [rounded corners={3pt}] -- (0.0,16.922) -- (19.8265,16.922) [sharp corners] -- (19.8265,-5.0) -- (41.3457,-5.0) [rounded corners={3pt}] -- (41.3457,16.922) -- (159.7763,16.922) [sharp corners] -- (159.7763,53.532) -- (140.6863,53.532) [rounded corners={3pt}] -- (140.6863,47.532) -- (19.8265,47.532) [sharp corners] -- (19.8265,75.175) -- cycle;
\path[unsampled host background] (0.0,75.175) [rounded corners={3pt}] -- (0.0,303.2216) --  (19.8265,303.2216) [sharp corners] -- (19.8265,75.175) -- cycle;
\path[host background] (75.7306,90.142) [rounded corners={3pt}] -- (75.7306,53.532) -- (140.6863,53.532) [sharp corners] -- (140.6863,53.532) -- (159.7763,53.532) [rounded corners={3pt}] -- (159.7763,53.532) -- (182.5113,53.532) [sharp corners] -- (182.5113,215.175) -- (159.7763,215.175) [rounded corners={3pt}] -- (159.7763,84.142) -- (112.9105,84.142) [sharp corners] -- (112.9105,90.142) -- cycle;
\path[host background] (41.3457,215.175) [rounded corners={3pt}] -- (41.3457,137.845) -- (75.7306,137.845) [sharp corners] -- (75.7306,90.142) -- (112.9105,90.142) [rounded corners={3pt}] -- (112.9105,137.845) -- (140.6863,137.845) [sharp corners] -- (140.6863,215.175) -- (112.9106,215.175) [rounded corners={3pt}] -- (112.9106,209.175) -- (75.7306,209.175) [sharp corners] -- (75.7306,215.175) -- cycle;
\path[host background] (41.3457,215.175) [rounded corners={3pt}] -- (41.3457,303.2216) -- node[host label] {A} (75.7306,303.2216) [sharp corners] -- (75.7306,215.175) -- cycle;
\path[host background] (112.9106,215.175) [rounded corners={3pt}] -- (112.9106,303.2216) -- node[host label] {B} (140.6863,303.2216) [sharp corners] -- (140.6863,215.175) -- cycle;
\path[host background] (159.7763,215.175) [rounded corners={3pt}] -- (159.7763,303.2216) -- node[host label] {C} (182.5113,303.2216) [sharp corners] -- (182.5113,215.175) -- cycle;
% branches
\path[branch] (31.8007,32.227) [rounded corners={1.5pt}] -- (150.2313,32.227) [sharp corners] -- (150.2313,68.837);
\path[branch] (31.8007,32.227) [rounded corners={1.5pt}] -- (9.9132,32.227) [sharp corners] -- (9.9132,248.4765);
\path[branch] (31.8007,4.4609) -- (31.8007,32.227) -- cycle;
\path[branch] (9.9132,248.4765) -- (9.9132,282.6116) -- cycle;
\path[transfer branch] (9.9132,248.4765) -- (47.7432,248.4765);\path[branch] (47.7432,248.4765) -- (47.7432,267.6116);
\path[branch] (150.2313,68.837) [rounded corners={1.5pt}] -- (94.3206,68.837) [sharp corners] -- (94.3206,95.1334);
\path[branch] (150.2313,68.837) [rounded corners={1.5pt}] -- (176.1738,68.837) [sharp corners] -- (176.1738,267.6116);
\path[branch] (103.3656,158.705) [rounded corners={1.5pt}] -- (69.3331,158.705) [sharp corners] -- (69.3331,245.1202);
\path[branch] (103.3656,158.705) [rounded corners={1.5pt}] -- (130.103,158.705) [sharp corners] -- (130.103,240.1433);
\path[branch] (85.2756,193.87) [rounded corners={1.5pt}] -- (58.5382,193.87) [sharp corners] -- (58.5382,267.6116);
\path[branch] (85.2756,193.87) [rounded corners={1.5pt}] -- (119.308,193.87) [sharp corners] -- (119.308,267.6116);
\path[branch] (94.3206,118.985) [rounded corners={1.5pt}] -- (103.3656,118.985) [sharp corners] -- (103.3656,158.705);
\path[branch] (94.3206,118.985) [rounded corners={1.5pt}] -- (85.2756,118.985) [sharp corners] -- (85.2756,193.87);
\path[branch] (94.3206,95.1334) -- (94.3206,118.985) -- cycle;
\path[branch] (69.3331,245.1202) -- (69.3331,267.6116) -- cycle;
\path[branch] (130.103,240.1433) -- (130.103,267.6116) -- cycle;
\path[transfer branch] (130.103,228.1433) -- (166.5413,228.1433);\path[branch] (166.5413,228.1433) -- (166.5413,267.61159999);
% gene transfers
% events
\node[speciation] at (31.80071,32.227000000000004) {1, 2};
\node[gain={black}{1, 2}] at (31.80071,4.46088) {};
\node[transfer] at (9.91325,248.47652999999997) {\underline{1}, \underline{2}};
\node[speciation] at (150.23129999999998,68.83699999999999) {1, 2};
\node[speciation] at (103.36556999999999,158.70499999999998) {1, 2, 3};
\node[speciation] at (85.27561,193.87) {2, 3};
\node[duplication] at (94.32059,107.98499999999999) {1, \underline{2}, \underline{3}};
\node[gain={black}{3}] at (94.32059,77.13337999999999) {};
\node[extant={black}{1, 3}] at (69.33313999999999,267.61159999999995) {};
\node[loss={black}{2}] at (69.33313999999999,245.12021999999996) {};
\node[extant={black}{2, 3}] at (58.53816,267.61159999999995) {};
\node[extant={black}{1, 2}] at (47.743179999999995,267.61159999999995) {};
\node[extant={black}{1, 2}] at (130.10302,267.61159999999995) {};
\node[transfer] at (130.10302,228.14329999999998) {1, 2 / 3};
\node[extant={black}{2, 3}] at (119.30804,267.61159999999995) {};
\node[extant={black}{3}] at (166.17377,267.61159999999) {};
\node[extant={black}{1, 2}] at (176.54126,267.61159999999995) {};
\end{tikzpicture}}

    \vspace{1ex}
    \(\Hist\)
    
    (2)
}

    \caption{
        (1)~A \emph{synteny tree} \(\syntenyTree\) (on the left) with an augmented \emph{species tree} \(\augS\) (on the right).
        Numbers represent single gene families and letters represent species.
        As indicated by the dashed gray lines connecting the two trees, the syntenies \(\{\{1,2\}, \{2,3\}, \{1,3\}\}\) belong to species~\textbf{A}, \(\{\{2,3\},\{1,2\}\}\) to~\textbf{B}, and \(\{\{3\},\{1,2\}\}\) to~\textbf{C}.
        Dotted lines in the augmented species tree represent unsampled edges.
        (2)~An output \emph{history} \(\Hist\) of Synesth when given \(\syntenyTree\) and \(\augS\) as input.
        The history tree is represented on top of the species tree drawn in gray in the background.
        The cross-hatched background represents a part of the history taking place in an unsampled lineage (corresponding to the edge to its left in \(\augS\)).
        Events are represented as follows: \enquote{\(\eventSpe\)} by ovals, \enquote{\(\eventDup\)} and \enquote{\(\eventCut\)} by rectangles, \enquote{\(\eventTrDup\)} and \enquote{\(\eventTrCut\)} by diamonds, \enquote{\(\eventLoss\)} by right half-circles and \enquote{\(\eventGain\)} by left half-circles.
        The synteny contents are written inside of each event, while the associated species is represented implicitly by the position of each event on top of the species tree.
        Duplicated or transferred genes are underlined, while fissions are represented by a separation in the synteny.
    }
    \label{fig:synesth:example}
\end{figure}

\section{Explicatory histories and optimization problems}\label{sec:synesth:opt-probs}

The goal of reconciliation is to infer histories which \emph{explicate} the topology of a synteny tree given a tree of the corresponding species.
Such histories are extensions of the synteny tree that map all leaves to the appropriate species and syntenies without introducing new extant leaves.
For example, in \cref{fig:synesth:example} the history in~(2) explicates the trees in~(1).

\begin{definition}[Visible leaves]
    A leaf \(l\) of a history \(\Hist = \tuple{H, \eventOf, \syntenyOf, \speciesOf}\) is said to be \emph{visible} if \(\syntenyOf(l) \neq \emptyset\) and \(\speciesOf(l) \in \treeNodes(\speciesTree)\), and \emph{invisible} otherwise.
    The set of visible leaves of \(\Hist\) is denoted \(\visLeaves(\Hist)\).
\end{definition}

In the history depicted in \cref{fig:synesth:example}, the only invisible leaf is the unnamed leaf in the cross-hatched region (representing an unsampled species). Notice that all losses in this history are partial losses; otherwise a full loss would also correspond to an invisible leaf.

The following is a formal definition of a history explicating a synteny and species tree given as input.

\begin{definition}[Explicatory histories]\label{def:synesth:expl-history}
    For a species tree \(\speciesTree\),
    a history \(\Hist = \tuple{H, \eventOf, \syntenyOf, \speciesOf}\) on \(\augS\) is said to \emph{explicate} a synteny tree~\(\syntenyTree = \tuple{T, \syntenyOf', \speciesOf'}\) and \(\speciesTree\) if:
    \begin{enumerate}
        \item \(H\) is an extension of \(T\).
        \item \(\syntenyOf' = \restr{\syntenyOf}{\treeLeaves(T)}\) and \(\speciesOf' = \restr{\speciesOf}{\treeLeaves(T)}\).
        \item \(\visLeaves(\Hist) = \treeLeaves(T)\).
        \item No \(\treeEdge{u}{v} \in \treeEdges(T)\) is such that \(\speciesOf(v) < \speciesOf(u)\).
        \item\label{def:synesth:expl-history-gain} No gain event is the parent of a node \(v \in \treeNodes(H) - \treeNodes(T)\).
        \item\label{def:synesth:expl-history-loss} No loss event is a child of a node \(v \in \treeNodes(H) - \treeNodes(T)\).
    \end{enumerate}
    The set of all such histories is denoted by \(\solSpace(\syntenyTree, \speciesTree)\).
\end{definition}

Condition (3) excludes histories introducing new visible leaves. With condition (4), we exclude locally cyclic solutions. This condition is a necessary, but not sufficient, condition for acyclicity. Imposing a full acyclicity condition would lead to computationally-intractable problems~\parencite{Tofigh2011SimultaneousIdentificationDuplications}.

As for conditions (5) and (6), they are introduced to avoid redundant histories : two histories with the same events, but with gains and losses distributed differently on branches. More precisely, condition (5) requires sifting gains down and merging them until they are the parent of a synteny tree node; conversely, condition (6) requires sifting losses up and merging them until they are the child of a synteny tree node.

Note that \(\solSpace(\syntenyTree, \speciesTree)\) is infinite: given any explicatory history, it is always possible to \emph{extend} it into a larger one by introducing superfluous duplications or transfers.
We next define a way to reduce this space to a finite one.

\begin{definition}[Event vector]
    Let \(\Hist = \tuple{H, \eventOf, \syntenyOf, \speciesOf}\) be a history.
    We define \(\eventVec(\Hist) = (c_{\eventDup},\allowbreak c_{\eventCut},\allowbreak c_{\eventTrDup},\allowbreak c_{\eventTrCut},\allowbreak c_{\eventLoss}) \in \setNat^5\) as the vector such that \(c_e\) is the number of events of type \(e \in \allEvents\) in \(\Hist\).
\end{definition}

As usual for reconciliation, our definition of an event vector excludes the number of speciations, since they do not allow to meaningfully distinguish histories.
Notice that the number of gains is also excluded, as needed to make the problem of \cref{sec:synesth:pareto-vectors} tractable.
That is, in addition to excluding convergent gains (\cref{def:synesth:hist-gain-once}), we also disregard taking advantage of segmental gains to reduce the number of individual gain events.

\begin{definition}[Order on vectors and histories]
    For two event vectors \((a_i)_{1\leq i \leq n}\) and \((b_i)_{1\leq i \leq n}\), \((a_i) \preceq (b_i)\) if for all \(i \in \{1, ..., n\}\), \(a_i \leq b_i\).
    This partial order induces another one on histories, namely \(\Hist \preceq \Hist'\) if \(\eventVec(\Hist) \preceq \eventVec(\Hist')\).
\end{definition}

\begin{definition}[Pareto optimality]
    For any set \(E\) with a partial order \(\preceq\), we define its Pareto subset as \(
        \paretoMin{\preceq}{E} = \{x \in E \mid \neg (\exists y \in E, y \preceq x \land x \neq y)\}
    \).
\end{definition}

As an example, in the set of vectors \(A = \{(1, 2), (2, 2), (3, 1)\} \subseteq \setNat^2\), we have \((1, 2) \preceq (2, 2)\) while \((1, 2)\) and \((3, 1)\) are not comparable, hence \(\paretoMin{\preceq}{A} = \{(1, 2), (3, 1)\}\).

Let us now consider the set \(\minSolSpace(\syntenyTree, \speciesTree) = \paretoMin{\preceq} \solSpace(\syntenyTree, \speciesTree)\) of histories whose event vectors are Pareto-optimal.
As opposed to \(\solSpace(\syntenyTree, \speciesTree)\), \(\minSolSpace\) is a finite set, as we will show with \cref{thm:synesth:number-min-hist}.
We can thus meaningfully define the problem of computing this set.

\begin{problem}{All Pareto-optimal histories\label{prob:synesth:all-hists}}
    \probinput{A synteny tree \(\syntenyTree\) and a species tree \(\speciesTree\).}
    \proboutput{The set \(\minSolSpace(\syntenyTree, \speciesTree)\) of Pareto-optimal histories explicating \(\syntenyTree\) and \(\speciesTree\).}
\end{problem}

Even though \(\minSolSpace\) is finite, it may contain a number of optimal histories exponential in \(|\treeNodes(T)|\). 
Rather, consider the set \(\minVecSet(\syntenyTree, \speciesTree) = \paretoMin{\preceq} \{\eventVec(\Hist) \mid \Hist \in \solSpace(\syntenyTree, \speciesTree)\}\) of Pareto-optimal event vectors.
As we will show later (in \cref{thm:synesth:all-vecs-complexity}), the number of optimal event vectors in \(\minVecSet\) is polynomial.
We consider the following problem of reduced complexity.

\begin{problem}{All Pareto-optimal vectors\label{prob:synesth:all-vecs}}
    \probinput{A synteny tree \(\syntenyTree\) and a species tree \(\speciesTree\).}
    \proboutput{The set \(\minVecSet(\syntenyTree, \speciesTree)\).}
\end{problem}

Now, given a vector of costs for each event type \(\eventCost = (\eventCost_\eventDup,\allowbreak \eventCost_\eventCut,\allowbreak \eventCost_\eventTrDup,\allowbreak \eventCost_\eventTrCut,\allowbreak \eventCost_\eventLoss) \in (\setReal^+ \cup \{\infty\})^5\), we can associate an overall cost \(\histCost(\Hist) = \eventCost \cdot \eventVec(\Hist)\) to each history.

\begin{problem}{Minimum cost\label{prob:synesth:min-cost}}
    \probinput{A synteny tree \(\syntenyTree\), a species tree \(\speciesTree\), and a vector \(\eventCost \in (\setReal^+ \cup \{\infty\})^5\).}
    \proboutput{The minimum cost \(\minCost(\syntenyTree, \speciesTree) = \min \{ \histCost(\Hist) \mid \Hist \in \solSpace(\syntenyTree, \speciesTree) \}\) of any history\newline explicating \(\syntenyTree\) and \(\speciesTree\).}
\end{problem}

Finally, we will call problems~\ref*{prob:synesth:all-vecs}' and~\ref*{prob:synesth:min-cost}' the versions of \cref{prob:synesth:all-vecs,prob:synesth:min-cost} where we additionally ask for one history corresponding to each returned optimal event vector (for \cref*{prob:synesth:all-vecs}') or one history for the returned minimum cost (for \cref*{prob:synesth:min-cost}').
For example, \cref{fig:synesth:example} shows in~(2) a possible solution for \cref*{prob:synesth:all-vecs}' for the synteny and species trees in~(1) for vector \((c_\eventDup = 1, c_\eventCut = 0, c_\eventTrDup = 1, c_\eventTrCut = 1, c_\eventLoss = 1)\), and for \cref*{prob:synesth:min-cost}' for cost \((\eventCost_\eventDup = 2, \eventCost_\eventCut = 2.5, \eventCost_\eventTrDup = 3, \eventCost_\eventTrCut = 3.5, \eventCost_\eventLoss = 1)\).

\section{Generating all pareto-optimal histories}\label{sec:synesth:pareto-histories}

We start by addressing \cref{prob:synesth:all-hists}, which asks to enumerate the \(\minSolSpace\) set.
This can be done inductively by building up histories from the extant leaves of the synteny tree up to its root.
This result is the basis for the dynamic programming formulations introduced for the other problems in the upcoming sections.

In the next two definitions, we introduce the notion of partial histories and a composition operator which is then used to build complete histories (see \cref{fig:synesth:history-builder}).
In the following, a node \(v\) of a history such that \(\syntenyOf(v) = X\) and \(\speciesOf(v) = \sigma\) is denoted \(\makeNode{X, \sigma}\), or \(\makeNode{e, X, \sigma}\) if the event \(e \in \allEvents\) associated to \(v\) is known.
If \(A\), \(B\) and \(C\) are three such nodes, we write \((A, B)C\) to refer to the binary tree with root \(C\) and \(A\) and \(B\) as leaves.

\begin{definition}[Partial histories]
    Let \(\syntenyTree = \tuple{T, \syntenyOf, \speciesOf}\) be a synteny tree on a species tree \(\speciesTree\).
    Let \(v \in \treeNodes(T)\), \(X \subseteq \geneFamilySet\), and \(\sigma \in \treeNodes(\augS)\).
    We define \(h(v, X, \sigma)\) as the set of minimal histories explicating the subtree of \(\syntenyTree\) rooted at \(v\) with \(v = [X, \sigma]\).
    We also define \(\optPath([X, \sigma], [Y, \gamma])\) as the set of minimal acyclic histories whose root is \([X, \sigma]\) and whose only visible leaf is \([Y, \gamma]\).
    Formally,
    \begin{align*}
        h(v, X, \sigma) = \{
            \tuple{H, \eventOf, \syntenyOf, \speciesOf}
                \in \minSolSpace(\syntenyTree[v], \speciesTree) \mid
            \begin{aligned}[t]
                & \rootOf(H) = v \ \land \syntenyOf(v) = X \land \speciesOf(v) = \sigma \}
            \end{aligned}
        \\
        \optPath([X, \sigma], [Y, \gamma]) = \paretoMin{\preceq}\{
            \Hist = \tuple{H, \eventOf, \syntenyOf, \speciesOf} \mid
            \begin{aligned}[t]
                & \visLeaves(\Hist) = \{l\} \land
                \syntenyOf(\rootOf(H)) = X \land \syntenyOf(l) = Y \\
                & \land \ \speciesOf(\rootOf(H)) = \speciesOf(l) = \sigma \land
                \Hist \text{ is acyclic} \}.
            \end{aligned}
    \end{align*}
\end{definition}

\newcommand{\mergeHist}{\mathcal{M}}

\begin{definition}[History composition]
    If \(\Hist = \tuple{H, \eventOf, \syntenyOf, \speciesOf}\)
    and \(\Hist' = \tuple{H', \eventOf', \syntenyOf', \speciesOf'}\) are two histories 
    such that there exists a leaf \(l \in \treeLeaves(H)\) with \(\syntenyOf(l) = \syntenyOf(\rootOf(H'))\) and \(\speciesOf(l) = \speciesOf(\rootOf(H'))\),
    then we define their \emph{composition} \(\Hist_1 \otimes \Hist_2\) as the history obtained by replacing \(l\) with \(\rootOf(H')\) in \(H\) and merging \(\syntenyOf\) with \(\syntenyOf'\), and \(\speciesOf\) with \(\speciesOf'\).
    When applied on two \emph{sets} of histories, this operation is defined as taking the cartesian product of the two sets and composing each obtained history pair together.

    For any \(e \in \allEvents\), \(X, Y, Y', Z, Z' \subseteq \geneFamilySet\) and \(\sigma, \gamma\leftChild, \gamma\leftChild', \gamma\rightChild, \gamma\rightChild' \in \treeNodes(\augS)\) we define:
    \begin{align*}
        &\mergeHist([e, X, \sigma], [Y', \gamma\leftChild'], [Y, \gamma\leftChild], [Z', \gamma\rightChild'], [Z, \gamma\rightChild])\\
        &= ([Y', \gamma\leftChild'], [Z', \gamma\rightChild'])[e, X, \sigma] \ \begin{aligned}[t]
            \otimes \ (\optPath([Y', \gamma\leftChild'], [Y, \gamma\leftChild])
            \otimes h(v\leftChild, Y, \gamma\leftChild)) \\
            \otimes \ (\optPath([Z', \gamma\rightChild'], [Z, \gamma\rightChild])
            \otimes h(v\rightChild, Z, \gamma\rightChild)).
        \end{aligned}
    \end{align*}
\end{definition}

For example, in \cref{fig:synesth:history-builder},
\(\Hist = \mergeHist([\eventSpe, X, \sigma], [X, \sigma\leftChild], [Y, \gamma\leftChild], [X, \sigma\rightChild], [Z, \gamma\rightChild])\).

\begin{figure}[htb]
    \centering
    \sidecols{.49\textwidth}{.49\textwidth}{
    \centering
    \begin{tikzpicture}[
        node distance=1.5cm and 1.2cm,
        label distance=0cm,
        node label/.style={font=\scriptsize},
    ]
        \coordinate[label={[node label]left:\(v\)}] (root) {};

        \coordinate [
            below left=of root,
            label={[node label]left:\(v\lchild\)},
        ] (node-left) {};

        \node [node label, below=.7cm of node-left]
            (node-left-subtree)
            {\(\subtree{\syntenyTree}{v\lchild}\)};

        \coordinate [
            below right=of root,
            label={[node label]right:\(v\rchild\)},
        ] (node-right) {};

        \node [node label, below=.7cm of node-right]
            (node-right-subtree)
            {\(\subtree{\syntenyTree}{v\rchild}\)};

        \draw (root) edge (node-left);
        \draw (root) edge (node-right);

        \path[draw, fill=white] (node-left.south)
            -- ([shift={(-.33, -.66)}]node-left.south)
            -- ([shift={(.33, -.66)}]node-left.south)
            -- cycle;

        \path[draw, fill=white] (node-right.south)
            -- ([shift={(-.33, -.66)}]node-right.south)
            -- ([shift={(.33, -.66)}]node-right.south)
            -- cycle;
    \end{tikzpicture}
    
    \(\syntenyTree\)
}{
    \centering
    \begin{tikzpicture}[
        node distance=.4cm and .2cm,
        label distance=0cm,
        node label/.style={font=\scriptsize},
    ]
        \node[
            speciation,
            label={[node label]left:\(v\)},
            label={[species label]right:\(\sigma\)},
        ] (root) {\(X\)};

        \node [
            generic,
            below left=of root,
            label={[species label]right:\(\sigma\lchild\)},
        ] (path-left) {\(X\)};

        \node [
            generic,
            below left=of path-left,
            label={[node label]left:\(v\lchild\)},
            label={[species label]right:\(\gamma\lchild\)},
        ] (node-left) {\(Y\)};

        \node [node label, below=.7cm of node-left]
            (node-left-subtree)
            {\(h(v\lchild, Y, \gamma\lchild)\)};

        \node [
            generic,
            below right=of root,
            label={[species label]right:\(\sigma\rchild\)},
        ] (path-right) {\(X\)};

        \node [
            generic,
            below right=of path-right,
            label={[node label]left:\(v\rchild\)},
            label={[species label]right:\(\gamma\rchild\)},
        ] (node-right) {\(Z\)};

        \node [node label, below=.7cm of node-right]
            (node-right-subtree)
            {\(h(v\rchild, Z, \gamma\rchild)\)};

        \draw (root) edge (path-left);
        \draw (root) edge (path-right);

        \draw[wavy edge] (path-left) edge (node-left);
        \draw[wavy edge] (path-right) edge (node-right);

        \path[draw, fill=white] (node-left.south)
            -- ([shift={(-.33, -.66)}]node-left.south)
            -- ([shift={(.33, -.66)}]node-left.south)
            -- cycle;

        \path[draw, fill=white] (node-right.south)
            -- ([shift={(-.33, -.66)}]node-right.south)
            -- ([shift={(.33, -.66)}]node-right.south)
            -- cycle;
    \end{tikzpicture}

    \(\Hist\)
}

    \caption{
        A history \(\Hist\) for a synteny tree \(\syntenyTree\) starting with a speciation at \(v = [X, \sigma]\), composed of two sub-histories \(h(v\leftChild, Y, \gamma\leftChild)\) and \(h(v\rightChild, Z, \gamma\rightChild)\)
        linked together by two paths (represented by wavy lines):
        \(\optPath([Y', \gamma\leftChild'], [Y, \gamma\leftChild])\)
        and \(\optPath([Z', \gamma\rightChild'], [Z, \gamma\rightChild])\).
    }
    \label{fig:synesth:history-builder}
\end{figure}

\pagebreak
This representation of histories as compositions of sub-histories and paths is now used to formulate the inductive definitions of \(h(v, X, \sigma)\) and \(\optPath\), and ultimately \(\minSolSpace(\syntenyTree, \speciesTree)\).
In the following, we use \(A \oplus B\) for \(\min(A \cup B)\).

\begin{theorem}[Inductive form of minimal histories]\label{thm:synesth:minimal-history}
    Let \(\syntenyTree = \tuple{T, \syntenyOf, \speciesOf}\) be a synteny tree on~\(\speciesTree\),
    \(v\) be an internal node of~\(T\),
    \(X \subseteq \geneFamilySet\), and
    \(\sigma\) be a node of \(\augS\).
    Then \(h(v, X, \sigma) = P_{\eventSpe} \oplus P_{\eventDup} \oplus P_{\eventTrDup} \oplus P_{\eventCut} \oplus P_{\eventTrCut}\), with:

    \begin{align*}
        P_{\eventSpe} = \bigoplus_{\substack{
            \gamma\leftChild, \gamma\rightChild \not< \sigma \\
            Y, Z \subseteq \geneFamilySet
        }} & \begin{aligned}[t]
            &\big(
                \mergeHist([\eventSpe, X, \sigma], [X, \sigma\leftChild], [Y, \gamma\leftChild], [X, \sigma\rightChild], [Z, \gamma\rightChild]) \\
            &\quad\oplus\,
                \mergeHist([\eventSpe, X, \sigma], [X, \sigma\rightChild], [Y, \gamma\leftChild], [X, \sigma\leftChild], [Z, \gamma\rightChild])\big)
        \end{aligned}
    \end{align*}

    \begin{align*}
        P_{\eventDup} = \bigoplus_{\substack{
            \gamma\leftChild, \gamma\rightChild \not< \sigma \\
            Y, Z \subseteq \geneFamilySet
        }} & \begin{aligned}[t]
            &\big(
                \mergeHist([\eventDup, X, \sigma], [X \cap Y, \sigma], [Y, \gamma\leftChild], [X, \sigma], [Z, \gamma\rightChild]) \\
            &\quad\oplus\,
                \mergeHist([\eventDup, X, \sigma], [X, \sigma], [Y, \gamma\leftChild], [X \cap Z, \sigma], [Z, \gamma\rightChild])\big)
        \end{aligned}
    \end{align*}

    \begin{align*}
        P_{\eventCut} = \bigoplus_{\substack{
            \gamma\leftChild, \gamma\rightChild \not< \sigma \\
            Y, Z \subseteq \geneFamilySet
        }} & \begin{aligned}[t]
            &\big(
                \mergeHist([\eventCut, X, \sigma], [X \cap Y, \sigma], [Y, \gamma\leftChild], [X - Y, \sigma], [Z, \gamma\rightChild]) \\
            &\quad\oplus\,
                \mergeHist([\eventCut, X, \sigma], [X - Z, \sigma], [Y, \gamma\leftChild], [X \cap Z, \sigma], [Z, \gamma\rightChild])\big)
        \end{aligned}
    \end{align*}
    
    \begin{align*}
        P_{\eventTrDup} = \bigoplus_{\substack{
            \gamma_i \separ \sigma,
            \gamma\targetChild \not< \gamma_i \\
            \gamma\conservChild \not< \sigma;~
            Y, Z \subseteq \geneFamilySet
        }} & \begin{aligned}[t]
            &\big(
                \mergeHist([\eventTrDup, X, \sigma], [X \cap Y, \gamma_i], [Y, \gamma\targetChild], [X, \sigma], [Z, \gamma\conservChild]) \\
            &\quad\oplus\,
                \mergeHist([\eventTrDup, X, \sigma], [X, \sigma], [Y, \gamma\conservChild], [X \cap Z, \gamma_i], [Z, \gamma\targetChild])\big)
        \end{aligned}
    \end{align*}

    \begin{align*}
        P_{\eventTrCut} = \bigoplus_{\substack{
            \gamma_i \separ \sigma,
            \gamma\targetChild \not< \gamma_i \\
            \gamma\conservChild \not< \sigma;~
            Y, Z \subseteq \geneFamilySet
        }} & \begin{aligned}[t]
            &\big(
                \mergeHist([\eventTrCut, X, \sigma], [X \cap Y, \gamma_i], [Y, \gamma\targetChild], [X - Y, \sigma], [Z, \gamma\conservChild]) \\
            &\quad\oplus\,
                \mergeHist([\eventTrCut, X, \sigma], [X - Z, \gamma_i], [Y, \gamma\targetChild], [X \cap Z, \sigma], [Z, \gamma\conservChild]) \\
            &\quad\oplus\,
                \mergeHist([\eventTrCut, X, \sigma], [X \cap Y, \sigma], [Y, \gamma\conservChild], [X - Y, \gamma_i], [Z, \gamma\targetChild]) \\
            &\quad\oplus\,
                \mergeHist([\eventTrCut, X, \sigma], [X - Z, \sigma], [Y, \gamma\conservChild], [X \cap Z, \gamma_i], [Z, \gamma\targetChild])\big)
        \end{aligned}
    \end{align*}
\end{theorem}
\begin{proof}
    First, notice that all histories of the \(P_e\) sets,
    for \(e \in \{\eventSpe,\allowbreak \eventDup,\allowbreak \eventCut,\allowbreak \eventTrDup,\allowbreak \eventTrCut\}\),
    explicate \(\syntenyTree[v]\) and \(\speciesTree\),
    have their root assigned to \(X\) and \(\sigma\)
    and are minimal.

    Let \(\Hist = \tuple{H, \eventOf, \syntenyOf, \speciesOf} \in h(v, X, \sigma)\).
    Since \(v\) is an internal node of \(T\), \(v\) must be binary in \(H\), hence \(\eventOf(v) \in \{\eventSpe, \eventDup, \eventCut, \eventTrDup, \eventTrCut\}\).
    If \(\eventOf(v) \in \{\eventSpe, \eventDup, \eventCut\}\), \(\gamma\leftChild\) and \(\gamma\rightChild\) cannot be strict ancestors of \(\sigma\), as otherwise \(\Hist\) would not be acyclic.
    Let \(Y = \syntenyOf(v\leftChild)\),
    \(Y' = \syntenyOf(w\leftChild)\),
    \(Z = \syntenyOf(v\rightChild)\),
    \(Z' = \syntenyOf(w\rightChild)\).

    If \(\eventOf(v) = \eventSpe\),
        \(\Hist \in P_{\eventSpe}\)
        by \cref{def:synesth:hist}~\cref{def:synesth:hist-spe}.

    Otherwise,
        let us first show that the synteny contents of both children are chosen appropriately.
        We have that \(X \cap Y \subseteq Y'\) (resp. \(X \cap Z \subseteq Z'\))
        since if any \(g \in X \cap Y\) (resp. \(g \in X \cap Z\))
        was not in \(Y'\) (resp. \(Z'\)),
        we would have multiple gains of \(g\) in \(H\),
        one above \(v\) (since \(g \in X\))
        and one between \(v\) and \(v_l\) (since \(g \in Y\), resp. \(v_r\) since \(g \in Z\)).
    \begin{itemize}
        \item If \(\eventOf(v) \in \{\eventDup, \eventTrDup\}\),
            by \cref{def:synesth:hist},
            at least one of \(Y'\) or \(Z'\) is equal to \(X\).
            Without loss of generality (w.l.o.g.), assume that \(Z' = X\).
            We should have \(Y' \subseteq X \cap Y\)
            as otherwise we would have at least one additional event in the path between \(v\) and \(v_l\)
            for loosing \(Y' - (X \cap Y)\),
            contradicting the minimality of \(\Hist\).
            Hence \(Y' = X \cap Y\).

        \item If \(\eventOf(v) \in \{\eventCut, \eventTrCut\}\),
            it cannot be that both \(Y' \not\subseteq X \cap Y\) and \(Z' \not\subseteq X \cap Z\) are true
            as such a scenario would lead to a loss on each path between \(v\) and \(v\leftChild\) and between \(v\) and \(v\rightChild\),
            whereas by choosing \(Y' \subseteq X \cap Y\) or \(Z' \subseteq X \cap Z\) we save at least one loss.
            Since \(Y'\) and \(Z'\) have to be a partition of \(X\), then \(Z' = X - Y\) if \(Y' = X \cap Y\), or \(Y' = X - Z\) otherwise.
    \end{itemize}

    We now show that the species of both children are also chosen appropriately.
    \begin{itemize}
        \item If \(\eventOf(v) \in \{\eventDup, \eventCut\}\), then by \cref{def:synesth:hist},
            \(\speciesOf(w\leftChild) = \speciesOf(w\rightChild) = \sigma\).

        \item If \(\eventOf(v) \in \{\eventTrDup, \eventTrCut\}\), 
            then by \cref{def:synesth:hist}
            either \(\speciesOf(w\leftChild) = \sigma\) and \(\speciesOf(w\rightChild) \separ \sigma\),
            or \(\speciesOf(w\rightChild) = \sigma\) and \(\speciesOf(w\leftChild) \separ \sigma\).

        \item If \(\eventOf(v) = \eventTrDup\),
            if \(\speciesOf(w\leftChild) = \sigma\) (resp. \(\speciesOf(w\rightChild) = \sigma\))
            then \(\syntenyOf(w\leftChild) = X\) (resp. \(\syntenyOf(w\rightChild) = Y\)). \qedhere
    \end{itemize}
\end{proof}

We next consider how to compute the set of histories \(\optPath([X, \sigma], [Y, \gamma])\) given any two syntenies \(X\) and \(Y\) and species \(\sigma\) and \(\gamma\).

\begin{theorem}[Minimal path histories]\label{thm:synesth:minimal-paths}
    Let \(X\) and \(Y\) be two syntenies and \(\sigma\) and \(\gamma\) be two species.
    If \(\gamma\) is a strict ancestor of \(\sigma\), then \(\optPath([X, \sigma], [Y, \gamma]) = \emptyset\).
    Otherwise, \(\optPath([X, \sigma], [Y, \gamma]) = P_0 \oplus P_1 \oplus P_2\), where:

    \begin{itemize}
        \item \(P_0 = \emptyset\) if \(\sigma \separ \gamma\).
            Otherwise, \(P_0\) contains the history made up of a chain of \(\treeDist[\augS]{\sigma}{\gamma}\) speciations and \(|\{(u, v) \in \treePath[\augS]{\sigma}{\gamma} \mid v \notin \treeNodes(\speciesTree)\}|\) full losses, preceded by a loss event if \(X \not\subseteq Y\) and followed by a gain event if \(Y \not\subseteq X\).
            If \(\gamma \in \treeLeaves(\augS) - \treeLeaves(\speciesTree)\), then \(P_0\) also includes a history where the initial partial loss event is replaced by a duplication and an unsampled leaf.
        \item \(P_1 = P_\eventTrDup \oplus  P_\eventTrCut^{\mathrm{Left}} \oplus P_\eventTrCut^{\mathrm{Right}}\).
            If \(\sigma \separ \gamma\), all three sets contain histories with a \(\eventTrAll\) event from \(\sigma' = \sigma\) to \(\gamma\).
            Otherwise, all three sets contain histories with a speciation at \(\sigma\), followed by a \(\eventTrAll\) event \(v\) from \(\sigma'\) (such that \(\sigma' \in \children(\sigma)\) and \(\sigma' \separ \gamma\)) to \(\gamma\).
            In both cases, the histories end with a gain event if \(Y \not\subseteq X\).
            For \(\Hist \in P_\eventTrDup\), \(\eventOf(v) = \eventTrDup\) and \(\syntenyOf(v\targetChild) = X \cap Y\), with a full loss at \(v\conservChild\) if \(\sigma' \notin \treeLeaves(\augS)\).
            For \(\Hist \in P_\eventTrCut^{\mathrm{Left}}\) or \(\Hist \in P_\eventTrCut^{\mathrm{Right}}\), \(\eventOf(v) = \eventTrCut\).
            For \(\Hist \in P_\eventTrCut^{\mathrm{Left}}\), \(\syntenyOf(v\targetChild) = X \cap Y\) and \(v\conservChild\) is a full loss if \(X \not\subseteq Y\) and \(\sigma' \notin \treeLeaves(\augS)\).
            For \(\Hist \in P_\eventTrCut^{\mathrm{Right}}\), \(\syntenyOf(v\conservChild) = \emptyset\) and there is a partial loss at \(v\targetChild\) if \(X \not\subseteq Y\).
        \item \(P_2 = \emptyset\) if \(\sigma \separ \gamma\) and \(X \subseteq Y\), or if \(\sigma \in \treeLeaves(\augS)\).
            Otherwise, \(P_2 = P^\eventTrDup_\eventTrDup \oplus P^\eventTrCut_\eventTrDup \oplus P^\eventTrCut_\eventTrCut\).
            Each of the three sets \(P^{e}_{e'}\) contain histories made up of two consecutive transfers \(v\) and \(v'\) such that \(\eventOf(v) = e\) and \(\eventOf(v') = e'\).
            The histories of \(P^\eventTrDup_\eventTrDup\) and \(P^\eventTrCut_\eventTrDup\) are such that \(\speciesOf(v') \in \treeLeaves(\augS)\).
            In \(P^\eventTrCut_\eventTrDup\) and \(P^\eventTrCut_\eventTrCut\), the first cut \(v\) is complete (i.e. with an empty leaf).
            In \(P^\eventTrCut_\eventTrCut\) if \(X \not\subseteq Y\) then \(\speciesOf(v) \in \treeLeaves(\augS)\), otherwise the second cut \(v'\) is complete if \(\speciesOf(v) \notin \treeLeaves(\augS)\).
            In all cases, the histories end with a gain event if \(Y \not\subseteq X\).
    \end{itemize}

    \threecols{.25\textwidth}{.45\textwidth}{.25\textwidth}{
        \centering
        \begin{tikzpicture}[x=.8cm, y=.5cm]
            \node[sideways loss={black}{\(-(X - Y)\)}, label={[species label]left:\(\sigma\)}] (s0) at (0, 1) {};
            \node[speciation, label={[species label]left:\(\sigma\)}] (s1) at (0, 0) {};
            \node[speciation] (s2) at (0, -1) {};
            \coordinate (s3) at (0, -2);
            \node[extant gene={\(Y\)}, label={[species label]left:\(\gamma\)}] (l3) at (0, -3) {};

            \path (s0) edge[draw] (s1);
            \path (s1) edge[draw] (s2);
            \path (s2) edge[draw, dotted] (s3);
            \path (s3) edge[draw] (l3);

            \coordinate (l1) at (1, 0);
            \path (s1) edge[draw] node[pos=1, loss={black}{}] {} (l1);
            \coordinate (l2) at (1, -1);
            \path (s2) edge[draw, dashed] (l2);

            \node[sideways gain={black}{\(+(Y - X)\)}] (p4) at (0, -2.5) {};
        \end{tikzpicture}

        \((P_0)\)
    }{
        \centering
        \begin{tikzpicture}[x=.3cm, y=.5cm]
            \node[transfer, label={[species label,xshift=1.5em]above left:\(\sigma = \sigma'\)}]
                (s1) {\(X\)};
            \coordinate (l1) at (0, -1) {};
            \coordinate (s2) at (3, 0);
            \node[extant gene={\(Y\)}, label={[species label]right:\(\gamma\)}]
                (l2) at (3, -1) {};
            \path (s1) edge[draw] node[pos=1, sideways loss={black}{}] {} (l1);
            \path (s2) edge[draw] (l2.center);
            \path (s1) edge[transfer branch, bend left=20] (s2);
        \end{tikzpicture}
        \begin{tikzpicture}[x=.3cm, y=.5cm]
            \node[speciation, label={[species label]left:\(\sigma\)}]
                (s1) {\(X\)};
            \coordinate (l1) at (1, -1);
            \node[transfer, label={[species label]left:\(\sigma'\)}]
                (s2) at (-1, -1) {\(X\)};
            \coordinate (l2) at (-2, -2);
            \coordinate (s3) at (1, -2.25);
            \node[extant gene={\(Y\)}, label={[species label]right:\(\gamma\)}]
                (l3) at (1, -3) {};

            \path (s1) edge[draw] node[pos=1, sloped, loss={black}{}] {} (l1);
            \path (s1) edge[draw] (s2);
            \path (s3) edge[draw] node[pos=.4, sideways loss={black}{}, label={right:\tiny\(\scriptstyle- (X - Y)\)}] {} (l3);
            \path (s2) edge[draw] node[pos=1, sloped, rotate=180, loss={black}{}] {} (l2);
            \path (s2) edge[transfer branch, bend left=10] (s3);
        \end{tikzpicture}

        \vspace{1ex}
        \((P_1)\)
    }{
        \centering

        \begin{tikzpicture}[x=.3cm, y=.5cm]
            \node[transfer, label={[species label]above left:\(\sigma\)}]
                (s1) {\(X\)};
            \node[sideways loss={black}{}] (l1) at (0, -1) {};
            \node[transfer, label={[species label]above:\(\sigma^*\)}]
                (s2) at (3, 0) {\(X\)};
            \node[extant gene={\(X\)}, label={[species label, label distance=-.3em]left:\(\sigma^*\)}]
                (l2) at (3, -1) {};
            \coordinate (s3) at (5, 0);
            \node[extant gene={\(Y\)}, label={[species label]right:\(\gamma\)}]
                (l3) at (5, -1) {};
            \path (s1) edge[draw] (l1);
            \path (s2) edge[draw] (l2.center);
            \path (s1) edge[transfer branch, bend left=20] (s2);
            \path (s3) edge[draw] (l3.center);
            \path (s2) edge[transfer branch, bend left=25] (s3);
        \end{tikzpicture}%

        \vspace{2ex}
        \((P_2)\)
    }
\end{theorem}
\begin{proof}
    If \(\gamma\) is a strict ancestor of \(\sigma\), there can be no acyclic history leading from \((X, \sigma)\) to \((Y, \gamma)\), therefore \(\optPath((X, \sigma), (Y, \gamma)) = \emptyset\).

    Otherwise, let \(\Hist = \tuple{H, \eventOf, \syntenyOf, \speciesOf} \in \optPath((X, \sigma), (Y, \gamma))\), \(u = \rootOf(H)\), and \(w\) be the only visible leaf of \(\Hist\).
    We call \(\treePath{u}{w}\) the \emph{main path} of \(\Hist\).
    We say that a subtree \(\subtree{H}{v}\) of \(H\) is \emph{invisible} if it contains only invisible leaves.

    First notice that all binary nodes of \(\Hist\) must be on the main path.
    In fact, if \(v\) is a binary node that is not on the main path, then both subtrees of \(v\) must be invisible.
    If \(\speciesOf(v) \notin \treeLeaves(\augS)\), then at least one subtree contains a full loss, hence we can replace \(v\) with a full loss.
    Otherwise, all subtrees contain only unsampled leaves, hence we can replace \(v\) with an unsampled leaf.
    Also notice that all binary nodes of the main path have exactly one child whose subtree is invisible.

    We show now that \(\Hist\)
    may only contain at most two \(\eventTrAll\) nodes.
    Assume that \(\Hist\) contains at least three consecutive \(\eventTrAll\) events that we denote \(v_1\), \(v_2\), \(v_3\).
    Let \(v_4\) be the child of \(v_3\) resulting from the transfer.
    Consider the history \(\Hist'\) obtained from \(\Hist\) by
    (1) removing all the nodes on the path from \(v_2\) to \(v_4\) and their subtrees, including in particular \(v_3\),
    (2) connecting \(v_2\) directly to \(v_4\),
    (3) remapping \(\speciesOf(v_2)\) to any unsampled species separated from \(\speciesOf(v_1)\) and \(\speciesOf(v_4)\), and
    (4) replacing the subtree below \(v_2\) by a single leaf \(l\) with \(\speciesOf(l) = \speciesOf(v_2)\) and \(\syntenyOf(l) = \syntenyOf(v_2)\) if \(\eventOf(v_2) = \eventTrDup\) or \(\syntenyOf(l) = \syntenyOf(v_2) - \syntenyOf(v_4)\) otherwise.
    Clearly, \(\Hist' \histExt \Hist\), \(\Hist' \neq \Hist\) and \(\Hist' \in \optPath((X, \sigma), (Y, \gamma))\), which contradicts the minimality of \(\Hist\).

    As per \cref{def:synesth:expl-history}~\cref{def:synesth:expl-history-gain,def:synesth:expl-history-loss}, any partial loss must be placed at the end of the history, and any gain must be placed at the start.
    Note that \(\Hist\) contains either zero (\(P_0\)), one (\(P_1\), \(P_1^*\), \(P_{1A}\) or \(P_{1A}^*\)) or two (\(P_2\)) transfers.

    \begin{itemize}
        \item {[\emph{\(\Hist\) contains no transfers.}]}
            Note that \(\sigma\) cannot be separated from \(\gamma\) since only transfers can reach separate species, thus \(\sigma\) is an ancestor of \(\gamma\).
            We start by showing that \(\Hist\) can only possibly contain a duplication if that duplication is on an unsampled leaf species.
            If \(v\) is a duplication node such that \(\speciesOf(v) \notin \treeLeaves(\augS)\), then the invisible subtree of \(v\) contains at least one full loss.
            In this case, we can remove the invisible subtree and turn \(v\) into a partial loss, if needed, otherwise replace \(v\) with its remaining child.
            Hence, if \(\Hist\) contains a duplication, it must be the last binary event on the main path.
            Thus all other binary events must be speciations, and we need exactly \(\treeDist[\augS]{\sigma}{\gamma}\) of them to reach \(\sigma\) from \(\gamma\).
            Those speciations must lead to one full loss for each species \(s \in \treeNodes(\speciesTree)\).
            Hence \(\Hist \in P_0\).

        \item {[\emph{\(\Hist\) contains exactly one transfer.}]}
            Let \(v\) be the only transfer node.
            If \(\sigma \leq \gamma\), then there must be a speciation above \(v\) so that \(\speciesOf(v)\) can be separated from~\(\gamma\), since \(\Hist\) cannot contain other transfers.
            Notice that \(\Hist\) contains no partial losses, duplications or cuts.
            In fact, any partial loss can be merged into the transfer.
            As for duplications and cuts, they can be removed if they do not save any loss, or merged into the transfer otherwise.
            Apart from the initial speciation if \(\sigma \leq \gamma\), \(\Hist\) contains no other speciations as those before (resp. after) the transfer can be removed by redirecting the transfer to start from a higher species (resp. to end at a lower species) that is still separated from \(\gamma\) (resp. \(\sigma\)).
            If \(\eventOf(v) = \eventTrDup\) or \(X \not\subseteq Y\), then the invisible subtree of \(v\) must contain at least one full loss, unless \(\speciesOf(v) \in \treeLeaves(\augS)\).
            If \(\speciesOf(v) \notin \treeLeaves(\augS)\) and \(\eventOf(v) = \eventTrCut\), then either subtree of \(v\) must contain at least one full or partial loss if \(\syntenyOf(v) \not\subseteq Y\).
            Hence \(\Hist \in P_1\).

        \item {[\emph{\(\Hist\) contains exactly two transfers.}]}
            Let \(v\) and \(v'\) be the only two transfer nodes.
            If \(\sigma \in \treeLeaves(\augS)\), then if \(\sigma \leq \gamma\), we can remove both transfers, and if \(\sigma \separ \gamma\) then we can remove \(v'\), in both cases without introducing additional events.
            If \(\sigma \separ \gamma\):
                if \(\eventOf(v) = \eventTrDup\) we can remove \(v'\) without adding new losses;
                if \(X \subseteq Y\) then \(\syntenyOf(v) \subseteq \syntenyOf(v') \subseteq \syntenyOf(w)\), hence we can also remove \(v'\).
            If either of \(v\) or \(v'\) is such that \(\eventOf(v) = \eventTrCut\) or \(\eventOf(v') = \eventTrCut\), then \(\eventOf(v) = \eventTrCut\), as otherwise \(\eventOf(v) = \eventTrDup\) and there is a full loss below \(v\), and we can exchange the transfer types of \(v\) and \(v'\) so that \(\eventOf(v) = \eventTrCut\), \(\eventOf(v') = \eventTrDup\) and \(\speciesOf(v') \in \treeLeaves(\augS)\) so as to save a full loss.
            If \(\syntenyOf(v') \not\subseteq Y\) or \(\eventOf(v') = \eventTrDup\), then \(\speciesOf(v') \in \treeLeaves(\augS)\) as otherwise we can reroute the first transfer towards an unsampled leaf so as to save a full loss. \qedhere
    \end{itemize}
\end{proof}

The set of all possible histories \(\minSolSpace(\syntenyTree, \speciesTree)\) can then be computed as the union of assignments of the root to all possible syntenies and species, starting with a path from an empty synteny (i.e. an initial gain).
In other words, \[\minSolSpace(\syntenyTree, \speciesTree) = \bigoplus_{\substack{X \subseteq \geneFamilySet\\\sigma \in \treeNodes(\augS)}}{\optPath([\emptyset, \sigma], [X, \sigma]) \otimes h(\rootOf(\syntenyTree), X, \sigma))}.\]

Hence, using \cref{thm:synesth:minimal-history,thm:synesth:minimal-paths}, one can derive a dynamic programming algorithm to solve \cref{prob:synesth:all-hists}.
Due to the exponential size of the set of solutions, the time complexity of that algorithm is exponential.

\begin{theorem}[Number of minimal histories]\label{thm:synesth:number-min-hist}
    Let \(\syntenyTree = \tuple{T, \syntenyOf, \speciesOf}\) be a synteny tree on a species tree \(\speciesTree\) with gene family set \(\geneFamilySet\).
    Then, \[
        |\minSolSpace(\syntenyTree, \speciesTree)| \in
            \mathcal{O}\left(
                (2^{|\geneFamilySet|}|\treeNodes(\speciesTree)|^3)^{|\treeNodes(T)|}
            \right).
    \]
\end{theorem}
\begin{proof}
    Using \cref{thm:synesth:minimal-history}, we obtain that \(|h(v, X, \sigma)| \in \mathcal{O}(f(v))\) with \[
        f(v) = \begin{cases}
            1 & \text{if } v \text{ is a leaf} \\
            (2^{|\geneFamilySet|} |\treeNodes(\speciesTree)|^2 p) \times f(v\leftChild) \times f(v\rightChild) & \text{otherwise},
        \end{cases}
    \]
    where \(p\) is an upper bound on the number of possible paths.
    This is because in each case of \cref{thm:synesth:minimal-history}, up to all subsets \(Y, Z \subseteq \geneFamilySet\) are tried, along with up to all possible species pairs \(\gamma\leftChild, \gamma\rightChild \in \treeNodes(\augS)^2\), which has a number of nodes directly proportional to \(\speciesTree\).
    Using \cref{thm:synesth:minimal-paths}, we get that \(p \in \mathcal{O}(|\treeNodes(\speciesTree)|)\).
    We get the desired result by solving the recurrence.
\end{proof}

\section{Polynomial-time computation of Pareto-optimal event vectors}\label{sec:synesth:pareto-vectors}

We now address \cref{prob:synesth:all-vecs}.
Given a synteny tree \(\syntenyTree = \tuple{T, \syntenyOf, \speciesOf}\) on a species tree~\(\speciesTree\),
similarly to the way \(h(v, X, \sigma)\) was previously used to compute \(\minSolSpace(\syntenyTree, \speciesTree)\), we define \(\Lambda(v, X, \sigma)\) for computing \(\minVecSet(\syntenyTree, \speciesTree)\)
as \(\Lambda(v, X, \sigma) = \{\eventVec(\Hist) \mid \Hist \in h(v, X, \sigma)\}\).

We can adapt the formulation of \cref{thm:synesth:minimal-history} to compute \(\Lambda(v, X, \sigma)\) using the following algebra of Pareto vectors: \begin{itemize}
    \item Replace \(A \oplus B\) by the union of vectors from \(A\) and \(B\), followed by retaining only the Pareto-optimal ones;
    \item Replace \(A \otimes B\) by an operation which sums the pairs of vectors from \(A\) and \(B\), followed by retaining the Pareto-optimal ones;
    \item The base cases are replaced with \(\{(0, 0, 0, 0, 0, 0)\}\) and~\(\emptyset\).
\end{itemize}

We start by showing that, when adapted for \cref{prob:synesth:all-vecs}, it is sufficient to try a constant number of subsets \(X, Y, Z \subseteq \geneFamilySet\) at each step of the recurrence of \cref{thm:synesth:minimal-history}.
First, let us show that it is sufficient to place the gain event for each gene family at the lowest common ancestor of the leaves they appear in.

\begin{definition}[Gain positions]
    Let
        \(\syntenyTree = \tuple{T, \syntenyOf, \speciesOf}\) be a synteny tree on \(\geneFamilySet\),
        \(\geneFamily \in \geneFamilySet\),
        \(v \in \treeNodes(T)\),
        and \(X \subseteq \geneFamilySet\).
    We define \begin{align*}
        \lca{}(\geneFamily) &= \lca{T}\{l \in \treeLeaves(T) \mid \geneFamily \in \syntenyOf(l)\}, \\
        f(v, X) &= \{\geneFamily \in X \mid \lca{}(\geneFamily) \not> v\}.
    \end{align*}
\end{definition}

\begin{lemma}[Gains at the LCA]\label{lem:gains-lca-vec}
    If
        \(\syntenyTree = \tuple{T, \syntenyOf, \speciesOf}\) is a synteny tree on \(\geneFamilySet\),
        \(v \in \treeNodes(T)\),
        and \(X \subseteq \geneFamilySet\),
    then
        \(\paretoMin{\preceq}(\Lambda(v, X, \sigma) \cup \Lambda(v, f(v, X), \sigma)) = \Lambda(v, f(v, X), \sigma)\).
\end{lemma}
\begin{proof}
    Let \(\Hist \in h(v, X, \sigma)\).
    Let \(\geneFamily \in X - f(v, X)\), if there is any.
    By definition, \(\lca{}(\geneFamily) > v\).
    Assume w.l.o.g. that \(\lca{}(\geneFamily) \geq v\leftChild\).
    Consider the history \(\Hist'\) in which
        \(\geneFamily\) is removed from \(\syntenyOf(v)\)
        and all \(\syntenyOf(w)\) for \(w \in \subtree{T}{v\rightChild}\),
        removing any invalid loss event created in the process,
        and in which the gain event for \(\geneFamily\) is moved to be the parent of \(v\rightChild\)
        (potentially merging it with other gains).
    Then \(\eventVec(\Hist') \preceq \eventVec(\Hist)\) since we only potentially removed losses from \(\Hist\) and the number of gains is not part of the event vector.
    Let \(\Hist^*\) be the history obtained after repeating this process for each \(\geneFamily \in X - f(v, X)\).
    Clearly \(\Hist^* \in h(v, f(v, X), \sigma)\), hence \(\eventOf(\Hist^*) \in \Lambda(v, f(v, X), \sigma)\) and \(\eventVec(\Hist^*) \preceq \eventVec(\Hist)\).
\end{proof}

In a similar way to lemma~5 from \textcite{Delabre2020EvolutionSegmentalDuplications}, we now show that only two synteny contents have to be tried at each step of the recurrence as any synteny larger than the minimal required gene families (as formally defined below) leads to the same set of optimal event vectors.

\begin{definition}[Minimal synteny contents]
    For any \(v \in \treeNodes(T)\), we define \begin{align*}
        \gainAt(v) &= \{\geneFamily \in \geneFamilySet \mid v = \lca{}(\geneFamily)\}, \\
        \syntenyMinAt(v) &= \begin{cases}
            \syntenyOf(v) & \text{if } v \text{ is a leaf,} \\
            \syntenyMinAt(v\leftChild) \cup \syntenyMinAt(v\rightChild) - (\gainAt(v\leftChild) \cup \gainAt(v\rightChild)) & \text{otherwise.} 
        \end{cases}
        \qedhere
    \end{align*}
\end{definition}

\begin{lemma}[Two choices of synteny contents]\label{lem:synesth:2choices-contents}
    Let \(\syntenyTree = \tuple{T, \syntenyOf, \speciesOf}\) be a synteny tree on a species tree \(\speciesTree\) and \(v\) be a node of \(T\).
    For any \(X, X' \supsetneq \syntenyMinAt(v)\) such that \(X = f(v, X)\) and \(X' = f(v, X')\),
    \(\Lambda(v, X, \sigma) = \Lambda(v, X', \sigma)\).
\end{lemma}
\begin{proof}
    We proceed by induction on the depth of \(v\) in \(T\).
    If \(v\) is a leaf, then \(\Lambda(v, X, \sigma) = \Lambda(v, X', \sigma) = \emptyset\).
    Otherwise, let \(\Hist = \tuple{H, \eventOf, \syntenyOf, \speciesOf} \in h(v, X, \sigma)\), denoting \(\Hist\leftChild \in h(v\leftChild, Y, \speciesOf(v\leftChild))\) (resp. \(\Hist\rightChild\) for \(v\rightChild\)) to be the subhistory of \(\Hist\) below \(v\leftChild\) (resp. \(v\rightChild\)).

    Let \(Y = \syntenyOf(v\leftChild)\), \(Z = \syntenyOf(v\rightChild)\),
        \(G = X - \syntenyMinAt(v)\), and \(G' = X' - \syntenyMinAt(v)\).
    Let \(\geneFamily \in G\),
        then \(\geneFamily \notin \syntenyMinAt(v)\),
        implying that \(\geneFamily \notin \syntenyMinAt(v\leftChild)\) and
            \(\geneFamily \notin \syntenyMinAt(v\rightChild)\),
        since \(\geneFamily \notin \gainAt(v\leftChild)\) and
            \(\geneFamily \notin \gainAt(v\rightChild)\)
        because \(\lca{}(\geneFamily) \not> v\).
    This implies that \(X \not\subseteq Y, Z\).
    Using the same argument for \(\geneFamily' \in G'\), we deduce that \(X' \not\subseteq Y, Z\).

    Suppose that \(Y \supsetneq \syntenyMinAt(v\leftChild)\).
    Let \(Y' = (Y - G) \cup G'\).
    Noting that \(Y' \supsetneq \syntenyMinAt(v\leftChild)\) and using the induction hypothesis, we see that \(\Lambda(v\leftChild, Y, \speciesOf(v\leftChild)) = \Lambda(v\leftChild, Y', \speciesOf(v\leftChild))\).
    There exists a history \(\Hist'\leftChild \in h(v\leftChild, Y', \speciesOf(v\leftChild))\) such that \(\eventVec(\Hist\leftChild) = \eventVec(\Hist\leftChild')\).
    The same argument applies to \(v\rightChild\), yielding a history \(\Hist'\rightChild\) such that \(\eventVec(\Hist\rightChild) = \eventVec(\Hist\rightChild')\).

    Suppose that \(Y = \syntenyMinAt(v\leftChild)\).
    Then there must be an event to lose \(X - Y\) on the path from \(v\) to \(v\leftChild\).
    That event can also be used to lose \(X' - Y\).
    In this case, we let \(\Hist\leftChild' = \Hist\leftChild\).
    The same argument applies to \(v\rightChild\).

    Finally, consider the history \(\Hist'\) obtained by replacing \(\Hist\leftChild\) by \(\Hist'\leftChild\), \(\Hist\rightChild\) by \(\Hist'\rightChild\), and setting \(\syntenyOf(v) = X'\).
    Clearly, \(\eventVec(\Hist) = \eventVec(\Hist')\) and \(\Hist' \in h(v, X', \sigma)\), hence \(\eventVec(\Hist) \in \Lambda(v, X', \sigma)\).
\end{proof}

Using the recurrence from~\cref{thm:synesth:minimal-history} adapted to use the algebra for Pareto vectors and simplified as shown in \cref{lem:gains-lca-vec,lem:synesth:2choices-contents}, we can obtain a dynamic programming algorithm for solving \cref{prob:synesth:all-vecs} and \cref*{prob:synesth:all-vecs}'.

\begin{theorem}\label{thm:synesth:all-vecs-complexity}
    \Cref{prob:synesth:all-vecs} can be solved
        in time \(\mathcal{O}((mn)^9\log(mn))\)
        and space \(\mathcal{O}((mn)^5)\),
    and \cref{prob:synesth:all-vecs}'
        in time \(\mathcal{O}((mn)^9\log(mn))\)
        and space \(\mathcal{O}((|\geneFamilySet|mn)^5)\).
\end{theorem}
\begin{proof}
    Let us first show that, letting \(m = |\treeNodes(T)|\) and \(n = |\treeNodes(\speciesTree)|\), \[
        |\minVecSet(\syntenyTree, \speciesTree)| \in \mathcal{O}((mn)^4).
    \]
    From \cref{thm:synesth:minimal-paths}, the number of events in a history \(\Hist \in \optPath((X, \sigma), (Y, \gamma))\) is bounded by~\(n\) (which can be attained for histories \(P_0\), while \(P_1\) and \(P_2\) always have a constant number of events).
    From \cref{thm:synesth:minimal-history}, histories \(\Hist \in \minSolSpace(\syntenyTree, \speciesTree)\) are extensions of the synteny tree \(T\) obtained by inserting such paths.
    Hence the number of events in such a history is bounded by \(mn\).
    As the vectors have 5~components, we obtain the desired result by adapting the argument of lemma~3.1 from \textcite{LibeskindHadas2014ParetoOptimalPhylogenetic}.

    We solve \cref{prob:synesth:all-vecs} using a dynamic programming table for \(\Lambda(v, X, \sigma)\).
    The number of entries in the table is
        \(m\) for \(v\),
        two for \(X\) (as per \cref{lem:synesth:2choices-contents}),
        and \(n\) for \(\sigma\).
    Hence, the space complexity result follows from the bound on the size of \(\minVecSet(\syntenyTree, \speciesTree)\) shown above.

    As for the time complexity, following the recurrence adapted from \cref{thm:synesth:minimal-history}, computing each entry needs to consider four options for \(Y\) and \(Z\), and up to \(n^3\) options for \(\gamma\leftChild\) and \(\gamma\rightChild\) (or \(\gamma\conservChild\) and \(\gamma\targetChild\)) and \(\gamma_i\).
    However, it is possible to reduce these \(n^3\) species options to a constant number of operations at each step by simultaneously computing three separate tables \(\mathit{in}\), \(\mathit{inAlt}\) and \(\mathit{out}\) as defined by~\textcite{Weiner2021ImprovedDuplicationTransfer} and explained in their algorithm~1 and the proof of their theorem~1.
    For each of these options, the \(\oplus\) and \(\otimes\) operators need to be used a constant number of times.

    The~\(\oplus\) and~\(\otimes\) operations on two sets containing \(k\) Pareto-vectors can be implemented in time \(\mathcal{O}(k\log k)\)~\parencite{Saule2015ParetoOptimizationAlgebraic} and \(\mathcal{O}(k^2\log k)\) respectively~\parencite{LibeskindHadas2014ParetoOptimalPhylogenetic}.
    In our case, it follows from the bound on the size of \(\minVecSet(\syntenyTree, \speciesTree)\) that both operators can be implemented in time \(\mathcal{O}((mn)^4\log(mn))\) and \(\mathcal{O}((mn)^8\log(mn))\) respectively, with \(m = |\treeNodes(T)|\) and \(n = |\treeNodes(\speciesTree)|\).

    To solve \cref*{prob:synesth:all-vecs}', we can extend the Pareto algebra used previously with a Pareto product of a history-generating algebra derived from \cref{thm:synesth:minimal-history}~\parencite{Saule2015ParetoOptimizationAlgebraic}, leading to the same time complexity as for \cref{prob:synesth:all-vecs}.
    As for the space complexity, we need to adjoin a history to each vector, and each history can be represented in space \(\mathcal{O}(|\geneFamilySet|mn)\).
\end{proof}

\section{Efficient computation of minimum-cost histories}\label{sec:synesth:min-cost-histories}

We finally address \cref{prob:synesth:min-cost}.
Let \(\syntenyTree = \tuple{T, \syntenyOf, \speciesOf}\) be a synteny tree on a species tree~\(\speciesTree\) and \(\eventCost \in (\setReal^+ \cup \{\infty\})^5\) be an event cost vector.
Similarly to the way \(h(v, X, \sigma)\) and \(\Lambda(v, X, \sigma)\) were previously used to compute \(\minSolSpace(\syntenyTree, \speciesTree)\) and \(\minVecSet(\syntenyTree, \speciesTree)\), we define \(c(v, X, \sigma) = \{\histCost(\Hist) \mid \Hist \in h(v, X, \sigma)\}\) to compute \(\minCost(\syntenyTree, \speciesTree)\).

We can again adapt the formulation of \cref{thm:synesth:minimal-history} to compute \(\Lambda(v, X, \sigma)\) using the so-called \enquote{min-plus} or \enquote{tropical} algebra: \begin{itemize}
    \item Replace \(A \oplus B\) by taking the minimum of the two values;
    \item Replace \(A \otimes B\) by the sum of the two values;
    \item The base cases are replaced with \(0\) and~\(\infty\).
\end{itemize}
Notice that \cref{lem:gains-lca-vec,lem:synesth:2choices-contents} still apply, since if \(\Lambda(v, X, \sigma)\) = \(\Lambda(v, X', \sigma)\), then \(c(v, X, \sigma) = c(v, X', \sigma)\).

\begin{theorem}
    \Cref{prob:synesth:min-cost} can be solved
        in time \(\mathcal{O}(mn)\)
        and space \(\mathcal{O}(mn)\),
    and \cref*{prob:synesth:min-cost}'
        in time \(\mathcal{O}(mn)\)
        and space \(\mathcal{O}((|\geneFamilySet|mn)^2)\).
\end{theorem}
\begin{proof}
    Both results follow from the proof of \cref{thm:synesth:all-vecs-complexity}, the fact that the minimum and sum operators can be computed in constant time, removing the \((mn)^{8}\log(mn)\) factor from the time complexity, and the fact that only a single value needs to be stored for each entry of the dynamic programming table, removing the \((mn)^4\) factor from the space complexity.
\end{proof}

\section{Results}\label{sec:synesth:results}

The CRISPR–Cas module is an adaptive system used by prokaryotes to defend against invading viruses and plasmids. 
Its fame is due to the development of the CRISPR-Cas~9 genome editing technology, which is one of the most reliable and accurate “molecular scissors”. An important part of a CRISPR-Cas system is an operon of associated Cas genes playing various roles in the defense machinery. As the microbial function of CRISPR-Cas systems highly depends on the syntenic organization of Cas genes, elucidating the evolution of these syntenies is crucial.

\begin{figure}[h]
    \centering
    \scalebox{.27}{\tikzset{
    y={(1pt,0)},
    x={(0,-1pt)},
    limit width/.style 2 args={
        execute at begin node={\begin{varwidth}{#1}#2},
        execute at end node={\end{varwidth}},
    },
    host background/.style={
        fill=host background color,
        draw=host background color,
        line width={.01pt},
    },
    unsampled host background/.style={
        preaction={fill=host background color},
        draw=host background color,
        pattern={Lines[angle=45, distance=3pt, line width=1.5pt]},
        pattern color=white,
    },
    host label/.style={
        font=\scriptsize,
        midway,
        anchor=west,
        limit width={9em}{\narrowragged},
    },
    branch/.style={
        draw={#1},
        line width={0.5pt},
    },
    loss branch/.style={
        branch,
        draw={black!40},
    },
    transfer branch/.style={
        branch={#1},
        densely dashed,
        -{Latex[length=0pt 8]},
    },
    invalid transfer branch/.style={
        transfer branch={red!80!black},
        out distance=20,
        in distance=20,
    },
    event/.style={
        draw={#1}, fill={host background color!50!white},
        font={\color{#1}},
        outer sep=0pt, inner xsep=0pt, inner ysep=2pt,
        limit width={6em}{\narrowragged},
        line width={0.5pt},
    },
    event/.default={black},
    extant/.code 2 args={
        \pgfkeysalso{
            circle, fill={#1},
            outer sep=0pt, inner sep=0pt,
            minimum size={3},
        }
        \def\cmp{#2}\ifx\cmp\empty\else
            \pgfkeysalso{
                label={[%
                    font={\color{#1}\vphantom{gb}},
                    % fill=host background color!35!white,
                    rounded corners=2,
                    label distance=2pt,
                    inner xsep=2pt, inner ysep=1pt,
                    outer xsep=0pt, outer ysep=0pt,
                    limit width={6em}{\narrowragged},
                ]right:#2},
            }
        \fi
    },
    extant/.default={black}{},
    unsampled/.style={
        circle, fill={#1},
        outer sep=0pt, inner sep=0pt,
        minimum size={3},
    },
    unsampled/.default={black!40},
    speciation/.style 2 args={
        event={#1}, rectangle, rounded corners,
        inner xsep=4pt,
        minimum width={8},
        minimum height={8},
        execute at begin node={#2},
    },
    duplication/.style 2 args={
        event={#1}, rectangle,
        inner xsep=4pt,
        minimum width={8},
        minimum height={8},
        execute at begin node={#2},
    },
    transfer/.code 2 args={
        \pgfkeysalso{
            event={#1},
        }
        \def\cmp{#2}\ifx\cmp\empty
            \pgfkeysalso{
                diamond,
                inner sep=1.5pt,
            }
        \else
            \pgfkeysalso{
                chamfered rectangle,
                chamfered rectangle sep={8 / 2.4},
                inner xsep=2pt,
                inner ysep=-1pt,
                minimum width={8},
                minimum height={8},
                execute at begin node={#2},
            }
        \fi
    },
    gain/.code 2 args={
        \pgfkeysalso{
            draw={#1}, fill={#1},
            semicircle, inner sep=1pt,
            shape border rotate={90},
        }
        \def\cmp{#2}\ifx\cmp\empty\else
            \pgfkeysalso{
                label={[%
                    font={\scriptsize\color{#1}\vphantom{gb}},
                    fill=host background color,
                    label distance=2pt,
                    inner xsep=0pt, inner ysep=0pt,
                    outer xsep=0pt, outer ysep=0pt,
                    limit width={6em}{\centering},
                ]above:#2},
            }
        \fi
    },
    loss/.code 2 args={
        \pgfkeysalso{
            draw={#1}, fill={#1},
            semicircle, inner sep=1pt,
            shape border rotate={180+90},
        }
        \def\cmp{#2}\ifx\cmp\empty
            \pgfkeysalso{draw={black!40}, fill={black!40}}
        \else
            \pgfkeysalso{
                label={[%
                    font={\scriptsize\color{#1}\vphantom{gb}},
                    fill=host background color,
                    label distance=2pt,
                    inner xsep=0pt, inner ysep=0pt,
                    outer xsep=0pt, outer ysep=0pt,
                    limit width={6em}{\centering},
                ]above:#2},
            }
        \fi
    },
}
\begin{tikzpicture}
% hosts
\path[host background] (0.0,-68.44) [rounded corners={4pt}] to (0.0,1.0) to node[host label] {Clostridium kluyveri DSM 555} (28.87,1.0) [sharp corners] to (28.87,-68.44) -- cycle;
\path[host background] (32.87,-67.61) [rounded corners={4pt}] to (32.87,1.0) to node[host label] {Bacillus halodurans C-125} (61.7399,1.0) [sharp corners] to (61.7399,-67.61) -- cycle;
\path[host background] (65.7399,-113.63) [rounded corners={4pt}] to (65.7399,1.0) to node[host label] {Staphylococcus epidermidis} (94.3299,1.0) [sharp corners] to (94.3299,-113.63) -- cycle;
\path[host background] (32.87,-33.805) [rounded corners={4pt}] to (32.87,-156.345) to (56.4199,-156.345) [sharp corners] to (56.4199,-160.345) to (70.9199,-160.345) [rounded corners={4pt}] to (70.9199,-156.345) to (94.3299,-156.345) [sharp corners] to (94.3299,-56.815) to (65.7399,-56.815) [rounded corners={4pt}] to (65.7399,-141.845) to (61.7399,-141.845) [sharp corners] to (61.7399,-33.805) -- cycle;
\path[host background] (0.0,-34.22) [rounded corners={4pt}] to (0.0,-174.845) to (31.8025,-174.845) [sharp corners] to (31.8025,-178.845) to (46.3025,-178.845) [rounded corners={4pt}] to (46.3025,-174.845) to (70.9199,-174.845) [sharp corners] to (70.9199,-151.095) to (56.4199,-151.095) [rounded corners={4pt}] to (56.4199,-160.345) to (28.87,-160.345) [sharp corners] to (28.87,-34.22) -- cycle;
\path[host background] (98.3299,-87.235) [rounded corners={4pt}] to (98.3299,1.0) to node[host label] {Rhodococcus jostii RHA1} (120.344,1.0) [sharp corners] to (120.344,-87.235) -- cycle;
\path[host background] (31.8025,-169.595) [rounded corners={4pt}] to (31.8025,-193.845) to (68.1983,-193.845) [sharp corners] to (68.1983,-197.845) to (82.6983,-197.845) [rounded corners={4pt}] to (82.6983,-193.845) to (120.344,-193.845) [sharp corners] to (120.344,-43.6175) to (98.3299,-43.6175) [rounded corners={4pt}] to (98.3299,-179.345) to (46.3025,-179.345) [sharp corners] to (46.3025,-169.595) -- cycle;
\path[host background] (124.3441,-76.252) [rounded corners={4pt}] to (124.3441,1.0) to node[host label] {Thermoflexia bacterium} (146.3582,1.0) [sharp corners] to (146.3582,-76.252) -- cycle;
\path[host background] (68.1983,-188.595) [rounded corners={4pt}] to (68.1983,-243.733) to (99.4032,-243.733) [sharp corners] to (99.4032,-247.733) to (113.9032,-247.733) [rounded corners={4pt}] to (113.9032,-243.733) to (146.3582,-243.733) [sharp corners] to (146.3582,-38.126) to (124.3441,-38.126) [rounded corners={4pt}] to (124.3441,-229.233) to (82.6983,-229.233) [sharp corners] to (82.6983,-188.595) -- cycle;
\path[host background] (150.3582,-229.233) [rounded corners={4pt}] to (150.3582,1.0) to node[host label] {Cyanothece sp.\ PCC 8802} (179.2282,1.0) [sharp corners] to (179.2282,-229.233) -- cycle;
\path[host background] (183.2282,-130.13) [rounded corners={4pt}] to (183.2282,1.0) to node[host label] {Synechocystis sp.\ 6803} (205.2423,1.0) [sharp corners] to (205.2423,-130.13) -- cycle;
\path[host background] (150.3582,-114.6165) [rounded corners={4pt}] to (150.3582,-243.733) to (170.3857,-243.733) [sharp corners] to (170.3857,-256.655) to (191.1499,-256.655) [rounded corners={4pt}] to (191.1499,-243.733) to (205.2423,-243.733) [sharp corners] to (205.2423,-65.065) to (183.2282,-65.065) [rounded corners={4pt}] to (183.2282,-229.233) to (179.2282,-229.233) [sharp corners] to (179.2282,-114.6165) -- cycle;
\path[host background] (99.4032,-238.483) [rounded corners={4pt}] to (99.4032,-271.155) to (130.996,-271.155) [sharp corners] to (130.996,-304.216) to (159.5572,-304.216) [rounded corners={4pt}] to (159.5572,-271.155) to (191.1499,-271.155) [sharp corners] to (191.1499,-242.944) to (170.3857,-242.944) [rounded corners={4pt}] to (170.3857,-256.655) to (113.9032,-256.655) [sharp corners] to (113.9032,-238.483) -- cycle;
\path[host background] (209.2423,-125.63) [rounded corners={4pt}] to (209.2423,1.0) to node[host label] {Thermotoga lettingae TMO} (226.2423,1.0) [sharp corners] to (226.2423,-125.63) -- cycle;
\path[host background] (130.996,-280.4355) [rounded corners={4pt}] to (130.996,-310.216) to (182.0247,-310.216) [sharp corners] to (182.0247,-314.216) to (188.0247,-314.216) [rounded corners={4pt}] to (188.0247,-310.216) to (226.2423,-310.216) [sharp corners] to (226.2423,-62.815) to (209.2423,-62.815) [rounded corners={4pt}] to (209.2423,-304.216) to (159.5572,-304.216) [sharp corners] to (159.5572,-280.4355) -- cycle;
\path[host background] (230.2423,-101.63) [rounded corners={4pt}] to (230.2423,1.0) to node[host label] {Candidatus Scalindua brodae} (247.2423,1.0) [sharp corners] to (247.2423,-101.63) -- cycle;
\path[host background] (251.2423,-179.345) [rounded corners={4pt}] to (251.2423,1.0) to node[host label] {Geobacter sulfurreducens} (280.1122,1.0) [sharp corners] to (280.1122,-179.345) -- cycle;
\path[host background] (284.1122,-82.735) [rounded corners={4pt}] to (284.1122,1.0) to node[host label] {Thioalkalivibrio sp.\ K90mix} (301.1122,1.0) [sharp corners] to (301.1122,-82.735) -- cycle;
\path[host background] (324.1122,-67.61) [rounded corners={4pt}] to (324.1122,1.0) to node[host label] {Shewanella putrefaciens CN-32} (352.9821,1.0) [sharp corners] to (352.9821,-67.61) -- cycle;
\path[host background] (356.9821,-90.157) [rounded corners={4pt}] to (356.9821,1.0) to node[host label] {Vibrio crassostreae J5 20} (378.9963,1.0) [sharp corners] to (378.9963,-90.157) -- cycle;
\path[host background] (382.9963,-67.61) [rounded corners={4pt}] to (382.9963,1.0) to node[host label] {Yersinia pseudotuberculosis} (411.8663,1.0) [sharp corners] to (411.8663,-67.61) -- cycle;
\path[host background] (415.8663,-73.44) [rounded corners={4pt}] to (415.8663,1.0) to node[host label] {Escherichia coli K12} (444.7362,1.0) [sharp corners] to (444.7362,-73.44) -- cycle;
\path[host background] (382.9963,-33.805) [rounded corners={4pt}] to (382.9963,-87.94) to (406.6163,-87.94) [sharp corners] to (406.6163,-91.94) to (421.1163,-91.94) [rounded corners={4pt}] to (421.1163,-87.94) to (444.7362,-87.94) [sharp corners] to (444.7362,-36.72) to (415.8663,-36.72) [rounded corners={4pt}] to (415.8663,-73.44) to (411.8663,-73.44) [sharp corners] to (411.8663,-33.805) -- cycle;
\path[host background] (356.9821,-45.0785) [rounded corners={4pt}] to (356.9821,-106.44) to (384.9313,-106.44) [sharp corners] to (384.9313,-110.44) to (399.4313,-110.44) [rounded corners={4pt}] to (399.4313,-106.44) to (421.1163,-106.44) [sharp corners] to (421.1163,-82.69) to (406.6163,-82.69) [rounded corners={4pt}] to (406.6163,-91.94) to (378.9963,-91.94) [sharp corners] to (378.9963,-45.0785) -- cycle;
\path[host background] (324.1122,-33.805) [rounded corners={4pt}] to (324.1122,-124.94) to (354.9821,-124.94) [sharp corners] to (354.9821,-141.845) to (375.7463,-141.845) [rounded corners={4pt}] to (375.7463,-124.94) to (399.4313,-124.94) [sharp corners] to (399.4313,-101.19) to (384.9313,-101.19) [rounded corners={4pt}] to (384.9313,-110.44) to (352.9821,-110.44) [sharp corners] to (352.9821,-33.805) -- cycle;
\path[host background] (305.1122,-36.97) [rounded corners={4pt}] to (305.1122,-156.345) to (333.3043,-156.345) [sharp corners] to (333.3043,-160.345) to (347.8043,-160.345) [rounded corners={4pt}] to (347.8043,-156.345) to (375.7463,-156.345) [sharp corners] to (375.7463,-126.1425) to (354.9821,-126.1425) [rounded corners={4pt}] to (354.9821,-141.845) to (320.1122,-141.845) [sharp corners] to (320.1122,-36.97) -- cycle;
\path[host background] (284.1122,-41.3675) [rounded corners={4pt}] to (284.1122,-174.845) to (309.3332,-174.845) [sharp corners] to (309.3332,-178.845) to (323.8332,-178.845) [rounded corners={4pt}] to (323.8332,-174.845) to (347.8043,-174.845) [sharp corners] to (347.8043,-151.095) to (333.3043,-151.095) [rounded corners={4pt}] to (333.3043,-160.345) to (301.1122,-160.345) [sharp corners] to (301.1122,-41.3675) -- cycle;
\path[host background] (251.2423,-89.6725) [rounded corners={4pt}] to (251.2423,-193.845) to (283.8802,-193.845) [sharp corners] to (283.8802,-197.845) to (298.3802,-197.845) [rounded corners={4pt}] to (298.3802,-193.845) to (323.8332,-193.845) [sharp corners] to (323.8332,-169.595) to (309.3332,-169.595) [rounded corners={4pt}] to (309.3332,-179.345) to (280.1122,-179.345) [sharp corners] to (280.1122,-89.6725) -- cycle;
\path[host background] (230.2423,-50.815) [rounded corners={4pt}] to (230.2423,-310.216) to (261.9363,-310.216) [sharp corners] to (261.9363,-314.216) to (267.9363,-314.216) [rounded corners={4pt}] to (267.9363,-310.216) to (298.3802,-310.216) [sharp corners] to (298.3802,-188.595) to (283.8802,-188.595) [rounded corners={4pt}] to (283.8802,-304.216) to (247.2423,-304.216) [sharp corners] to (247.2423,-50.815) -- cycle;
\path[host background] (182.0247,-309.216) [rounded corners={4pt}] to (182.0247,-320.216) to (221.9805,-320.216) [sharp corners] to (221.9805,-324.216) to (227.9805,-324.216) [rounded corners={4pt}] to (227.9805,-320.216) to (267.9363,-320.216) [sharp corners] to (267.9363,-309.216) to (261.9363,-309.216) [rounded corners={4pt}] to (261.9363,-314.216) to (188.0247,-314.216) [sharp corners] to (188.0247,-309.216) -- cycle;
% unsampled hosts
\path[unsampled host background] (305.1122,-133.94) [rounded corners={2pt}] to (305.1122,1.0) to  (320.1122,1.0) [sharp corners] to (320.1122,-133.94) -- cycle;
% branches
\path[branch] (47.3049,-65.6014) to (47.3049,-59.11);
\path[branch] (80.0349,-113.63) to (80.0349,-51.61);
\path[branch] (63.6699,-149.095) [rounded corners={1.5pt}] to (47.3049,-149.095) [sharp corners] to (47.3049,-65.6014);
\path[loss branch] (63.6699,-149.095) to (78.3049,-149.095) node[sideways loss] {};
\path[branch] (39.0525,-167.595) [rounded corners={1.5pt}] to (14.435,-167.595) [sharp corners] to (14.435,-59.94);
\path[branch] (39.0525,-167.595) [rounded corners={1.5pt}] to (63.6699,-167.595) [sharp corners] to (63.6699,-149.095);
\path[branch] (39.0525,-174.845) to (39.0525,-167.595);
\path[branch] (111.8441,-66.7824) to (111.8441,-53.83);
\path[branch] (111.8441,-82.735) to (111.8441,-66.7824);
\path[transfer branch] (111.8441,-82.735) to (292.6122,-82.735);
\path[branch] (75.4483,-186.595) [rounded corners={1.5pt}] to (111.8441,-186.595) [sharp corners] to (111.8441,-82.735);
\path[loss branch] (75.4483,-186.595) to (61.8441,-186.595) node[base loss={black!40}{}{0}{right}] {};;
\path[branch] (137.8582,-69.7909) to (137.8582,-58.83);
\path[branch] (106.6532,-236.483) [rounded corners={1.5pt}] to (75.4483,-236.483) [sharp corners] to (75.4483,-186.595);
\path[branch] (106.6532,-236.483) [rounded corners={1.5pt}] to (137.8582,-236.483) [sharp corners] to (137.8582,-69.7909);
\path[branch] (164.7932,-73.6949) to (164.7932,-54.39);
\path[branch] (164.7932,-193.845) to (164.7932,-73.6949);
\path[transfer branch] (164.7932,-193.845) to (291.1302,-193.845);
\path[branch] (164.7932,-215.2891) to (164.7932,-193.845);
\path[branch] (196.7424,-78.7299) to (196.7424,-58.83);
\path[branch] (196.7424,-101.63) to (196.7424,-78.7299);
\path[transfer branch] (196.7424,-101.63) to (238.7423,-101.63);
\path[branch] (196.7424,-113.63) to (196.7424,-101.63);
\path[transfer branch] (196.7424,-113.63) to (80.0349,-113.63);
\path[branch] (196.7424,-125.63) to (196.7424,-113.63);
\path[transfer branch] (196.7424,-125.63) to (217.7423,-125.63);
\path[branch] (183.8999,-236.483) [rounded corners={1.5pt}] to (196.7424,-236.483) [sharp corners] to (196.7424,-125.63);
\path[branch] (183.8999,-236.483) [rounded corners={1.5pt}] to (164.7932,-236.483) [sharp corners] to (164.7932,-215.2891);
\path[branch] (183.8999,-250.1941) to (183.8999,-236.483);
\path[branch] (152.3072,-263.905) [rounded corners={1.5pt}] to (106.6532,-263.905) [sharp corners] to (106.6532,-236.483);
\path[branch] (152.3072,-263.905) [rounded corners={1.5pt}] to (183.8999,-263.905) [sharp corners] to (183.8999,-250.1941);
\path[branch] (152.3072,-287.6856) to (152.3072,-263.905);
\path[branch] (217.7423,-125.63) to (217.7423,-58.83);
\path[branch] (238.7423,-101.63) to (238.7423,-53.83);
\path[branch] (265.6772,-68.4229) to (265.6772,-59.94);
\path[branch] (265.6772,-174.845) to (265.6772,-68.4229);
\path[transfer branch] (265.6772,-174.845) to (39.0525,-174.845);
\path[branch] (292.6122,-82.735) to (292.6122,-48.83);
\path[branch] (312.6122,-69.44) to (312.6122,-4.5);
\path[transfer branch] (312.6122,-69.44) to (430.3013,-69.44);
\path[branch] (338.5472,-65.6014) to (338.5472,-59.11);
\path[branch] (370.4963,-71.7434) to (370.4963,-48.83);
\path[branch] (430.3013,-69.44) to (430.3013,-64.94);
\path[branch] (413.8663,-80.69) [rounded corners={1.5pt}] to (397.4313,-80.69) [sharp corners] to (397.4313,-59.11);
\path[loss branch] (413.8663,-80.69) to (428.4313,-80.69) node[sideways loss] {};
\path[branch] (392.1813,-99.19) [rounded corners={1.5pt}] to (370.4963,-99.19) [sharp corners] to (370.4963,-71.7434);
\path[branch] (392.1813,-99.19) [rounded corners={1.5pt}] to (413.8663,-99.19) [sharp corners] to (413.8663,-80.69);
\path[branch] (368.4963,-117.69) [rounded corners={1.5pt}] to (338.5472,-117.69) [sharp corners] to (338.5472,-65.6014);
\path[branch] (368.4963,-117.69) [rounded corners={1.5pt}] to (392.1813,-117.69) [sharp corners] to (392.1813,-99.19);
\path[branch] (368.4963,-133.3924) to (368.4963,-117.69);
\path[branch] (340.5543,-149.095) [rounded corners={1.5pt}] to (368.4963,-149.095) [sharp corners] to (368.4963,-133.3924);
\path[branch] (340.5543,-149.095) [rounded corners={1.5pt}] to (312.6122,-149.095) [sharp corners] to (312.6122,-69.44);
\path[branch] (316.5832,-167.595) [rounded corners={1.5pt}] to (340.5543,-167.595) [sharp corners] to (340.5543,-149.095);
\path[loss branch] (316.5832,-167.595) to (301.5543,-167.595) node[base loss={black!40}{}{0}{right}] {};;
\path[branch] (291.1302,-186.595) [rounded corners={1.5pt}] to (265.6772,-186.595) [sharp corners] to (265.6772,-174.845);
\path[branch] (291.1302,-186.595) [rounded corners={1.5pt}] to (316.5832,-186.595) [sharp corners] to (316.5832,-167.595);
\path[branch] (291.1302,-193.845) to (291.1302,-186.595);
% gene transfers
% events
\node[extant={black}{1, 2, 3", 3',\\4, 5, 6, 7, 8}] at (14.435,-59.94) {};
\node[extant={black}{1, 2, 3", 3', 4, 5, 7, 8}] at (47.3049,-59.11) {};
\node[loss={black}{6}] at (47.3049,-65.6014) {};
\node[extant={black}{1, 2, 5, 6, 7, 10, 11}] at (80.0349,-51.61) {};
\node[speciation={black}{}] at (63.6699,-149.095) {};
\node[speciation={black}{}] at (39.0525,-167.595) {};
\node[extant={black}{5, 7, 8, 11}] at (111.8441,-53.83) {};
\node[loss={black}{6, 10}] at (111.8441,-66.7824) {};
\node[transfer={black}{}] at (111.8441,-82.735) {};
\node[speciation={black}{}] at (75.4483,-186.595) {};
\node[extant={black}{5, 7, 10, 11}] at (137.8582,-58.83) {};
\node[loss={black}{6, 8}] at (137.8582,-69.7909) {};
\node[speciation={black}{}] at (106.6532,-236.483) {};
\node[extant={black}{1, 2, 3', 4,\\5, 6, 7, 10}] at (164.7932,-54.39) {};
\node[loss={black}{3", 8, 11}] at (164.7932,-73.6949) {};
\node[transfer={black}{}] at (164.7932,-193.845) {};
\node[gain={black}{3", 3', 4}] at (164.7932,-215.2891) {};
\node[extant={black}{5, 7, 10, 11}] at (196.7424,-58.83) {};
\node[loss={black}{1, 2, 6, 8}] at (196.7424,-78.7299) {};
\node[transfer={black}{}] at (196.7424,-101.63) {};
\node[transfer={black}{}] at (196.7424,-113.63) {};
\node[transfer={black}{}] at (196.7424,-125.63) {};
\node[speciation={black}{}] at (183.8999,-236.483) {};
\node[gain={black}{1, 2}] at (183.8999,-250.1941) {};
\node[speciation={black}{}] at (152.3072,-263.905) {};
\node[gain={black}{5, 6, 7,\\8, 10, 11}] at (152.3072,-287.6856) {};
\node[extant={black}{5, 7, 10, 11}] at (217.7423,-58.83) {};
\node[extant={black}{1, 2, 7, 11}] at (238.7423,-53.83) {};
\node[extant={black}{1, 2, 3", 3',\\4, 5, 6, 7, 8}] at (265.6772,-59.94) {};
\node[loss={black}{11}] at (265.6772,-68.4229) {};
\node[transfer={black}{}] at (265.6772,-174.845) {};
\node[extant={black}{5, 6, 7, 8}] at (292.6122,-48.83) {};
\node[unsampled] at (312.6122,-4.5) {};
\node[transfer={black}{}] at (312.6122,-69.44) {};
\node[extant={black}{1, 2, 3", 3', 5, 6, 7}] at (338.5472,-59.11) {};
\node[loss={black}{8}] at (338.5472,-65.6014) {};
\node[extant={black}{5, 6, 7, 8}] at (370.4963,-48.83) {};
\node[loss={black}{1, 2, 3", 3'}] at (370.4963,-71.7434) {};
\node[extant={black}{1, 2, 3", 3', 5, 6, 7, 8}] at (397.4313,-59.11) {};
\node[extant={black}{1, 2, 3", 3', 5, 6, 7, 8, 11}] at (430.3013,-64.94) {};
\node[speciation={black}{}] at (413.8663,-80.69) {};
\node[speciation={black}{}] at (392.1813,-99.19) {};
\node[speciation={black}{}] at (368.4963,-117.69) {};
\node[loss={black}{4, 11}] at (368.4963,-133.3924) {};
\node[speciation={black}{}] at (340.5543,-149.095) {};
\node[speciation={black}{}] at (316.5832,-167.595) {};
\node[speciation={black}{}] at (291.1302,-186.595) {};
\end{tikzpicture}
}
    \caption{Output of Synesth for the CRISPR-Cas Class 1 dataset, when asked for one minimum-cost history for event costs \(\eventCost(\eventLoss)=1\), \(\eventCost(\eventDup)=2\), \(\eventCost(\eventCut) = 2.5\), \(\eventCost(\eventTrDup)=4\) and \(\eventCost(\eventTrCut) = 4.5\).}
    \label{fig:synesth:crispr-sol}
\end{figure}

In \cref{chap:superdtl}, taking the CRISPR-Cas classification of Class~1 from \textcite{Makarova2020EvolutionaryClassificationCRISPR} as the synteny tree, and considering a dataset of 15 bacterial species as well as the species tree topology inferred in \textcite{Coleman2021RootedPhylogenyResolves}, we recovered an evolutionary scenario which is broadly consistent with that of \textcite{Koonin2022EvolutionaryPlasticityFunctional} with the CRISPR-Cas emergence inferred at the root of Terrabacteria. However, some inconsistencies were observed. For example, in the Proteobacteria subtree, the \emph{SuperDTL} algorithm inferred an unlikely scenario with an ancestral synteny duplication before the LCA of \emph{Shewanella putrefaciens}, \emph{Vibrio crassostreae}, \emph{Yersinia pseudotuberculosis}, and  \emph{Escherichia coli} resulting in a succession of three consecutive full synteny losses along the branch to \emph{E. coli}.

We used Synesth with the same event costs for Loss, Dup and TrDup, and setting costs for \(\eventCut\) and \(\eventTrCut\)
such that \((\eventCost_\eventDup, \eventCost_\eventCut, \eventCost_\eventTrDup, \eventCost_\eventTrCut, \eventCost_\eventLoss) = (2, 2.5, 4, 4.5, 1)\).
Almost the same history was obtained, but with the above unlikely scenario replaced by a speciation (on the branch separating the ancestor of \emph{Geobacter sulfurreducens} to the LCA of \emph{Thioalkalivibrio} and \emph{Shewanella putrefaciens}) copying the ancestral synteny to an unsampled species, which is later transferred back to \emph{E. coli} (see Figure~\ref{fig:synesth:crispr-sol}).

Choosing the appropriate event cost is one of the main challenge of tree reconciliation. 
To visualize the set of Pareto-optimal vectors computed as explained before, we use a similar approach as developed by \textcite{LibeskindHadas2014ParetoOptimalPhylogenetic}, to display a summary of the solution space with loss events normalized to cost 1, and the other cost events varying in a given range. In order to represent the solutions in a 2D plot, we set an equal cost for \(\eventDup\) and \(\eventCut\) and for \(\eventTrDup\) and \(\eventTrCut\), and we represent a solution by the condensed event cost vector \(c_\eventDup + c_\eventCut, c_\eventTrDup + c_\eventTrCut, c_\eventLoss)\).
The resulting plot is given in \cref{fig:synesth:crispr-pareto}, in which each color-coded region corresponds to a set of event costs that give rise to exactly the same set of minimal histories.

\begin{figure}[h]
    \centering
    \begingroup%
\makeatletter%
\begin{pgfpicture}%
\pgfpathrectangle{\pgfpointorigin}{\pgfqpoint{4.743371in}{3.994767in}}%
\pgfusepath{use as bounding box, clip}%
\begin{pgfscope}%
\pgfsetbuttcap%
\pgfsetmiterjoin%
\definecolor{currentfill}{rgb}{1.000000,1.000000,1.000000}%
\pgfsetfillcolor{currentfill}%
\pgfsetlinewidth{0.000000pt}%
\definecolor{currentstroke}{rgb}{1.000000,1.000000,1.000000}%
\pgfsetstrokecolor{currentstroke}%
\pgfsetdash{}{0pt}%
\pgfpathmoveto{\pgfqpoint{0.000000in}{0.000000in}}%
\pgfpathlineto{\pgfqpoint{4.743371in}{0.000000in}}%
\pgfpathlineto{\pgfqpoint{4.743371in}{3.994767in}}%
\pgfpathlineto{\pgfqpoint{0.000000in}{3.994767in}}%
\pgfpathlineto{\pgfqpoint{0.000000in}{0.000000in}}%
\pgfpathclose%
\pgfusepath{fill}%
\end{pgfscope}%
\begin{pgfscope}%
\pgfsetbuttcap%
\pgfsetmiterjoin%
\definecolor{currentfill}{rgb}{1.000000,1.000000,1.000000}%
\pgfsetfillcolor{currentfill}%
\pgfsetlinewidth{0.000000pt}%
\definecolor{currentstroke}{rgb}{0.000000,0.000000,0.000000}%
\pgfsetstrokecolor{currentstroke}%
\pgfsetstrokeopacity{0.000000}%
\pgfsetdash{}{0pt}%
\pgfpathmoveto{\pgfqpoint{0.636060in}{1.830025in}}%
\pgfpathlineto{\pgfqpoint{4.107310in}{1.830025in}}%
\pgfpathlineto{\pgfqpoint{4.107310in}{3.846542in}}%
\pgfpathlineto{\pgfqpoint{0.636060in}{3.846542in}}%
\pgfpathlineto{\pgfqpoint{0.636060in}{1.830025in}}%
\pgfpathclose%
\pgfusepath{fill}%
\end{pgfscope}%
\begin{pgfscope}%
\pgfpathrectangle{\pgfqpoint{0.636060in}{1.830025in}}{\pgfqpoint{3.471250in}{2.016517in}}%
\pgfusepath{clip}%
\pgfsetbuttcap%
\pgfsetmiterjoin%
\definecolor{currentfill}{rgb}{0.904762,0.365079,0.301587}%
\pgfsetfillcolor{currentfill}%
\pgfsetlinewidth{1.003750pt}%
\definecolor{currentstroke}{rgb}{0.904762,0.365079,0.301587}%
\pgfsetstrokecolor{currentstroke}%
\pgfsetdash{}{0pt}%
\pgfpathmoveto{\pgfqpoint{2.718810in}{2.838284in}}%
\pgfpathlineto{\pgfqpoint{4.107310in}{2.838284in}}%
\pgfpathlineto{\pgfqpoint{4.107310in}{3.241587in}}%
\pgfpathlineto{\pgfqpoint{2.718810in}{2.838284in}}%
\pgfpathclose%
\pgfusepath{stroke,fill}%
\end{pgfscope}%
\begin{pgfscope}%
\pgfpathrectangle{\pgfqpoint{0.636060in}{1.830025in}}{\pgfqpoint{3.471250in}{2.016517in}}%
\pgfusepath{clip}%
\pgfsetbuttcap%
\pgfsetmiterjoin%
\definecolor{currentfill}{rgb}{1.000000,0.603175,0.444445}%
\pgfsetfillcolor{currentfill}%
\pgfsetlinewidth{1.003750pt}%
\definecolor{currentstroke}{rgb}{1.000000,0.603175,0.444445}%
\pgfsetstrokecolor{currentstroke}%
\pgfsetdash{}{0pt}%
\pgfpathmoveto{\pgfqpoint{1.330310in}{2.636632in}}%
\pgfpathlineto{\pgfqpoint{4.107310in}{2.636632in}}%
\pgfpathlineto{\pgfqpoint{4.107310in}{2.838284in}}%
\pgfpathlineto{\pgfqpoint{2.718810in}{2.838284in}}%
\pgfpathlineto{\pgfqpoint{1.330310in}{2.636632in}}%
\pgfpathclose%
\pgfusepath{stroke,fill}%
\end{pgfscope}%
\begin{pgfscope}%
\pgfpathrectangle{\pgfqpoint{0.636060in}{1.830025in}}{\pgfqpoint{3.471250in}{2.016517in}}%
\pgfusepath{clip}%
\pgfsetbuttcap%
\pgfsetmiterjoin%
\definecolor{currentfill}{rgb}{0.761905,0.507936,0.206349}%
\pgfsetfillcolor{currentfill}%
\pgfsetlinewidth{1.003750pt}%
\definecolor{currentstroke}{rgb}{0.761905,0.507936,0.206349}%
\pgfsetstrokecolor{currentstroke}%
\pgfsetdash{}{0pt}%
\pgfpathmoveto{\pgfqpoint{0.636060in}{2.233329in}}%
\pgfpathlineto{\pgfqpoint{4.107310in}{2.233329in}}%
\pgfpathlineto{\pgfqpoint{4.107310in}{2.636632in}}%
\pgfpathlineto{\pgfqpoint{1.330310in}{2.636632in}}%
\pgfpathlineto{\pgfqpoint{0.636060in}{2.434980in}}%
\pgfpathlineto{\pgfqpoint{0.636060in}{2.233329in}}%
\pgfpathclose%
\pgfusepath{stroke,fill}%
\end{pgfscope}%
\begin{pgfscope}%
\pgfpathrectangle{\pgfqpoint{0.636060in}{1.830025in}}{\pgfqpoint{3.471250in}{2.016517in}}%
\pgfusepath{clip}%
\pgfsetbuttcap%
\pgfsetmiterjoin%
\definecolor{currentfill}{rgb}{0.936508,0.698413,0.063492}%
\pgfsetfillcolor{currentfill}%
\pgfsetlinewidth{1.003750pt}%
\definecolor{currentstroke}{rgb}{0.936508,0.698413,0.063492}%
\pgfsetstrokecolor{currentstroke}%
\pgfsetdash{}{0pt}%
\pgfpathmoveto{\pgfqpoint{0.636060in}{2.166112in}}%
\pgfpathlineto{\pgfqpoint{4.107310in}{2.166112in}}%
\pgfpathlineto{\pgfqpoint{4.107310in}{2.233329in}}%
\pgfpathlineto{\pgfqpoint{0.636060in}{2.233329in}}%
\pgfpathlineto{\pgfqpoint{0.636060in}{2.166112in}}%
\pgfpathclose%
\pgfusepath{stroke,fill}%
\end{pgfscope}%
\begin{pgfscope}%
\pgfpathrectangle{\pgfqpoint{0.636060in}{1.830025in}}{\pgfqpoint{3.471250in}{2.016517in}}%
\pgfusepath{clip}%
\pgfsetbuttcap%
\pgfsetmiterjoin%
\definecolor{currentfill}{rgb}{0.666667,0.793651,0.412698}%
\pgfsetfillcolor{currentfill}%
\pgfsetlinewidth{1.003750pt}%
\definecolor{currentstroke}{rgb}{0.666667,0.793651,0.412698}%
\pgfsetstrokecolor{currentstroke}%
\pgfsetdash{}{0pt}%
\pgfpathmoveto{\pgfqpoint{0.636060in}{1.830025in}}%
\pgfpathlineto{\pgfqpoint{4.107310in}{1.830025in}}%
\pgfpathlineto{\pgfqpoint{4.107310in}{1.930851in}}%
\pgfpathlineto{\pgfqpoint{0.983185in}{1.930851in}}%
\pgfpathlineto{\pgfqpoint{0.636060in}{1.830025in}}%
\pgfpathclose%
\pgfusepath{stroke,fill}%
\end{pgfscope}%
\begin{pgfscope}%
\pgfpathrectangle{\pgfqpoint{0.636060in}{1.830025in}}{\pgfqpoint{3.471250in}{2.016517in}}%
\pgfusepath{clip}%
\pgfsetbuttcap%
\pgfsetmiterjoin%
\definecolor{currentfill}{rgb}{0.269841,0.666667,0.222222}%
\pgfsetfillcolor{currentfill}%
\pgfsetlinewidth{1.003750pt}%
\definecolor{currentstroke}{rgb}{0.269841,0.666667,0.222222}%
\pgfsetstrokecolor{currentstroke}%
\pgfsetdash{}{0pt}%
\pgfpathmoveto{\pgfqpoint{0.636060in}{2.031677in}}%
\pgfpathlineto{\pgfqpoint{4.107310in}{2.031677in}}%
\pgfpathlineto{\pgfqpoint{4.107310in}{2.166112in}}%
\pgfpathlineto{\pgfqpoint{0.636060in}{2.166112in}}%
\pgfpathlineto{\pgfqpoint{0.636060in}{2.031677in}}%
\pgfpathclose%
\pgfusepath{stroke,fill}%
\end{pgfscope}%
\begin{pgfscope}%
\pgfpathrectangle{\pgfqpoint{0.636060in}{1.830025in}}{\pgfqpoint{3.471250in}{2.016517in}}%
\pgfusepath{clip}%
\pgfsetbuttcap%
\pgfsetmiterjoin%
\definecolor{currentfill}{rgb}{0.079365,0.666667,0.571429}%
\pgfsetfillcolor{currentfill}%
\pgfsetlinewidth{1.003750pt}%
\definecolor{currentstroke}{rgb}{0.079365,0.666667,0.571429}%
\pgfsetstrokecolor{currentstroke}%
\pgfsetdash{}{0pt}%
\pgfpathmoveto{\pgfqpoint{0.983185in}{1.930851in}}%
\pgfpathlineto{\pgfqpoint{4.107310in}{1.930851in}}%
\pgfpathlineto{\pgfqpoint{4.107310in}{2.031677in}}%
\pgfpathlineto{\pgfqpoint{0.636060in}{2.031677in}}%
\pgfpathlineto{\pgfqpoint{0.983185in}{1.930851in}}%
\pgfpathclose%
\pgfusepath{stroke,fill}%
\end{pgfscope}%
\begin{pgfscope}%
\pgfpathrectangle{\pgfqpoint{0.636060in}{1.830025in}}{\pgfqpoint{3.471250in}{2.016517in}}%
\pgfusepath{clip}%
\pgfsetbuttcap%
\pgfsetmiterjoin%
\definecolor{currentfill}{rgb}{0.000000,0.761905,0.857143}%
\pgfsetfillcolor{currentfill}%
\pgfsetlinewidth{1.003750pt}%
\definecolor{currentstroke}{rgb}{0.000000,0.761905,0.857143}%
\pgfsetstrokecolor{currentstroke}%
\pgfsetdash{}{0pt}%
\pgfpathmoveto{\pgfqpoint{1.330310in}{2.636632in}}%
\pgfpathlineto{\pgfqpoint{2.718810in}{2.838284in}}%
\pgfpathlineto{\pgfqpoint{4.107310in}{3.241587in}}%
\pgfpathlineto{\pgfqpoint{4.107310in}{3.443239in}}%
\pgfpathlineto{\pgfqpoint{1.330310in}{2.636632in}}%
\pgfpathclose%
\pgfusepath{stroke,fill}%
\end{pgfscope}%
\begin{pgfscope}%
\pgfpathrectangle{\pgfqpoint{0.636060in}{1.830025in}}{\pgfqpoint{3.471250in}{2.016517in}}%
\pgfusepath{clip}%
\pgfsetbuttcap%
\pgfsetmiterjoin%
\definecolor{currentfill}{rgb}{0.000001,0.873016,0.650794}%
\pgfsetfillcolor{currentfill}%
\pgfsetlinewidth{1.003750pt}%
\definecolor{currentstroke}{rgb}{0.000001,0.873016,0.650794}%
\pgfsetstrokecolor{currentstroke}%
\pgfsetdash{}{0pt}%
\pgfpathmoveto{\pgfqpoint{0.636060in}{1.830025in}}%
\pgfpathlineto{\pgfqpoint{0.983185in}{1.930851in}}%
\pgfpathlineto{\pgfqpoint{0.636060in}{2.031677in}}%
\pgfpathlineto{\pgfqpoint{0.636060in}{1.830025in}}%
\pgfpathclose%
\pgfusepath{stroke,fill}%
\end{pgfscope}%
\begin{pgfscope}%
\pgfpathrectangle{\pgfqpoint{0.636060in}{1.830025in}}{\pgfqpoint{3.471250in}{2.016517in}}%
\pgfusepath{clip}%
\pgfsetbuttcap%
\pgfsetmiterjoin%
\definecolor{currentfill}{rgb}{0.269841,0.603175,0.968254}%
\pgfsetfillcolor{currentfill}%
\pgfsetlinewidth{1.003750pt}%
\definecolor{currentstroke}{rgb}{0.269841,0.603175,0.968254}%
\pgfsetstrokecolor{currentstroke}%
\pgfsetdash{}{0pt}%
\pgfpathmoveto{\pgfqpoint{0.636060in}{3.039935in}}%
\pgfpathlineto{\pgfqpoint{0.636060in}{2.636632in}}%
\pgfpathlineto{\pgfqpoint{1.330310in}{2.636632in}}%
\pgfpathlineto{\pgfqpoint{4.107310in}{3.443239in}}%
\pgfpathlineto{\pgfqpoint{4.107310in}{3.846542in}}%
\pgfpathlineto{\pgfqpoint{3.413060in}{3.846542in}}%
\pgfpathlineto{\pgfqpoint{0.636060in}{3.039935in}}%
\pgfpathclose%
\pgfusepath{stroke,fill}%
\end{pgfscope}%
\begin{pgfscope}%
\pgfpathrectangle{\pgfqpoint{0.636060in}{1.830025in}}{\pgfqpoint{3.471250in}{2.016517in}}%
\pgfusepath{clip}%
\pgfsetbuttcap%
\pgfsetmiterjoin%
\definecolor{currentfill}{rgb}{0.619048,0.476190,1.000000}%
\pgfsetfillcolor{currentfill}%
\pgfsetlinewidth{1.003750pt}%
\definecolor{currentstroke}{rgb}{0.619048,0.476190,1.000000}%
\pgfsetstrokecolor{currentstroke}%
\pgfsetdash{}{0pt}%
\pgfpathmoveto{\pgfqpoint{0.636060in}{2.434980in}}%
\pgfpathlineto{\pgfqpoint{1.330310in}{2.636632in}}%
\pgfpathlineto{\pgfqpoint{0.636060in}{2.636632in}}%
\pgfpathlineto{\pgfqpoint{0.636060in}{2.434980in}}%
\pgfpathclose%
\pgfusepath{stroke,fill}%
\end{pgfscope}%
\begin{pgfscope}%
\pgfpathrectangle{\pgfqpoint{0.636060in}{1.830025in}}{\pgfqpoint{3.471250in}{2.016517in}}%
\pgfusepath{clip}%
\pgfsetbuttcap%
\pgfsetmiterjoin%
\definecolor{currentfill}{rgb}{0.793651,0.396825,0.857143}%
\pgfsetfillcolor{currentfill}%
\pgfsetlinewidth{1.003750pt}%
\definecolor{currentstroke}{rgb}{0.793651,0.396825,0.857143}%
\pgfsetstrokecolor{currentstroke}%
\pgfsetdash{}{0pt}%
\pgfpathmoveto{\pgfqpoint{0.636060in}{3.342413in}}%
\pgfpathlineto{\pgfqpoint{0.636060in}{3.039935in}}%
\pgfpathlineto{\pgfqpoint{3.413060in}{3.846542in}}%
\pgfpathlineto{\pgfqpoint{2.024560in}{3.846542in}}%
\pgfpathlineto{\pgfqpoint{1.330310in}{3.644890in}}%
\pgfpathlineto{\pgfqpoint{0.636060in}{3.342413in}}%
\pgfpathclose%
\pgfusepath{stroke,fill}%
\end{pgfscope}%
\begin{pgfscope}%
\pgfpathrectangle{\pgfqpoint{0.636060in}{1.830025in}}{\pgfqpoint{3.471250in}{2.016517in}}%
\pgfusepath{clip}%
\pgfsetbuttcap%
\pgfsetmiterjoin%
\definecolor{currentfill}{rgb}{0.888889,0.349206,0.603175}%
\pgfsetfillcolor{currentfill}%
\pgfsetlinewidth{1.003750pt}%
\definecolor{currentstroke}{rgb}{0.888889,0.349206,0.603175}%
\pgfsetstrokecolor{currentstroke}%
\pgfsetdash{}{0pt}%
\pgfpathmoveto{\pgfqpoint{1.330310in}{3.644890in}}%
\pgfpathlineto{\pgfqpoint{2.024560in}{3.846542in}}%
\pgfpathlineto{\pgfqpoint{1.677435in}{3.846542in}}%
\pgfpathlineto{\pgfqpoint{1.330310in}{3.644890in}}%
\pgfpathclose%
\pgfusepath{stroke,fill}%
\end{pgfscope}%
\begin{pgfscope}%
\pgfpathrectangle{\pgfqpoint{0.636060in}{1.830025in}}{\pgfqpoint{3.471250in}{2.016517in}}%
\pgfusepath{clip}%
\pgfsetbuttcap%
\pgfsetmiterjoin%
\definecolor{currentfill}{rgb}{1.000000,0.587302,0.888889}%
\pgfsetfillcolor{currentfill}%
\pgfsetlinewidth{1.003750pt}%
\definecolor{currentstroke}{rgb}{1.000000,0.587302,0.888889}%
\pgfsetstrokecolor{currentstroke}%
\pgfsetdash{}{0pt}%
\pgfpathmoveto{\pgfqpoint{0.636060in}{3.342413in}}%
\pgfpathlineto{\pgfqpoint{1.330310in}{3.644890in}}%
\pgfpathlineto{\pgfqpoint{1.677435in}{3.846542in}}%
\pgfpathlineto{\pgfqpoint{0.636060in}{3.846542in}}%
\pgfpathlineto{\pgfqpoint{0.636060in}{3.342413in}}%
\pgfpathclose%
\pgfusepath{stroke,fill}%
\end{pgfscope}%
\begin{pgfscope}%
\pgfpathrectangle{\pgfqpoint{0.636060in}{1.830025in}}{\pgfqpoint{3.471250in}{2.016517in}}%
\pgfusepath{clip}%
\pgfsetbuttcap%
\pgfsetroundjoin%
\pgfsetlinewidth{0.803000pt}%
\definecolor{currentstroke}{rgb}{0.000000,0.000000,1.000000}%
\pgfsetstrokecolor{currentstroke}%
\pgfsetstrokeopacity{0.200000}%
\pgfsetdash{{8.000000pt}{4.000000pt}}{0.000000pt}%
\pgfpathmoveto{\pgfqpoint{0.636060in}{1.830025in}}%
\pgfpathlineto{\pgfqpoint{0.636060in}{3.846542in}}%
\pgfusepath{stroke}%
\end{pgfscope}%
\begin{pgfscope}%
\pgfsetbuttcap%
\pgfsetroundjoin%
\definecolor{currentfill}{rgb}{0.000000,0.000000,0.000000}%
\pgfsetfillcolor{currentfill}%
\pgfsetlinewidth{0.803000pt}%
\definecolor{currentstroke}{rgb}{0.000000,0.000000,0.000000}%
\pgfsetstrokecolor{currentstroke}%
\pgfsetdash{}{0pt}%
\pgfsys@defobject{currentmarker}{\pgfqpoint{0.000000in}{-0.048611in}}{\pgfqpoint{0.000000in}{0.000000in}}{%
\pgfpathmoveto{\pgfqpoint{0.000000in}{0.000000in}}%
\pgfpathlineto{\pgfqpoint{0.000000in}{-0.048611in}}%
\pgfusepath{stroke,fill}%
}%
\begin{pgfscope}%
\pgfsys@transformshift{0.636060in}{1.830025in}%
\pgfsys@useobject{currentmarker}{}%
\end{pgfscope}%
\end{pgfscope}%
\begin{pgfscope}%
\definecolor{textcolor}{rgb}{0.000000,0.000000,0.000000}%
\pgfsetstrokecolor{textcolor}%
\pgfsetfillcolor{textcolor}%
\pgftext[x=0.636060in,y=1.732803in,,top]{\color{textcolor}{\sffamily\fontsize{10.000000}{12.000000}\selectfont\catcode`\^=\active\def^{\ifmmode\sp\else\^{}\fi}\catcode`\%=\active\def%{\%}0}}%
\end{pgfscope}%
\begin{pgfscope}%
\pgfpathrectangle{\pgfqpoint{0.636060in}{1.830025in}}{\pgfqpoint{3.471250in}{2.016517in}}%
\pgfusepath{clip}%
\pgfsetbuttcap%
\pgfsetroundjoin%
\pgfsetlinewidth{0.803000pt}%
\definecolor{currentstroke}{rgb}{0.000000,0.000000,1.000000}%
\pgfsetstrokecolor{currentstroke}%
\pgfsetstrokeopacity{0.200000}%
\pgfsetdash{{8.000000pt}{4.000000pt}}{0.000000pt}%
\pgfpathmoveto{\pgfqpoint{1.330310in}{1.830025in}}%
\pgfpathlineto{\pgfqpoint{1.330310in}{3.846542in}}%
\pgfusepath{stroke}%
\end{pgfscope}%
\begin{pgfscope}%
\pgfsetbuttcap%
\pgfsetroundjoin%
\definecolor{currentfill}{rgb}{0.000000,0.000000,0.000000}%
\pgfsetfillcolor{currentfill}%
\pgfsetlinewidth{0.803000pt}%
\definecolor{currentstroke}{rgb}{0.000000,0.000000,0.000000}%
\pgfsetstrokecolor{currentstroke}%
\pgfsetdash{}{0pt}%
\pgfsys@defobject{currentmarker}{\pgfqpoint{0.000000in}{-0.048611in}}{\pgfqpoint{0.000000in}{0.000000in}}{%
\pgfpathmoveto{\pgfqpoint{0.000000in}{0.000000in}}%
\pgfpathlineto{\pgfqpoint{0.000000in}{-0.048611in}}%
\pgfusepath{stroke,fill}%
}%
\begin{pgfscope}%
\pgfsys@transformshift{1.330310in}{1.830025in}%
\pgfsys@useobject{currentmarker}{}%
\end{pgfscope}%
\end{pgfscope}%
\begin{pgfscope}%
\definecolor{textcolor}{rgb}{0.000000,0.000000,0.000000}%
\pgfsetstrokecolor{textcolor}%
\pgfsetfillcolor{textcolor}%
\pgftext[x=1.330310in,y=1.732803in,,top]{\color{textcolor}{\sffamily\fontsize{10.000000}{12.000000}\selectfont\catcode`\^=\active\def^{\ifmmode\sp\else\^{}\fi}\catcode`\%=\active\def%{\%}1}}%
\end{pgfscope}%
\begin{pgfscope}%
\pgfpathrectangle{\pgfqpoint{0.636060in}{1.830025in}}{\pgfqpoint{3.471250in}{2.016517in}}%
\pgfusepath{clip}%
\pgfsetbuttcap%
\pgfsetroundjoin%
\pgfsetlinewidth{0.803000pt}%
\definecolor{currentstroke}{rgb}{0.000000,0.000000,1.000000}%
\pgfsetstrokecolor{currentstroke}%
\pgfsetstrokeopacity{0.200000}%
\pgfsetdash{{8.000000pt}{4.000000pt}}{0.000000pt}%
\pgfpathmoveto{\pgfqpoint{2.024560in}{1.830025in}}%
\pgfpathlineto{\pgfqpoint{2.024560in}{3.846542in}}%
\pgfusepath{stroke}%
\end{pgfscope}%
\begin{pgfscope}%
\pgfsetbuttcap%
\pgfsetroundjoin%
\definecolor{currentfill}{rgb}{0.000000,0.000000,0.000000}%
\pgfsetfillcolor{currentfill}%
\pgfsetlinewidth{0.803000pt}%
\definecolor{currentstroke}{rgb}{0.000000,0.000000,0.000000}%
\pgfsetstrokecolor{currentstroke}%
\pgfsetdash{}{0pt}%
\pgfsys@defobject{currentmarker}{\pgfqpoint{0.000000in}{-0.048611in}}{\pgfqpoint{0.000000in}{0.000000in}}{%
\pgfpathmoveto{\pgfqpoint{0.000000in}{0.000000in}}%
\pgfpathlineto{\pgfqpoint{0.000000in}{-0.048611in}}%
\pgfusepath{stroke,fill}%
}%
\begin{pgfscope}%
\pgfsys@transformshift{2.024560in}{1.830025in}%
\pgfsys@useobject{currentmarker}{}%
\end{pgfscope}%
\end{pgfscope}%
\begin{pgfscope}%
\definecolor{textcolor}{rgb}{0.000000,0.000000,0.000000}%
\pgfsetstrokecolor{textcolor}%
\pgfsetfillcolor{textcolor}%
\pgftext[x=2.024560in,y=1.732803in,,top]{\color{textcolor}{\sffamily\fontsize{10.000000}{12.000000}\selectfont\catcode`\^=\active\def^{\ifmmode\sp\else\^{}\fi}\catcode`\%=\active\def%{\%}2}}%
\end{pgfscope}%
\begin{pgfscope}%
\pgfpathrectangle{\pgfqpoint{0.636060in}{1.830025in}}{\pgfqpoint{3.471250in}{2.016517in}}%
\pgfusepath{clip}%
\pgfsetbuttcap%
\pgfsetroundjoin%
\pgfsetlinewidth{0.803000pt}%
\definecolor{currentstroke}{rgb}{0.000000,0.000000,1.000000}%
\pgfsetstrokecolor{currentstroke}%
\pgfsetstrokeopacity{0.200000}%
\pgfsetdash{{8.000000pt}{4.000000pt}}{0.000000pt}%
\pgfpathmoveto{\pgfqpoint{2.718810in}{1.830025in}}%
\pgfpathlineto{\pgfqpoint{2.718810in}{3.846542in}}%
\pgfusepath{stroke}%
\end{pgfscope}%
\begin{pgfscope}%
\pgfsetbuttcap%
\pgfsetroundjoin%
\definecolor{currentfill}{rgb}{0.000000,0.000000,0.000000}%
\pgfsetfillcolor{currentfill}%
\pgfsetlinewidth{0.803000pt}%
\definecolor{currentstroke}{rgb}{0.000000,0.000000,0.000000}%
\pgfsetstrokecolor{currentstroke}%
\pgfsetdash{}{0pt}%
\pgfsys@defobject{currentmarker}{\pgfqpoint{0.000000in}{-0.048611in}}{\pgfqpoint{0.000000in}{0.000000in}}{%
\pgfpathmoveto{\pgfqpoint{0.000000in}{0.000000in}}%
\pgfpathlineto{\pgfqpoint{0.000000in}{-0.048611in}}%
\pgfusepath{stroke,fill}%
}%
\begin{pgfscope}%
\pgfsys@transformshift{2.718810in}{1.830025in}%
\pgfsys@useobject{currentmarker}{}%
\end{pgfscope}%
\end{pgfscope}%
\begin{pgfscope}%
\definecolor{textcolor}{rgb}{0.000000,0.000000,0.000000}%
\pgfsetstrokecolor{textcolor}%
\pgfsetfillcolor{textcolor}%
\pgftext[x=2.718810in,y=1.732803in,,top]{\color{textcolor}{\sffamily\fontsize{10.000000}{12.000000}\selectfont\catcode`\^=\active\def^{\ifmmode\sp\else\^{}\fi}\catcode`\%=\active\def%{\%}3}}%
\end{pgfscope}%
\begin{pgfscope}%
\pgfpathrectangle{\pgfqpoint{0.636060in}{1.830025in}}{\pgfqpoint{3.471250in}{2.016517in}}%
\pgfusepath{clip}%
\pgfsetbuttcap%
\pgfsetroundjoin%
\pgfsetlinewidth{0.803000pt}%
\definecolor{currentstroke}{rgb}{0.000000,0.000000,1.000000}%
\pgfsetstrokecolor{currentstroke}%
\pgfsetstrokeopacity{0.200000}%
\pgfsetdash{{8.000000pt}{4.000000pt}}{0.000000pt}%
\pgfpathmoveto{\pgfqpoint{3.413060in}{1.830025in}}%
\pgfpathlineto{\pgfqpoint{3.413060in}{3.846542in}}%
\pgfusepath{stroke}%
\end{pgfscope}%
\begin{pgfscope}%
\pgfsetbuttcap%
\pgfsetroundjoin%
\definecolor{currentfill}{rgb}{0.000000,0.000000,0.000000}%
\pgfsetfillcolor{currentfill}%
\pgfsetlinewidth{0.803000pt}%
\definecolor{currentstroke}{rgb}{0.000000,0.000000,0.000000}%
\pgfsetstrokecolor{currentstroke}%
\pgfsetdash{}{0pt}%
\pgfsys@defobject{currentmarker}{\pgfqpoint{0.000000in}{-0.048611in}}{\pgfqpoint{0.000000in}{0.000000in}}{%
\pgfpathmoveto{\pgfqpoint{0.000000in}{0.000000in}}%
\pgfpathlineto{\pgfqpoint{0.000000in}{-0.048611in}}%
\pgfusepath{stroke,fill}%
}%
\begin{pgfscope}%
\pgfsys@transformshift{3.413060in}{1.830025in}%
\pgfsys@useobject{currentmarker}{}%
\end{pgfscope}%
\end{pgfscope}%
\begin{pgfscope}%
\definecolor{textcolor}{rgb}{0.000000,0.000000,0.000000}%
\pgfsetstrokecolor{textcolor}%
\pgfsetfillcolor{textcolor}%
\pgftext[x=3.413060in,y=1.732803in,,top]{\color{textcolor}{\sffamily\fontsize{10.000000}{12.000000}\selectfont\catcode`\^=\active\def^{\ifmmode\sp\else\^{}\fi}\catcode`\%=\active\def%{\%}4}}%
\end{pgfscope}%
\begin{pgfscope}%
\pgfpathrectangle{\pgfqpoint{0.636060in}{1.830025in}}{\pgfqpoint{3.471250in}{2.016517in}}%
\pgfusepath{clip}%
\pgfsetbuttcap%
\pgfsetroundjoin%
\pgfsetlinewidth{0.803000pt}%
\definecolor{currentstroke}{rgb}{0.000000,0.000000,1.000000}%
\pgfsetstrokecolor{currentstroke}%
\pgfsetstrokeopacity{0.200000}%
\pgfsetdash{{8.000000pt}{4.000000pt}}{0.000000pt}%
\pgfpathmoveto{\pgfqpoint{4.107310in}{1.830025in}}%
\pgfpathlineto{\pgfqpoint{4.107310in}{3.846542in}}%
\pgfusepath{stroke}%
\end{pgfscope}%
\begin{pgfscope}%
\pgfsetbuttcap%
\pgfsetroundjoin%
\definecolor{currentfill}{rgb}{0.000000,0.000000,0.000000}%
\pgfsetfillcolor{currentfill}%
\pgfsetlinewidth{0.803000pt}%
\definecolor{currentstroke}{rgb}{0.000000,0.000000,0.000000}%
\pgfsetstrokecolor{currentstroke}%
\pgfsetdash{}{0pt}%
\pgfsys@defobject{currentmarker}{\pgfqpoint{0.000000in}{-0.048611in}}{\pgfqpoint{0.000000in}{0.000000in}}{%
\pgfpathmoveto{\pgfqpoint{0.000000in}{0.000000in}}%
\pgfpathlineto{\pgfqpoint{0.000000in}{-0.048611in}}%
\pgfusepath{stroke,fill}%
}%
\begin{pgfscope}%
\pgfsys@transformshift{4.107310in}{1.830025in}%
\pgfsys@useobject{currentmarker}{}%
\end{pgfscope}%
\end{pgfscope}%
\begin{pgfscope}%
\definecolor{textcolor}{rgb}{0.000000,0.000000,0.000000}%
\pgfsetstrokecolor{textcolor}%
\pgfsetfillcolor{textcolor}%
\pgftext[x=4.107310in,y=1.732803in,,top]{\color{textcolor}{\sffamily\fontsize{10.000000}{12.000000}\selectfont\catcode`\^=\active\def^{\ifmmode\sp\else\^{}\fi}\catcode`\%=\active\def%{\%}5}}%
\end{pgfscope}%
\begin{pgfscope}%
\definecolor{textcolor}{rgb}{0.000000,0.000000,0.000000}%
\pgfsetstrokecolor{textcolor}%
\pgfsetfillcolor{textcolor}%
    \pgftext[x=2.371685in,y=1.553791in,,top]{\color{textcolor}{\sffamily\fontsize{10.000000}{12.000000}\selectfont\catcode`\^=\active\def^{\ifmmode\sp\else\^{}\fi}\catcode`\%=\active\def%{\%}\(\delta_{\mathrm{Dup}} = \delta_{\mathrm{Cut}}\)}}%
\end{pgfscope}%
\begin{pgfscope}%
\pgfpathrectangle{\pgfqpoint{0.636060in}{1.830025in}}{\pgfqpoint{3.471250in}{2.016517in}}%
\pgfusepath{clip}%
\pgfsetbuttcap%
\pgfsetroundjoin%
\pgfsetlinewidth{0.803000pt}%
\definecolor{currentstroke}{rgb}{0.000000,0.000000,1.000000}%
\pgfsetstrokecolor{currentstroke}%
\pgfsetstrokeopacity{0.200000}%
\pgfsetdash{{8.000000pt}{4.000000pt}}{0.000000pt}%
\pgfpathmoveto{\pgfqpoint{0.636060in}{1.830025in}}%
\pgfpathlineto{\pgfqpoint{4.107310in}{1.830025in}}%
\pgfusepath{stroke}%
\end{pgfscope}%
\begin{pgfscope}%
\pgfsetbuttcap%
\pgfsetroundjoin%
\definecolor{currentfill}{rgb}{0.000000,0.000000,0.000000}%
\pgfsetfillcolor{currentfill}%
\pgfsetlinewidth{0.803000pt}%
\definecolor{currentstroke}{rgb}{0.000000,0.000000,0.000000}%
\pgfsetstrokecolor{currentstroke}%
\pgfsetdash{}{0pt}%
\pgfsys@defobject{currentmarker}{\pgfqpoint{-0.048611in}{0.000000in}}{\pgfqpoint{-0.000000in}{0.000000in}}{%
\pgfpathmoveto{\pgfqpoint{-0.000000in}{0.000000in}}%
\pgfpathlineto{\pgfqpoint{-0.048611in}{0.000000in}}%
\pgfusepath{stroke,fill}%
}%
\begin{pgfscope}%
\pgfsys@transformshift{0.636060in}{1.830025in}%
\pgfsys@useobject{currentmarker}{}%
\end{pgfscope}%
\end{pgfscope}%
\begin{pgfscope}%
\definecolor{textcolor}{rgb}{0.000000,0.000000,0.000000}%
\pgfsetstrokecolor{textcolor}%
\pgfsetfillcolor{textcolor}%
\pgftext[x=0.469393in, y=1.781800in, left, base]{\color{textcolor}{\sffamily\fontsize{10.000000}{12.000000}\selectfont\catcode`\^=\active\def^{\ifmmode\sp\else\^{}\fi}\catcode`\%=\active\def%{\%}0}}%
\end{pgfscope}%
\begin{pgfscope}%
\pgfpathrectangle{\pgfqpoint{0.636060in}{1.830025in}}{\pgfqpoint{3.471250in}{2.016517in}}%
\pgfusepath{clip}%
\pgfsetbuttcap%
\pgfsetroundjoin%
\pgfsetlinewidth{0.803000pt}%
\definecolor{currentstroke}{rgb}{0.000000,0.000000,1.000000}%
\pgfsetstrokecolor{currentstroke}%
\pgfsetstrokeopacity{0.200000}%
\pgfsetdash{{8.000000pt}{4.000000pt}}{0.000000pt}%
\pgfpathmoveto{\pgfqpoint{0.636060in}{2.233329in}}%
\pgfpathlineto{\pgfqpoint{4.107310in}{2.233329in}}%
\pgfusepath{stroke}%
\end{pgfscope}%
\begin{pgfscope}%
\pgfsetbuttcap%
\pgfsetroundjoin%
\definecolor{currentfill}{rgb}{0.000000,0.000000,0.000000}%
\pgfsetfillcolor{currentfill}%
\pgfsetlinewidth{0.803000pt}%
\definecolor{currentstroke}{rgb}{0.000000,0.000000,0.000000}%
\pgfsetstrokecolor{currentstroke}%
\pgfsetdash{}{0pt}%
\pgfsys@defobject{currentmarker}{\pgfqpoint{-0.048611in}{0.000000in}}{\pgfqpoint{-0.000000in}{0.000000in}}{%
\pgfpathmoveto{\pgfqpoint{-0.000000in}{0.000000in}}%
\pgfpathlineto{\pgfqpoint{-0.048611in}{0.000000in}}%
\pgfusepath{stroke,fill}%
}%
\begin{pgfscope}%
\pgfsys@transformshift{0.636060in}{2.233329in}%
\pgfsys@useobject{currentmarker}{}%
\end{pgfscope}%
\end{pgfscope}%
\begin{pgfscope}%
\definecolor{textcolor}{rgb}{0.000000,0.000000,0.000000}%
\pgfsetstrokecolor{textcolor}%
\pgfsetfillcolor{textcolor}%
\pgftext[x=0.469393in, y=2.185103in, left, base]{\color{textcolor}{\sffamily\fontsize{10.000000}{12.000000}\selectfont\catcode`\^=\active\def^{\ifmmode\sp\else\^{}\fi}\catcode`\%=\active\def%{\%}2}}%
\end{pgfscope}%
\begin{pgfscope}%
\pgfpathrectangle{\pgfqpoint{0.636060in}{1.830025in}}{\pgfqpoint{3.471250in}{2.016517in}}%
\pgfusepath{clip}%
\pgfsetbuttcap%
\pgfsetroundjoin%
\pgfsetlinewidth{0.803000pt}%
\definecolor{currentstroke}{rgb}{0.000000,0.000000,1.000000}%
\pgfsetstrokecolor{currentstroke}%
\pgfsetstrokeopacity{0.200000}%
\pgfsetdash{{8.000000pt}{4.000000pt}}{0.000000pt}%
\pgfpathmoveto{\pgfqpoint{0.636060in}{2.636632in}}%
\pgfpathlineto{\pgfqpoint{4.107310in}{2.636632in}}%
\pgfusepath{stroke}%
\end{pgfscope}%
\begin{pgfscope}%
\pgfsetbuttcap%
\pgfsetroundjoin%
\definecolor{currentfill}{rgb}{0.000000,0.000000,0.000000}%
\pgfsetfillcolor{currentfill}%
\pgfsetlinewidth{0.803000pt}%
\definecolor{currentstroke}{rgb}{0.000000,0.000000,0.000000}%
\pgfsetstrokecolor{currentstroke}%
\pgfsetdash{}{0pt}%
\pgfsys@defobject{currentmarker}{\pgfqpoint{-0.048611in}{0.000000in}}{\pgfqpoint{-0.000000in}{0.000000in}}{%
\pgfpathmoveto{\pgfqpoint{-0.000000in}{0.000000in}}%
\pgfpathlineto{\pgfqpoint{-0.048611in}{0.000000in}}%
\pgfusepath{stroke,fill}%
}%
\begin{pgfscope}%
\pgfsys@transformshift{0.636060in}{2.636632in}%
\pgfsys@useobject{currentmarker}{}%
\end{pgfscope}%
\end{pgfscope}%
\begin{pgfscope}%
\definecolor{textcolor}{rgb}{0.000000,0.000000,0.000000}%
\pgfsetstrokecolor{textcolor}%
\pgfsetfillcolor{textcolor}%
\pgftext[x=0.469393in, y=2.588407in, left, base]{\color{textcolor}{\sffamily\fontsize{10.000000}{12.000000}\selectfont\catcode`\^=\active\def^{\ifmmode\sp\else\^{}\fi}\catcode`\%=\active\def%{\%}4}}%
\end{pgfscope}%
\begin{pgfscope}%
\pgfpathrectangle{\pgfqpoint{0.636060in}{1.830025in}}{\pgfqpoint{3.471250in}{2.016517in}}%
\pgfusepath{clip}%
\pgfsetbuttcap%
\pgfsetroundjoin%
\pgfsetlinewidth{0.803000pt}%
\definecolor{currentstroke}{rgb}{0.000000,0.000000,1.000000}%
\pgfsetstrokecolor{currentstroke}%
\pgfsetstrokeopacity{0.200000}%
\pgfsetdash{{8.000000pt}{4.000000pt}}{0.000000pt}%
\pgfpathmoveto{\pgfqpoint{0.636060in}{3.039935in}}%
\pgfpathlineto{\pgfqpoint{4.107310in}{3.039935in}}%
\pgfusepath{stroke}%
\end{pgfscope}%
\begin{pgfscope}%
\pgfsetbuttcap%
\pgfsetroundjoin%
\definecolor{currentfill}{rgb}{0.000000,0.000000,0.000000}%
\pgfsetfillcolor{currentfill}%
\pgfsetlinewidth{0.803000pt}%
\definecolor{currentstroke}{rgb}{0.000000,0.000000,0.000000}%
\pgfsetstrokecolor{currentstroke}%
\pgfsetdash{}{0pt}%
\pgfsys@defobject{currentmarker}{\pgfqpoint{-0.048611in}{0.000000in}}{\pgfqpoint{-0.000000in}{0.000000in}}{%
\pgfpathmoveto{\pgfqpoint{-0.000000in}{0.000000in}}%
\pgfpathlineto{\pgfqpoint{-0.048611in}{0.000000in}}%
\pgfusepath{stroke,fill}%
}%
\begin{pgfscope}%
\pgfsys@transformshift{0.636060in}{3.039935in}%
\pgfsys@useobject{currentmarker}{}%
\end{pgfscope}%
\end{pgfscope}%
\begin{pgfscope}%
\definecolor{textcolor}{rgb}{0.000000,0.000000,0.000000}%
\pgfsetstrokecolor{textcolor}%
\pgfsetfillcolor{textcolor}%
\pgftext[x=0.469393in, y=2.991710in, left, base]{\color{textcolor}{\sffamily\fontsize{10.000000}{12.000000}\selectfont\catcode`\^=\active\def^{\ifmmode\sp\else\^{}\fi}\catcode`\%=\active\def%{\%}6}}%
\end{pgfscope}%
\begin{pgfscope}%
\pgfpathrectangle{\pgfqpoint{0.636060in}{1.830025in}}{\pgfqpoint{3.471250in}{2.016517in}}%
\pgfusepath{clip}%
\pgfsetbuttcap%
\pgfsetroundjoin%
\pgfsetlinewidth{0.803000pt}%
\definecolor{currentstroke}{rgb}{0.000000,0.000000,1.000000}%
\pgfsetstrokecolor{currentstroke}%
\pgfsetstrokeopacity{0.200000}%
\pgfsetdash{{8.000000pt}{4.000000pt}}{0.000000pt}%
\pgfpathmoveto{\pgfqpoint{0.636060in}{3.443239in}}%
\pgfpathlineto{\pgfqpoint{4.107310in}{3.443239in}}%
\pgfusepath{stroke}%
\end{pgfscope}%
\begin{pgfscope}%
\pgfsetbuttcap%
\pgfsetroundjoin%
\definecolor{currentfill}{rgb}{0.000000,0.000000,0.000000}%
\pgfsetfillcolor{currentfill}%
\pgfsetlinewidth{0.803000pt}%
\definecolor{currentstroke}{rgb}{0.000000,0.000000,0.000000}%
\pgfsetstrokecolor{currentstroke}%
\pgfsetdash{}{0pt}%
\pgfsys@defobject{currentmarker}{\pgfqpoint{-0.048611in}{0.000000in}}{\pgfqpoint{-0.000000in}{0.000000in}}{%
\pgfpathmoveto{\pgfqpoint{-0.000000in}{0.000000in}}%
\pgfpathlineto{\pgfqpoint{-0.048611in}{0.000000in}}%
\pgfusepath{stroke,fill}%
}%
\begin{pgfscope}%
\pgfsys@transformshift{0.636060in}{3.443239in}%
\pgfsys@useobject{currentmarker}{}%
\end{pgfscope}%
\end{pgfscope}%
\begin{pgfscope}%
\definecolor{textcolor}{rgb}{0.000000,0.000000,0.000000}%
\pgfsetstrokecolor{textcolor}%
\pgfsetfillcolor{textcolor}%
\pgftext[x=0.469393in, y=3.395013in, left, base]{\color{textcolor}{\sffamily\fontsize{10.000000}{12.000000}\selectfont\catcode`\^=\active\def^{\ifmmode\sp\else\^{}\fi}\catcode`\%=\active\def%{\%}8}}%
\end{pgfscope}%
\begin{pgfscope}%
\pgfpathrectangle{\pgfqpoint{0.636060in}{1.830025in}}{\pgfqpoint{3.471250in}{2.016517in}}%
\pgfusepath{clip}%
\pgfsetbuttcap%
\pgfsetroundjoin%
\pgfsetlinewidth{0.803000pt}%
\definecolor{currentstroke}{rgb}{0.000000,0.000000,1.000000}%
\pgfsetstrokecolor{currentstroke}%
\pgfsetstrokeopacity{0.200000}%
\pgfsetdash{{8.000000pt}{4.000000pt}}{0.000000pt}%
\pgfpathmoveto{\pgfqpoint{0.636060in}{3.846542in}}%
\pgfpathlineto{\pgfqpoint{4.107310in}{3.846542in}}%
\pgfusepath{stroke}%
\end{pgfscope}%
\begin{pgfscope}%
\pgfsetbuttcap%
\pgfsetroundjoin%
\definecolor{currentfill}{rgb}{0.000000,0.000000,0.000000}%
\pgfsetfillcolor{currentfill}%
\pgfsetlinewidth{0.803000pt}%
\definecolor{currentstroke}{rgb}{0.000000,0.000000,0.000000}%
\pgfsetstrokecolor{currentstroke}%
\pgfsetdash{}{0pt}%
\pgfsys@defobject{currentmarker}{\pgfqpoint{-0.048611in}{0.000000in}}{\pgfqpoint{-0.000000in}{0.000000in}}{%
\pgfpathmoveto{\pgfqpoint{-0.000000in}{0.000000in}}%
\pgfpathlineto{\pgfqpoint{-0.048611in}{0.000000in}}%
\pgfusepath{stroke,fill}%
}%
\begin{pgfscope}%
\pgfsys@transformshift{0.636060in}{3.846542in}%
\pgfsys@useobject{currentmarker}{}%
\end{pgfscope}%
\end{pgfscope}%
\begin{pgfscope}%
\definecolor{textcolor}{rgb}{0.000000,0.000000,0.000000}%
\pgfsetstrokecolor{textcolor}%
\pgfsetfillcolor{textcolor}%
\pgftext[x=0.399949in, y=3.798317in, left, base]{\color{textcolor}{\sffamily\fontsize{10.000000}{12.000000}\selectfont\catcode`\^=\active\def^{\ifmmode\sp\else\^{}\fi}\catcode`\%=\active\def%{\%}10}}%
\end{pgfscope}%
\begin{pgfscope}%
\definecolor{textcolor}{rgb}{0.000000,0.000000,0.000000}%
\pgfsetstrokecolor{textcolor}%
\pgfsetfillcolor{textcolor}%
    \pgftext[x=0.344393in,y=2.838284in,,bottom,rotate=90.000000]{\color{textcolor}{\sffamily\fontsize{10.000000}{12.000000}\selectfont\catcode`\^=\active\def^{\ifmmode\sp\else\^{}\fi}\catcode`\%=\active\def%{\%}\(\delta_{\mathrm{TrDup}} = \delta_{\mathrm{TrCut}}\)}}%
\end{pgfscope}%
\begin{pgfscope}%
\pgfpathrectangle{\pgfqpoint{0.636060in}{1.830025in}}{\pgfqpoint{3.471250in}{2.016517in}}%
\pgfusepath{clip}%
\pgfsetrectcap%
\pgfsetroundjoin%
\pgfsetlinewidth{1.505625pt}%
\definecolor{currentstroke}{rgb}{0.222222,0.190476,0.698413}%
\pgfsetstrokecolor{currentstroke}%
\pgfsetdash{}{0pt}%
\pgfpathmoveto{\pgfqpoint{0.636060in}{2.434980in}}%
\pgfpathlineto{\pgfqpoint{4.107310in}{3.443239in}}%
\pgfusepath{stroke}%
\end{pgfscope}%
\begin{pgfscope}%
\pgfpathrectangle{\pgfqpoint{0.636060in}{1.830025in}}{\pgfqpoint{3.471250in}{2.016517in}}%
\pgfusepath{clip}%
\pgfsetrectcap%
\pgfsetroundjoin%
\pgfsetlinewidth{1.505625pt}%
\definecolor{currentstroke}{rgb}{0.507937,0.031746,0.095238}%
\pgfsetstrokecolor{currentstroke}%
\pgfsetdash{}{0pt}%
\pgfpathmoveto{\pgfqpoint{0.636060in}{3.039935in}}%
\pgfpathlineto{\pgfqpoint{3.447489in}{3.856542in}}%
\pgfusepath{stroke}%
\end{pgfscope}%
\begin{pgfscope}%
\pgfsetrectcap%
\pgfsetmiterjoin%
\pgfsetlinewidth{0.803000pt}%
\definecolor{currentstroke}{rgb}{0.000000,0.000000,0.000000}%
\pgfsetstrokecolor{currentstroke}%
\pgfsetdash{}{0pt}%
\pgfpathmoveto{\pgfqpoint{0.636060in}{1.830025in}}%
\pgfpathlineto{\pgfqpoint{0.636060in}{3.846542in}}%
\pgfusepath{stroke}%
\end{pgfscope}%
\begin{pgfscope}%
\pgfsetrectcap%
\pgfsetmiterjoin%
\pgfsetlinewidth{0.803000pt}%
\definecolor{currentstroke}{rgb}{0.000000,0.000000,0.000000}%
\pgfsetstrokecolor{currentstroke}%
\pgfsetdash{}{0pt}%
\pgfpathmoveto{\pgfqpoint{4.107310in}{1.830025in}}%
\pgfpathlineto{\pgfqpoint{4.107310in}{3.846542in}}%
\pgfusepath{stroke}%
\end{pgfscope}%
\begin{pgfscope}%
\pgfsetrectcap%
\pgfsetmiterjoin%
\pgfsetlinewidth{0.803000pt}%
\definecolor{currentstroke}{rgb}{0.000000,0.000000,0.000000}%
\pgfsetstrokecolor{currentstroke}%
\pgfsetdash{}{0pt}%
\pgfpathmoveto{\pgfqpoint{0.636060in}{1.830025in}}%
\pgfpathlineto{\pgfqpoint{4.107310in}{1.830025in}}%
\pgfusepath{stroke}%
\end{pgfscope}%
\begin{pgfscope}%
\pgfsetrectcap%
\pgfsetmiterjoin%
\pgfsetlinewidth{0.803000pt}%
\definecolor{currentstroke}{rgb}{0.000000,0.000000,0.000000}%
\pgfsetstrokecolor{currentstroke}%
\pgfsetdash{}{0pt}%
\pgfpathmoveto{\pgfqpoint{0.636060in}{3.846542in}}%
\pgfpathlineto{\pgfqpoint{4.107310in}{3.846542in}}%
\pgfusepath{stroke}%
\end{pgfscope}%
\begin{pgfscope}%
\pgfpathrectangle{\pgfqpoint{0.636060in}{1.830025in}}{\pgfqpoint{3.471250in}{2.016517in}}%
\pgfusepath{clip}%
\pgfsetbuttcap%
\pgfsetroundjoin%
\definecolor{currentfill}{rgb}{0.000000,0.317460,0.000000}%
\pgfsetfillcolor{currentfill}%
\pgfsetlinewidth{2.007500pt}%
\definecolor{currentstroke}{rgb}{0.000000,0.317460,0.000000}%
\pgfsetstrokecolor{currentstroke}%
\pgfsetdash{}{0pt}%
\pgfsys@defobject{currentmarker}{\pgfqpoint{-0.041667in}{-0.041667in}}{\pgfqpoint{0.041667in}{0.041667in}}{%
\pgfpathmoveto{\pgfqpoint{-0.041667in}{-0.041667in}}%
\pgfpathlineto{\pgfqpoint{0.041667in}{0.041667in}}%
\pgfpathmoveto{\pgfqpoint{-0.041667in}{0.041667in}}%
\pgfpathlineto{\pgfqpoint{0.041667in}{-0.041667in}}%
\pgfusepath{stroke,fill}%
}%
\begin{pgfscope}%
\pgfsys@transformshift{1.330310in}{2.636632in}%
\pgfsys@useobject{currentmarker}{}%
\end{pgfscope}%
\end{pgfscope}%
\begin{pgfscope}%
\pgfsetbuttcap%
\pgfsetmiterjoin%
\definecolor{currentfill}{rgb}{0.904762,0.365079,0.301587}%
\pgfsetfillcolor{currentfill}%
\pgfsetlinewidth{1.003750pt}%
\definecolor{currentstroke}{rgb}{0.904762,0.365079,0.301587}%
\pgfsetstrokecolor{currentstroke}%
\pgfsetdash{}{0pt}%
\pgfpathmoveto{\pgfqpoint{0.138889in}{1.186727in}}%
\pgfpathlineto{\pgfqpoint{0.333333in}{1.186727in}}%
\pgfpathlineto{\pgfqpoint{0.333333in}{1.254782in}}%
\pgfpathlineto{\pgfqpoint{0.138889in}{1.254782in}}%
\pgfpathlineto{\pgfqpoint{0.138889in}{1.186727in}}%
\pgfpathclose%
\pgfusepath{stroke,fill}%
\end{pgfscope}%
\begin{pgfscope}%
\definecolor{textcolor}{rgb}{0.000000,0.000000,0.000000}%
\pgfsetstrokecolor{textcolor}%
\pgfsetfillcolor{textcolor}%
\pgftext[x=0.411111in, y=1.246140in, left, base]{\color{textcolor}{\sffamily\fontsize{7.000000}{8.400000}\selectfont\catcode`\^=\active\def^{\ifmmode\sp\else\^{}\fi}\catcode`\%=\active\def%{\%}(0, 6, 18)}}%
\end{pgfscope}%
\begin{pgfscope}%
\definecolor{textcolor}{rgb}{0.000000,0.000000,0.000000}%
\pgfsetstrokecolor{textcolor}%
\pgfsetfillcolor{textcolor}%
\pgftext[x=0.411111in, y=1.140816in, left, base]{\color{textcolor}{\sffamily\fontsize{7.000000}{8.400000}\selectfont\catcode`\^=\active\def^{\ifmmode\sp\else\^{}\fi}\catcode`\%=\active\def%{\%}\(n=3{,}278{,}464\)}}%
\end{pgfscope}%
\begin{pgfscope}%
\pgfsetbuttcap%
\pgfsetmiterjoin%
\definecolor{currentfill}{rgb}{1.000000,0.603175,0.444445}%
\pgfsetfillcolor{currentfill}%
\pgfsetlinewidth{1.003750pt}%
\definecolor{currentstroke}{rgb}{1.000000,0.603175,0.444445}%
\pgfsetstrokecolor{currentstroke}%
\pgfsetdash{}{0pt}%
\pgfpathmoveto{\pgfqpoint{0.138889in}{0.940971in}}%
\pgfpathlineto{\pgfqpoint{0.333333in}{0.940971in}}%
\pgfpathlineto{\pgfqpoint{0.333333in}{1.009027in}}%
\pgfpathlineto{\pgfqpoint{0.138889in}{1.009027in}}%
\pgfpathlineto{\pgfqpoint{0.138889in}{0.940971in}}%
\pgfpathclose%
\pgfusepath{stroke,fill}%
\end{pgfscope}%
\begin{pgfscope}%
\definecolor{textcolor}{rgb}{0.000000,0.000000,0.000000}%
\pgfsetstrokecolor{textcolor}%
\pgfsetfillcolor{textcolor}%
\pgftext[x=0.411111in, y=1.000384in, left, base]{\color{textcolor}{\sffamily\fontsize{7.000000}{8.400000}\selectfont\catcode`\^=\active\def^{\ifmmode\sp\else\^{}\fi}\catcode`\%=\active\def%{\%}(0, 7, 13)}}%
\end{pgfscope}%
\begin{pgfscope}%
\definecolor{textcolor}{rgb}{0.000000,0.000000,0.000000}%
\pgfsetstrokecolor{textcolor}%
\pgfsetfillcolor{textcolor}%
\pgftext[x=0.411111in, y=0.895061in, left, base]{\color{textcolor}{\sffamily\fontsize{7.000000}{8.400000}\selectfont\catcode`\^=\active\def^{\ifmmode\sp\else\^{}\fi}\catcode`\%=\active\def%{\%}\(n=7{,}447{,}552\)}}%
\end{pgfscope}%
\begin{pgfscope}%
\pgfsetbuttcap%
\pgfsetmiterjoin%
\definecolor{currentfill}{rgb}{0.761905,0.507936,0.206349}%
\pgfsetfillcolor{currentfill}%
\pgfsetlinewidth{1.003750pt}%
\definecolor{currentstroke}{rgb}{0.761905,0.507936,0.206349}%
\pgfsetstrokecolor{currentstroke}%
\pgfsetdash{}{0pt}%
\pgfpathmoveto{\pgfqpoint{0.138889in}{0.695215in}}%
\pgfpathlineto{\pgfqpoint{0.333333in}{0.695215in}}%
\pgfpathlineto{\pgfqpoint{0.333333in}{0.763271in}}%
\pgfpathlineto{\pgfqpoint{0.138889in}{0.763271in}}%
\pgfpathlineto{\pgfqpoint{0.138889in}{0.695215in}}%
\pgfpathclose%
\pgfusepath{stroke,fill}%
\end{pgfscope}%
\begin{pgfscope}%
\definecolor{textcolor}{rgb}{0.000000,0.000000,0.000000}%
\pgfsetstrokecolor{textcolor}%
\pgfsetfillcolor{textcolor}%
\pgftext[x=0.411111in, y=0.754629in, left, base]{\color{textcolor}{\sffamily\fontsize{7.000000}{8.400000}\selectfont\catcode`\^=\active\def^{\ifmmode\sp\else\^{}\fi}\catcode`\%=\active\def%{\%}(0, 8, 9)}}%
\end{pgfscope}%
\begin{pgfscope}%
\definecolor{textcolor}{rgb}{0.000000,0.000000,0.000000}%
\pgfsetstrokecolor{textcolor}%
\pgfsetfillcolor{textcolor}%
\pgftext[x=0.411111in, y=0.649305in, left, base]{\color{textcolor}{\sffamily\fontsize{7.000000}{8.400000}\selectfont\catcode`\^=\active\def^{\ifmmode\sp\else\^{}\fi}\catcode`\%=\active\def%{\%}\(n=135{,}168\)}}%
\end{pgfscope}%
\begin{pgfscope}%
\pgfsetbuttcap%
\pgfsetmiterjoin%
\definecolor{currentfill}{rgb}{0.936508,0.698413,0.063492}%
\pgfsetfillcolor{currentfill}%
\pgfsetlinewidth{1.003750pt}%
\definecolor{currentstroke}{rgb}{0.936508,0.698413,0.063492}%
\pgfsetstrokecolor{currentstroke}%
\pgfsetdash{}{0pt}%
\pgfpathmoveto{\pgfqpoint{0.138889in}{0.449459in}}%
\pgfpathlineto{\pgfqpoint{0.333333in}{0.449459in}}%
\pgfpathlineto{\pgfqpoint{0.333333in}{0.517515in}}%
\pgfpathlineto{\pgfqpoint{0.138889in}{0.517515in}}%
\pgfpathlineto{\pgfqpoint{0.138889in}{0.449459in}}%
\pgfpathclose%
\pgfusepath{stroke,fill}%
\end{pgfscope}%
\begin{pgfscope}%
\definecolor{textcolor}{rgb}{0.000000,0.000000,0.000000}%
\pgfsetstrokecolor{textcolor}%
\pgfsetfillcolor{textcolor}%
\pgftext[x=0.411111in, y=0.508873in, left, base]{\color{textcolor}{\sffamily\fontsize{7.000000}{8.400000}\selectfont\catcode`\^=\active\def^{\ifmmode\sp\else\^{}\fi}\catcode`\%=\active\def%{\%}(0, 9, 7)}}%
\end{pgfscope}%
\begin{pgfscope}%
\definecolor{textcolor}{rgb}{0.000000,0.000000,0.000000}%
\pgfsetstrokecolor{textcolor}%
\pgfsetfillcolor{textcolor}%
\pgftext[x=0.411111in, y=0.403549in, left, base]{\color{textcolor}{\sffamily\fontsize{7.000000}{8.400000}\selectfont\catcode`\^=\active\def^{\ifmmode\sp\else\^{}\fi}\catcode`\%=\active\def%{\%}\(n=1{,}859{,}904\)}}%
\end{pgfscope}%
\begin{pgfscope}%
\pgfsetbuttcap%
\pgfsetmiterjoin%
\definecolor{currentfill}{rgb}{0.666667,0.793651,0.412698}%
\pgfsetfillcolor{currentfill}%
\pgfsetlinewidth{1.003750pt}%
\definecolor{currentstroke}{rgb}{0.666667,0.793651,0.412698}%
\pgfsetstrokecolor{currentstroke}%
\pgfsetdash{}{0pt}%
\pgfpathmoveto{\pgfqpoint{0.138889in}{0.203704in}}%
\pgfpathlineto{\pgfqpoint{0.333333in}{0.203704in}}%
\pgfpathlineto{\pgfqpoint{0.333333in}{0.271759in}}%
\pgfpathlineto{\pgfqpoint{0.138889in}{0.271759in}}%
\pgfpathlineto{\pgfqpoint{0.138889in}{0.203704in}}%
\pgfpathclose%
\pgfusepath{stroke,fill}%
\end{pgfscope}%
\begin{pgfscope}%
\definecolor{textcolor}{rgb}{0.000000,0.000000,0.000000}%
\pgfsetstrokecolor{textcolor}%
\pgfsetfillcolor{textcolor}%
\pgftext[x=0.411111in, y=0.263117in, left, base]{\color{textcolor}{\sffamily\fontsize{7.000000}{8.400000}\selectfont\catcode`\^=\active\def^{\ifmmode\sp\else\^{}\fi}\catcode`\%=\active\def%{\%}(0, 15, 0)}}%
\end{pgfscope}%
\begin{pgfscope}%
\definecolor{textcolor}{rgb}{0.000000,0.000000,0.000000}%
\pgfsetstrokecolor{textcolor}%
\pgfsetfillcolor{textcolor}%
\pgftext[x=0.411111in, y=0.157793in, left, base]{\color{textcolor}{\sffamily\fontsize{7.000000}{8.400000}\selectfont\catcode`\^=\active\def^{\ifmmode\sp\else\^{}\fi}\catcode`\%=\active\def%{\%}\(n=262{,}209{,}728\)}}%
\end{pgfscope}%
\begin{pgfscope}%
\pgfsetbuttcap%
\pgfsetmiterjoin%
\definecolor{currentfill}{rgb}{0.269841,0.666667,0.222222}%
\pgfsetfillcolor{currentfill}%
\pgfsetlinewidth{1.003750pt}%
\definecolor{currentstroke}{rgb}{0.269841,0.666667,0.222222}%
\pgfsetstrokecolor{currentstroke}%
\pgfsetdash{}{0pt}%
\pgfpathmoveto{\pgfqpoint{1.386943in}{1.186727in}}%
\pgfpathlineto{\pgfqpoint{1.581387in}{1.186727in}}%
\pgfpathlineto{\pgfqpoint{1.581387in}{1.254782in}}%
\pgfpathlineto{\pgfqpoint{1.386943in}{1.254782in}}%
\pgfpathlineto{\pgfqpoint{1.386943in}{1.186727in}}%
\pgfpathclose%
\pgfusepath{stroke,fill}%
\end{pgfscope}%
\begin{pgfscope}%
\definecolor{textcolor}{rgb}{0.000000,0.000000,0.000000}%
\pgfsetstrokecolor{textcolor}%
\pgfsetfillcolor{textcolor}%
\pgftext[x=1.659165in, y=1.246140in, left, base]{\color{textcolor}{\sffamily\fontsize{7.000000}{8.400000}\selectfont\catcode`\^=\active\def^{\ifmmode\sp\else\^{}\fi}\catcode`\%=\active\def%{\%}(0, 12, 2)}}%
\end{pgfscope}%
\begin{pgfscope}%
\definecolor{textcolor}{rgb}{0.000000,0.000000,0.000000}%
\pgfsetstrokecolor{textcolor}%
\pgfsetfillcolor{textcolor}%
\pgftext[x=1.659165in, y=1.140816in, left, base]{\color{textcolor}{\sffamily\fontsize{7.000000}{8.400000}\selectfont\catcode`\^=\active\def^{\ifmmode\sp\else\^{}\fi}\catcode`\%=\active\def%{\%}\(n=55{,}296\)}}%
\end{pgfscope}%
\begin{pgfscope}%
\pgfsetbuttcap%
\pgfsetmiterjoin%
\definecolor{currentfill}{rgb}{0.079365,0.666667,0.571429}%
\pgfsetfillcolor{currentfill}%
\pgfsetlinewidth{1.003750pt}%
\definecolor{currentstroke}{rgb}{0.079365,0.666667,0.571429}%
\pgfsetstrokecolor{currentstroke}%
\pgfsetdash{}{0pt}%
\pgfpathmoveto{\pgfqpoint{1.386943in}{0.940971in}}%
\pgfpathlineto{\pgfqpoint{1.581387in}{0.940971in}}%
\pgfpathlineto{\pgfqpoint{1.581387in}{1.009027in}}%
\pgfpathlineto{\pgfqpoint{1.386943in}{1.009027in}}%
\pgfpathlineto{\pgfqpoint{1.386943in}{0.940971in}}%
\pgfpathclose%
\pgfusepath{stroke,fill}%
\end{pgfscope}%
\begin{pgfscope}%
\definecolor{textcolor}{rgb}{0.000000,0.000000,0.000000}%
\pgfsetstrokecolor{textcolor}%
\pgfsetfillcolor{textcolor}%
\pgftext[x=1.659165in, y=1.000384in, left, base]{\color{textcolor}{\sffamily\fontsize{7.000000}{8.400000}\selectfont\catcode`\^=\active\def^{\ifmmode\sp\else\^{}\fi}\catcode`\%=\active\def%{\%}(0, 13, 1)}}%
\end{pgfscope}%
\begin{pgfscope}%
\definecolor{textcolor}{rgb}{0.000000,0.000000,0.000000}%
\pgfsetstrokecolor{textcolor}%
\pgfsetfillcolor{textcolor}%
\pgftext[x=1.659165in, y=0.895061in, left, base]{\color{textcolor}{\sffamily\fontsize{7.000000}{8.400000}\selectfont\catcode`\^=\active\def^{\ifmmode\sp\else\^{}\fi}\catcode`\%=\active\def%{\%}\(n=76{,}544\)}}%
\end{pgfscope}%
\begin{pgfscope}%
\pgfsetbuttcap%
\pgfsetmiterjoin%
\definecolor{currentfill}{rgb}{0.000000,0.761905,0.857143}%
\pgfsetfillcolor{currentfill}%
\pgfsetlinewidth{1.003750pt}%
\definecolor{currentstroke}{rgb}{0.000000,0.761905,0.857143}%
\pgfsetstrokecolor{currentstroke}%
\pgfsetdash{}{0pt}%
\pgfpathmoveto{\pgfqpoint{1.386943in}{0.695215in}}%
\pgfpathlineto{\pgfqpoint{1.581387in}{0.695215in}}%
\pgfpathlineto{\pgfqpoint{1.581387in}{0.763271in}}%
\pgfpathlineto{\pgfqpoint{1.386943in}{0.763271in}}%
\pgfpathlineto{\pgfqpoint{1.386943in}{0.695215in}}%
\pgfpathclose%
\pgfusepath{stroke,fill}%
\end{pgfscope}%
\begin{pgfscope}%
\definecolor{textcolor}{rgb}{0.000000,0.000000,0.000000}%
\pgfsetstrokecolor{textcolor}%
\pgfsetfillcolor{textcolor}%
\pgftext[x=1.659165in, y=0.754629in, left, base]{\color{textcolor}{\sffamily\fontsize{7.000000}{8.400000}\selectfont\catcode`\^=\active\def^{\ifmmode\sp\else\^{}\fi}\catcode`\%=\active\def%{\%}(1, 5, 20)}}%
\end{pgfscope}%
\begin{pgfscope}%
\definecolor{textcolor}{rgb}{0.000000,0.000000,0.000000}%
\pgfsetstrokecolor{textcolor}%
\pgfsetfillcolor{textcolor}%
\pgftext[x=1.659165in, y=0.649305in, left, base]{\color{textcolor}{\sffamily\fontsize{7.000000}{8.400000}\selectfont\catcode`\^=\active\def^{\ifmmode\sp\else\^{}\fi}\catcode`\%=\active\def%{\%}\(n=880{,}768\)}}%
\end{pgfscope}%
\begin{pgfscope}%
\pgfsetbuttcap%
\pgfsetroundjoin%
\definecolor{currentfill}{rgb}{0.000000,0.317460,0.000000}%
\pgfsetfillcolor{currentfill}%
\pgfsetlinewidth{2.007500pt}%
\definecolor{currentstroke}{rgb}{0.000000,0.317460,0.000000}%
\pgfsetstrokecolor{currentstroke}%
\pgfsetdash{}{0pt}%
\pgfsys@defobject{currentmarker}{\pgfqpoint{-0.041667in}{-0.041667in}}{\pgfqpoint{0.041667in}{0.041667in}}{%
\pgfpathmoveto{\pgfqpoint{-0.041667in}{-0.041667in}}%
\pgfpathlineto{\pgfqpoint{0.041667in}{0.041667in}}%
\pgfpathmoveto{\pgfqpoint{-0.041667in}{0.041667in}}%
\pgfpathlineto{\pgfqpoint{0.041667in}{-0.041667in}}%
\pgfusepath{stroke,fill}%
}%
\begin{pgfscope}%
\pgfsys@transformshift{1.484165in}{0.474980in}%
\pgfsys@useobject{currentmarker}{}%
\end{pgfscope}%
\end{pgfscope}%
\begin{pgfscope}%
\definecolor{textcolor}{rgb}{0.000000,0.000000,0.000000}%
\pgfsetstrokecolor{textcolor}%
\pgfsetfillcolor{textcolor}%
\pgftext[x=1.659165in, y=0.508873in, left, base]{\color{textcolor}{\sffamily\fontsize{7.000000}{8.400000}\selectfont\catcode`\^=\active\def^{\ifmmode\sp\else\^{}\fi}\catcode`\%=\active\def%{\%}(1, 6, 16)}}%
\end{pgfscope}%
\begin{pgfscope}%
\definecolor{textcolor}{rgb}{0.000000,0.000000,0.000000}%
\pgfsetstrokecolor{textcolor}%
\pgfsetfillcolor{textcolor}%
\pgftext[x=1.659165in, y=0.403549in, left, base]{\color{textcolor}{\sffamily\fontsize{7.000000}{8.400000}\selectfont\catcode`\^=\active\def^{\ifmmode\sp\else\^{}\fi}\catcode`\%=\active\def%{\%}\(n=5{,}467{,}264\)}}%
\end{pgfscope}%
\begin{pgfscope}%
\pgfsetrectcap%
\pgfsetroundjoin%
\pgfsetlinewidth{1.505625pt}%
\definecolor{currentstroke}{rgb}{0.222222,0.190476,0.698413}%
\pgfsetstrokecolor{currentstroke}%
\pgfsetdash{}{0pt}%
\pgfpathmoveto{\pgfqpoint{2.524271in}{1.220755in}}%
\pgfpathlineto{\pgfqpoint{2.621493in}{1.220755in}}%
\pgfpathlineto{\pgfqpoint{2.718715in}{1.220755in}}%
\pgfusepath{stroke}%
\end{pgfscope}%
\begin{pgfscope}%
\definecolor{textcolor}{rgb}{0.000000,0.000000,0.000000}%
\pgfsetstrokecolor{textcolor}%
\pgfsetfillcolor{textcolor}%
\pgftext[x=2.796493in, y=1.246140in, left, base]{\color{textcolor}{\sffamily\fontsize{7.000000}{8.400000}\selectfont\catcode`\^=\active\def^{\ifmmode\sp\else\^{}\fi}\catcode`\%=\active\def%{\%}(1, 7, 12)}}%
\end{pgfscope}%
\begin{pgfscope}%
\definecolor{textcolor}{rgb}{0.000000,0.000000,0.000000}%
\pgfsetstrokecolor{textcolor}%
\pgfsetfillcolor{textcolor}%
\pgftext[x=2.796493in, y=1.140816in, left, base]{\color{textcolor}{\sffamily\fontsize{7.000000}{8.400000}\selectfont\catcode`\^=\active\def^{\ifmmode\sp\else\^{}\fi}\catcode`\%=\active\def%{\%}\(n=1{,}056{,}640\)}}%
\end{pgfscope}%
\begin{pgfscope}%
\pgfsetbuttcap%
\pgfsetmiterjoin%
\definecolor{currentfill}{rgb}{0.000001,0.873016,0.650794}%
\pgfsetfillcolor{currentfill}%
\pgfsetlinewidth{1.003750pt}%
\definecolor{currentstroke}{rgb}{0.000001,0.873016,0.650794}%
\pgfsetstrokecolor{currentstroke}%
\pgfsetdash{}{0pt}%
\pgfpathmoveto{\pgfqpoint{2.524271in}{0.940971in}}%
\pgfpathlineto{\pgfqpoint{2.718715in}{0.940971in}}%
\pgfpathlineto{\pgfqpoint{2.718715in}{1.009027in}}%
\pgfpathlineto{\pgfqpoint{2.524271in}{1.009027in}}%
\pgfpathlineto{\pgfqpoint{2.524271in}{0.940971in}}%
\pgfpathclose%
\pgfusepath{stroke,fill}%
\end{pgfscope}%
\begin{pgfscope}%
\definecolor{textcolor}{rgb}{0.000000,0.000000,0.000000}%
\pgfsetstrokecolor{textcolor}%
\pgfsetfillcolor{textcolor}%
\pgftext[x=2.796493in, y=1.000384in, left, base]{\color{textcolor}{\sffamily\fontsize{7.000000}{8.400000}\selectfont\catcode`\^=\active\def^{\ifmmode\sp\else\^{}\fi}\catcode`\%=\active\def%{\%}(1, 14, 0)}}%
\end{pgfscope}%
\begin{pgfscope}%
\definecolor{textcolor}{rgb}{0.000000,0.000000,0.000000}%
\pgfsetstrokecolor{textcolor}%
\pgfsetfillcolor{textcolor}%
\pgftext[x=2.796493in, y=0.895061in, left, base]{\color{textcolor}{\sffamily\fontsize{7.000000}{8.400000}\selectfont\catcode`\^=\active\def^{\ifmmode\sp\else\^{}\fi}\catcode`\%=\active\def%{\%}\(n=794{,}368\)}}%
\end{pgfscope}%
\begin{pgfscope}%
\pgfsetbuttcap%
\pgfsetmiterjoin%
\definecolor{currentfill}{rgb}{0.269841,0.603175,0.968254}%
\pgfsetfillcolor{currentfill}%
\pgfsetlinewidth{1.003750pt}%
\definecolor{currentstroke}{rgb}{0.269841,0.603175,0.968254}%
\pgfsetstrokecolor{currentstroke}%
\pgfsetdash{}{0pt}%
\pgfpathmoveto{\pgfqpoint{2.524271in}{0.695215in}}%
\pgfpathlineto{\pgfqpoint{2.718715in}{0.695215in}}%
\pgfpathlineto{\pgfqpoint{2.718715in}{0.763271in}}%
\pgfpathlineto{\pgfqpoint{2.524271in}{0.763271in}}%
\pgfpathlineto{\pgfqpoint{2.524271in}{0.695215in}}%
\pgfpathclose%
\pgfusepath{stroke,fill}%
\end{pgfscope}%
\begin{pgfscope}%
\definecolor{textcolor}{rgb}{0.000000,0.000000,0.000000}%
\pgfsetstrokecolor{textcolor}%
\pgfsetfillcolor{textcolor}%
\pgftext[x=2.796493in, y=0.754629in, left, base]{\color{textcolor}{\sffamily\fontsize{7.000000}{8.400000}\selectfont\catcode`\^=\active\def^{\ifmmode\sp\else\^{}\fi}\catcode`\%=\active\def%{\%}(2, 4, 23)}}%
\end{pgfscope}%
\begin{pgfscope}%
\definecolor{textcolor}{rgb}{0.000000,0.000000,0.000000}%
\pgfsetstrokecolor{textcolor}%
\pgfsetfillcolor{textcolor}%
\pgftext[x=2.796493in, y=0.649305in, left, base]{\color{textcolor}{\sffamily\fontsize{7.000000}{8.400000}\selectfont\catcode`\^=\active\def^{\ifmmode\sp\else\^{}\fi}\catcode`\%=\active\def%{\%}\(n=491{,}520\)}}%
\end{pgfscope}%
\begin{pgfscope}%
\pgfsetbuttcap%
\pgfsetmiterjoin%
\definecolor{currentfill}{rgb}{0.619048,0.476190,1.000000}%
\pgfsetfillcolor{currentfill}%
\pgfsetlinewidth{1.003750pt}%
\definecolor{currentstroke}{rgb}{0.619048,0.476190,1.000000}%
\pgfsetstrokecolor{currentstroke}%
\pgfsetdash{}{0pt}%
\pgfpathmoveto{\pgfqpoint{2.524271in}{0.449459in}}%
\pgfpathlineto{\pgfqpoint{2.718715in}{0.449459in}}%
\pgfpathlineto{\pgfqpoint{2.718715in}{0.517515in}}%
\pgfpathlineto{\pgfqpoint{2.524271in}{0.517515in}}%
\pgfpathlineto{\pgfqpoint{2.524271in}{0.449459in}}%
\pgfpathclose%
\pgfusepath{stroke,fill}%
\end{pgfscope}%
\begin{pgfscope}%
\definecolor{textcolor}{rgb}{0.000000,0.000000,0.000000}%
\pgfsetstrokecolor{textcolor}%
\pgfsetfillcolor{textcolor}%
\pgftext[x=2.796493in, y=0.508873in, left, base]{\color{textcolor}{\sffamily\fontsize{7.000000}{8.400000}\selectfont\catcode`\^=\active\def^{\ifmmode\sp\else\^{}\fi}\catcode`\%=\active\def%{\%}(2, 6, 15)}}%
\end{pgfscope}%
\begin{pgfscope}%
\definecolor{textcolor}{rgb}{0.000000,0.000000,0.000000}%
\pgfsetstrokecolor{textcolor}%
\pgfsetfillcolor{textcolor}%
\pgftext[x=2.796493in, y=0.403549in, left, base]{\color{textcolor}{\sffamily\fontsize{7.000000}{8.400000}\selectfont\catcode`\^=\active\def^{\ifmmode\sp\else\^{}\fi}\catcode`\%=\active\def%{\%}\(n=378{,}880\)}}%
\end{pgfscope}%
\begin{pgfscope}%
\pgfsetrectcap%
\pgfsetroundjoin%
\pgfsetlinewidth{1.505625pt}%
\definecolor{currentstroke}{rgb}{0.507937,0.031746,0.095238}%
\pgfsetstrokecolor{currentstroke}%
\pgfsetdash{}{0pt}%
\pgfpathmoveto{\pgfqpoint{3.661598in}{1.220755in}}%
\pgfpathlineto{\pgfqpoint{3.758821in}{1.220755in}}%
\pgfpathlineto{\pgfqpoint{3.856043in}{1.220755in}}%
\pgfusepath{stroke}%
\end{pgfscope}%
\begin{pgfscope}%
\definecolor{textcolor}{rgb}{0.000000,0.000000,0.000000}%
\pgfsetstrokecolor{textcolor}%
\pgfsetfillcolor{textcolor}%
\pgftext[x=3.933821in, y=1.246140in, left, base]{\color{textcolor}{\sffamily\fontsize{7.000000}{8.400000}\selectfont\catcode`\^=\active\def^{\ifmmode\sp\else\^{}\fi}\catcode`\%=\active\def%{\%}(3, 3, 29)}}%
\end{pgfscope}%
\begin{pgfscope}%
\definecolor{textcolor}{rgb}{0.000000,0.000000,0.000000}%
\pgfsetstrokecolor{textcolor}%
\pgfsetfillcolor{textcolor}%
\pgftext[x=3.933821in, y=1.140816in, left, base]{\color{textcolor}{\sffamily\fontsize{7.000000}{8.400000}\selectfont\catcode`\^=\active\def^{\ifmmode\sp\else\^{}\fi}\catcode`\%=\active\def%{\%}\(n=1{,}851{,}392\)}}%
\end{pgfscope}%
\begin{pgfscope}%
\pgfsetbuttcap%
\pgfsetmiterjoin%
\definecolor{currentfill}{rgb}{0.793651,0.396825,0.857143}%
\pgfsetfillcolor{currentfill}%
\pgfsetlinewidth{1.003750pt}%
\definecolor{currentstroke}{rgb}{0.793651,0.396825,0.857143}%
\pgfsetstrokecolor{currentstroke}%
\pgfsetdash{}{0pt}%
\pgfpathmoveto{\pgfqpoint{3.661598in}{0.940971in}}%
\pgfpathlineto{\pgfqpoint{3.856043in}{0.940971in}}%
\pgfpathlineto{\pgfqpoint{3.856043in}{1.009027in}}%
\pgfpathlineto{\pgfqpoint{3.661598in}{1.009027in}}%
\pgfpathlineto{\pgfqpoint{3.661598in}{0.940971in}}%
\pgfpathclose%
\pgfusepath{stroke,fill}%
\end{pgfscope}%
\begin{pgfscope}%
\definecolor{textcolor}{rgb}{0.000000,0.000000,0.000000}%
\pgfsetstrokecolor{textcolor}%
\pgfsetfillcolor{textcolor}%
\pgftext[x=3.933821in, y=1.000384in, left, base]{\color{textcolor}{\sffamily\fontsize{7.000000}{8.400000}\selectfont\catcode`\^=\active\def^{\ifmmode\sp\else\^{}\fi}\catcode`\%=\active\def%{\%}(4, 2, 35)}}%
\end{pgfscope}%
\begin{pgfscope}%
\definecolor{textcolor}{rgb}{0.000000,0.000000,0.000000}%
\pgfsetstrokecolor{textcolor}%
\pgfsetfillcolor{textcolor}%
\pgftext[x=3.933821in, y=0.895061in, left, base]{\color{textcolor}{\sffamily\fontsize{7.000000}{8.400000}\selectfont\catcode`\^=\active\def^{\ifmmode\sp\else\^{}\fi}\catcode`\%=\active\def%{\%}\(n=485{,}376\)}}%
\end{pgfscope}%
\begin{pgfscope}%
\pgfsetbuttcap%
\pgfsetmiterjoin%
\definecolor{currentfill}{rgb}{0.888889,0.349206,0.603175}%
\pgfsetfillcolor{currentfill}%
\pgfsetlinewidth{1.003750pt}%
\definecolor{currentstroke}{rgb}{0.888889,0.349206,0.603175}%
\pgfsetstrokecolor{currentstroke}%
\pgfsetdash{}{0pt}%
\pgfpathmoveto{\pgfqpoint{3.661598in}{0.695215in}}%
\pgfpathlineto{\pgfqpoint{3.856043in}{0.695215in}}%
\pgfpathlineto{\pgfqpoint{3.856043in}{0.763271in}}%
\pgfpathlineto{\pgfqpoint{3.661598in}{0.763271in}}%
\pgfpathlineto{\pgfqpoint{3.661598in}{0.695215in}}%
\pgfpathclose%
\pgfusepath{stroke,fill}%
\end{pgfscope}%
\begin{pgfscope}%
\definecolor{textcolor}{rgb}{0.000000,0.000000,0.000000}%
\pgfsetstrokecolor{textcolor}%
\pgfsetfillcolor{textcolor}%
\pgftext[x=3.933821in, y=0.754629in, left, base]{\color{textcolor}{\sffamily\fontsize{7.000000}{8.400000}\selectfont\catcode`\^=\active\def^{\ifmmode\sp\else\^{}\fi}\catcode`\%=\active\def%{\%}(5, 1, 43)}}%
\end{pgfscope}%
\begin{pgfscope}%
\definecolor{textcolor}{rgb}{0.000000,0.000000,0.000000}%
\pgfsetstrokecolor{textcolor}%
\pgfsetfillcolor{textcolor}%
\pgftext[x=3.933821in, y=0.649305in, left, base]{\color{textcolor}{\sffamily\fontsize{7.000000}{8.400000}\selectfont\catcode`\^=\active\def^{\ifmmode\sp\else\^{}\fi}\catcode`\%=\active\def%{\%}\(n=204{,}800\)}}%
\end{pgfscope}%
\begin{pgfscope}%
\pgfsetbuttcap%
\pgfsetmiterjoin%
\definecolor{currentfill}{rgb}{1.000000,0.587302,0.888889}%
\pgfsetfillcolor{currentfill}%
\pgfsetlinewidth{1.003750pt}%
\definecolor{currentstroke}{rgb}{1.000000,0.587302,0.888889}%
\pgfsetstrokecolor{currentstroke}%
\pgfsetdash{}{0pt}%
\pgfpathmoveto{\pgfqpoint{3.661598in}{0.449459in}}%
\pgfpathlineto{\pgfqpoint{3.856043in}{0.449459in}}%
\pgfpathlineto{\pgfqpoint{3.856043in}{0.517515in}}%
\pgfpathlineto{\pgfqpoint{3.661598in}{0.517515in}}%
\pgfpathlineto{\pgfqpoint{3.661598in}{0.449459in}}%
\pgfpathclose%
\pgfusepath{stroke,fill}%
\end{pgfscope}%
\begin{pgfscope}%
\definecolor{textcolor}{rgb}{0.000000,0.000000,0.000000}%
\pgfsetstrokecolor{textcolor}%
\pgfsetfillcolor{textcolor}%
\pgftext[x=3.933821in, y=0.508873in, left, base]{\color{textcolor}{\sffamily\fontsize{7.000000}{8.400000}\selectfont\catcode`\^=\active\def^{\ifmmode\sp\else\^{}\fi}\catcode`\%=\active\def%{\%}(7, 0, 50)}}%
\end{pgfscope}%
\begin{pgfscope}%
\definecolor{textcolor}{rgb}{0.000000,0.000000,0.000000}%
\pgfsetstrokecolor{textcolor}%
\pgfsetfillcolor{textcolor}%
\pgftext[x=3.933821in, y=0.403549in, left, base]{\color{textcolor}{\sffamily\fontsize{7.000000}{8.400000}\selectfont\catcode`\^=\active\def^{\ifmmode\sp\else\^{}\fi}\catcode`\%=\active\def%{\%}\(n=47{,}104\)}}%
\end{pgfscope}%
\end{pgfpicture}%
\makeatother%
\endgroup%

    \caption{The event cost landscape for the solutions returned by Synesth for Class~1 Cas gene synteny dataset (see text). The cost of a loss event is fixed to 1. For each color region, the legend shows a condensed event count vector of the form \((c_\eventDup + c_\eventCut, c_\eventTrDup + c_\eventTrCut, c_\eventLoss)\). “\(n\)” indicates the number of distinct minimal histories for any set of event costs in that region.}
    \label{fig:synesth:crispr-pareto}
\end{figure}

\section{Conclusion}

Synesth is a flexible tool for tree reconciliation, allowing for a wide range of segmental events, addressing the inevitable incompleteness of the input dataset in terms of unsampled species, and offering various optimization and output criteria to the user.
Moreover, its time complexity brings it up to the level of the most time-efficient reconciliation algorithms such as \emph{ecceTERA}~\parencite{Jacox2016EcceteraComprehensiveGene} and \emph{RANGER-DTLx}~\parencite{Weiner2021ImprovedDuplicationTransfer}, but for an evolutionary model with events involving sets of genes rather than single genes.

The inductive characterization of Pareto optimal histories allows for an exhaustive exploration of the solution space.
Taking advantage of this flexibility, future extensions of this work may address the problem of formally characterizing this space in terms of constructing equivalence classes or normalized histories, uniformly sampling the space of histories, and assigning confidence values to predicted histories.

Moreover, further extensions of the model are required, such as accounting for gene orders in syntenies or considering segmental gains as an additional parameter for selecting Pareto optimal histories.
However, probably the most questionable limitation of the model is the absence of synteny fusions, while synteny fissions are allowed, favouring large syntenies up to the root of the tree.
Note however that including fusions, as well as including tandem duplications (duplications inside a given synteny) will require adding reticulated nodes.
It will be interesting to see whether our dynamic programming scheme can be generalized to such phylogenetic networks, or if it makes the problem NP-hard, in which case tree decomposition methods may be explored~\parencite{Scornavacca2022TreewidthBasedAlgorithms}.

Another important challenge not addressed in this paper is how to obtain an input synteny tree.
In fact, phylogenetic methods rather output sets of gene trees, one for each gene family.
If the individual gene trees are \enquote{consistent}, i.e., with no contradictory phylogenetic information, then a tree displaying them all can be obtained.
However, even in this case, there may be an exponential number of such supertrees.
In~\textcite{Delabre2020EvolutionSegmentalDuplications}, the suggested solution was to test each possible supertree, and retain the one leading to the most parsimonious reconciliation.
An alternative would be to simultaneously construct and reconcile a supertree with a given species tree. This opens the door to interesting future investigations.
